\documentclass[a4paper]{book}
\usepackage[times,inconsolata,hyper]{Rd}
\usepackage{makeidx}
\usepackage[utf8]{inputenc} % @SET ENCODING@
% \usepackage{graphicx} % @USE GRAPHICX@
\makeindex{}
\begin{document}
\chapter*{}
\begin{center}
{\textbf{\huge Package `rSHUD'}}
\par\bigskip{\large \today}
\end{center}
\inputencoding{utf8}
\ifthenelse{\boolean{Rd@use@hyper}}{\hypersetup{pdftitle = {rSHUD: Toolbox For SHUD Model System}}}{}\ifthenelse{\boolean{Rd@use@hyper}}{\hypersetup{pdfauthor = {Lele Shu}}}{}\begin{description}
\raggedright{}
\item[Type]\AsIs{Package}
\item[Title]\AsIs{Toolbox For SHUD Model System}
\item[Version]\AsIs{2.2.0}
\item[Date]\AsIs{2024-11-23}
\item[Maintainer]\AsIs{Lele Shu }\email{shulele@lzb.ac.cn}\AsIs{}
\item[Description]\AsIs{The Simulator for Hydrologic Unstructured Domains (SHUD) is a multiprocess,
multi-scale hydrologic model where the major hydrological processes are fully coupled
using the finite volume method. This toolbox provides comprehensive tools to convert
geospatial data into SHUD format, read/write input files, read SHUD output files,
perform hydrological analysis, and spatially visualize SHUD results. It supports
watershed delineation, mesh generation, river network processing, and various
hydrological calculations including PET, melt factor, and water balance analysis.}
\item[URL]\AsIs{}\url{https://www.shud.xyz/}\AsIs{}
\item[BugReports]\AsIs{}\url{https://shud-system.github.io/}\AsIs{}
\item[License]\AsIs{MIT + file LICENSE}
\item[LazyData]\AsIs{TRUE}
\item[Depends]\AsIs{R (>= 4.0.0)}
\item[Encoding]\AsIs{UTF-8}
\item[Imports]\AsIs{Rcpp,
reshape2,
ggplot2,
gridExtra,
grid,
fields,
xts,
hydroGOF,
zoo,
terra (>= 1.7-0),
sf (>= 1.0-0),
RTriangle,
proj4,
gstat,
abind,
utils,
lubridate,
doParallel,
geometry,
methods}
\item[LinkingTo]\AsIs{Rcpp}
\item[RoxygenNote]\AsIs{7.2.3}
\item[Remotes]\AsIs{shulele/RTriangle/pkg}
\item[Suggests]\AsIs{testthat,
knitr,
rmarkdown,
deldir,
interp,
whitebox,
ncdf4,
raster,
sp,
colorspace}
\item[VignetteBuilder]\AsIs{knitr}
\item[Roxygen]\AsIs{list(markdown = TRUE)}
\item[Language]\AsIs{en-US}
\end{description}
\Rdcontents{\R{} topics documented:}
\inputencoding{utf8}
\HeaderA{.shud}{This is .shud environment.}{.shud}
\keyword{datasets}{.shud}
%
\begin{Description}\relax
This is .shud environment.
\end{Description}
%
\begin{Usage}
\begin{verbatim}
.shud
\end{verbatim}
\end{Usage}
%
\begin{Format}
An object of class \code{environment} of length 0.
\end{Format}
\inputencoding{utf8}
\HeaderA{AddHoleToPolygon}{Add holes into Polygons \code{AddHoleToPolygon}}{AddHoleToPolygon}
%
\begin{Description}\relax
Add holes into Polygons
\code{AddHoleToPolygon}
\end{Description}
%
\begin{Usage}
\begin{verbatim}
AddHoleToPolygon(poly, hole)
\end{verbatim}
\end{Usage}
%
\begin{Arguments}
\begin{ldescription}
\item[\code{poly}] SpatialPolygons

\item[\code{hole}] Hole Polygon
\end{ldescription}
\end{Arguments}
\inputencoding{utf8}
\HeaderA{AerodynamicResistance}{Aerodynamic Resistance. Eq 4.2.25 in David R Maidment, Handbook of Hydrology \code{AerodynamicResistance}}{AerodynamicResistance}
%
\begin{Description}\relax
Aerodynamic Resistance. Eq 4.2.25 in David R Maidment, Handbook of Hydrology
\code{AerodynamicResistance}
\end{Description}
%
\begin{Usage}
\begin{verbatim}
AerodynamicResistance(Uz, hc, Z_u, Z_e, VON_KARMAN = 0.4)
\end{verbatim}
\end{Usage}
%
\begin{Arguments}
\begin{ldescription}
\item[\code{Uz}] Wind speed at height z \LinkA{m s-1}{m s.Rdash.1}

\item[\code{hc}] Height of crop \LinkA{m}{m}

\item[\code{Z\_u}] height of wind measurements \LinkA{m}{m}

\item[\code{Z\_e}] height of humidity measurements \LinkA{m}{m}

\item[\code{VON\_KARMAN}] von Karman's constant. Default = 0.41 \LinkA{-}{.Rdash.}
\end{ldescription}
\end{Arguments}
%
\begin{Value}
Aerodynamic resistance \LinkA{s m-1}{s m.Rdash.1}
\end{Value}
\inputencoding{utf8}
\HeaderA{AirDensity}{Air Density, Eq 4.2.4 in David R Maidment, Handbook of Hydrology \code{shud.ic}}{AirDensity}
%
\begin{Description}\relax
Air Density, Eq 4.2.4 in David R Maidment, Handbook of Hydrology
\code{shud.ic}
\end{Description}
%
\begin{Usage}
\begin{verbatim}
AirDensity(P, Temp)
\end{verbatim}
\end{Usage}
%
\begin{Arguments}
\begin{ldescription}
\item[\code{P}] Air pressure \LinkA{kP}{kP}

\item[\code{Temp}] Air temperature \LinkA{C}{C}
\end{ldescription}
\end{Arguments}
%
\begin{Value}
Density of Air. \LinkA{kg m-3}{kg m.Rdash.3}
\end{Value}
\inputencoding{utf8}
\HeaderA{as\_shud\_river}{Convert River Network to SHUD.RIVER S4 Object}{as.Rul.shud.Rul.river}
%
\begin{Description}\relax
Converts a river network (from build\_river\_network) to the SHUD.RIVER
S4 class format for compatibility with existing SHUD workflows.
\end{Description}
%
\begin{Usage}
\begin{verbatim}
as_shud_river(network_list)
\end{verbatim}
\end{Usage}
%
\begin{Arguments}
\begin{ldescription}
\item[\code{network\_list}] List output from \code{build\_river\_network()}
\end{ldescription}
\end{Arguments}
%
\begin{Details}\relax
This function creates a SHUD.RIVER object that includes both the legacy
data.frame format (for backward compatibility) and the modern sf format
(stored in the network slot).
\end{Details}
%
\begin{Value}
SHUD.RIVER S4 object
\end{Value}
%
\begin{Examples}
\begin{ExampleCode}
## Not run: 
library(sf)
library(terra)

rivers <- st_read("rivers.shp")
dem <- rast("dem.tif")
network_list <- build_river_network(rivers, dem)
shud_river <- as_shud_river(network_list)

## End(Not run)
\end{ExampleCode}
\end{Examples}
\inputencoding{utf8}
\HeaderA{auto\_build\_model}{Automatically Build SHUD Model}{auto.Rul.build.Rul.model}
%
\begin{Description}\relax
Automated workflow to build a complete SHUD hydrological model from
geospatial data. This modernized function uses terra and sf exclusively
for all spatial operations.
\end{Description}
%
\begin{Usage}
\begin{verbatim}
auto_build_model(
  project_name,
  domain,
  rivers = NULL,
  dem,
  soil = NULL,
  geology = NULL,
  landcover = NULL,
  forcing_sites = NULL,
  forcing_files = NULL,
  output_dir = getwd(),
  mesh_options = list(),
  river_options = list(),
  aquifer_depth = 10,
  years = 2000:2010,
  clean = TRUE,
  parameters = NULL,
  calibration = NULL,
  melt_factor = NULL,
  remove_outliers = TRUE,
  verbose = TRUE
)
\end{verbatim}
\end{Usage}
%
\begin{Arguments}
\begin{ldescription}
\item[\code{project\_name}] Character, project name for output files

\item[\code{domain}] sf object with POLYGON geometry defining watershed boundary

\item[\code{rivers}] sf object with LINESTRING geometry for river network (optional)

\item[\code{dem}] SpatRaster object with elevation data

\item[\code{soil}] SpatRaster object with soil data (optional)

\item[\code{geology}] SpatRaster object with geology data (optional)

\item[\code{landcover}] SpatRaster object with land cover data (optional)

\item[\code{forcing\_sites}] sf object with POINT geometry for forcing locations (optional)

\item[\code{forcing\_files}] Character vector of forcing data filenames

\item[\code{output\_dir}] Character, output directory path (default: current directory)

\item[\code{mesh\_options}] List of mesh generation options

\item[\code{river\_options}] List of river processing options

\item[\code{aquifer\_depth}] Numeric, depth of aquifer bottom below surface (default: 10 m)

\item[\code{years}] Integer vector, years for LAI/melt factor time series

\item[\code{clean}] Logical, whether to clean existing files in output directory

\item[\code{parameters}] List or data.frame, model parameters (default: shud.para())

\item[\code{calibration}] List or data.frame, calibration parameters (default: shud.calib())

\item[\code{melt\_factor}] xts object, melt factor time series (optional)

\item[\code{remove\_outliers}] Logical, remove outliers in soil/geology data

\item[\code{verbose}] Logical, display progress messages (default: TRUE)
\end{ldescription}
\end{Arguments}
%
\begin{Value}
List with components:
\begin{ldescription}
\item[\code{mesh}] SHUD.MESH object
\item[\code{river}] SHUD.RIVER object
\item[\code{attributes}] Attribute data frame
\item[\code{files}] Named vector of output file paths
\end{ldescription}
\end{Value}
%
\begin{Section}{Migration Note}

This function replaces \code{autoBuildModel()} with modern spatial libraries.
Key differences:
\begin{itemize}

\item{} Only accepts terra::SpatRaster and sf objects
\item{} Rejects legacy raster/sp objects with clear error messages
\item{} Uses snake\_case parameter names (e.g., project\_name vs prjname)
\item{} Returns structured list instead of just mesh object
\item{} Provides verbose progress reporting

\end{itemize}


To migrate from legacy code:
\begin{alltt}
# Old code (raster/sp):
# wbd <- readOGR("boundary.shp")
# dem <- raster("dem.tif")

# New code (sf/terra):
wbd <- sf::st_read("boundary.shp")
dem <- terra::rast("dem.tif")
\end{alltt}

\end{Section}
%
\begin{Examples}
\begin{ExampleCode}
## Not run: 
library(sf)
library(terra)

# Load data using modern libraries
domain <- st_read("watershed.shp")
rivers <- st_read("rivers.shp")
dem <- rast("dem.tif")
soil <- rast("soil.tif")
lc <- rast("landcover.tif")

# Build model
model <- auto_build_model(
  project_name = "my_watershed",
  domain = domain,
  rivers = rivers,
  dem = dem,
  soil = soil,
  landcover = lc,
  output_dir = "model_output",
  years = 2000:2010
)

# Access results
mesh <- model$mesh
river_network <- model$river

## End(Not run)
\end{ExampleCode}
\end{Examples}
\inputencoding{utf8}
\HeaderA{autoBuildModel}{Automatically run SHUDgis based on the input data and parameters. \code{autoBuildModel}}{autoBuildModel}
%
\begin{Description}\relax
Automatically run SHUDgis based on the input data and parameters.
\code{autoBuildModel}
\end{Description}
%
\begin{Usage}
\begin{verbatim}
autoBuildModel(
  indata,
  forcfiles,
  prjname = "sac",
  outdir = getwd(),
  a.max = 1e+06 * 0.2,
  q.min = 33,
  tol.riv = 0,
  tol.wb = 200,
  tol.len = 0,
  AqDepth = 10,
  years = 2000:2010,
  clean = TRUE,
  cfg.para = shud.para(nday = 365 * length(years) + round(length(years)/4)),
  cfg.calib = shud.calib(),
  mf = MeltFactor(years = years),
  rm.outlier = TRUE,
  quiet = FALSE
)
\end{verbatim}
\end{Usage}
%
\begin{Arguments}
\begin{ldescription}
\item[\code{indata}] Input data, list of data.

\item[\code{forcfiles}] Filenames of .csv forcing data

\item[\code{prjname}] Projectname

\item[\code{outdir}] Output directory

\item[\code{a.max}] Maximum area of triangles. Unit=m2

\item[\code{q.min}] Minimum of angles for each triangle

\item[\code{tol.riv}] Tolerance to simplify river lines, in meter

\item[\code{tol.wb}] Tolerance to simplify watershed boundary, in meter

\item[\code{tol.len}] Tolerance to simplyfy river by longituanal length, meter.

\item[\code{AqDepth}] Depth of aquifer bottom, which is surface elemvation minus bedrock elevation. meter.

\item[\code{years}] Years to generate the LAI, Meltfactor, Roughness Length, which must be same as the period of forcing data.

\item[\code{clean}] Whether clean the existing model files in output directory.

\item[\code{cfg.para}] model configuration, parameter

\item[\code{cfg.calib}] model calibration

\item[\code{mf}] Meltfactor

\item[\code{rm.outlier}] Whether to remove the outlier in soil/geol;

\item[\code{quiet}] Whether to ask confirmation when multiple outlets exist.
\end{ldescription}
\end{Arguments}
%
\begin{Value}
\code{SHUD.mesh}
\end{Value}
%
\begin{Examples}
\begin{ExampleCode}
data(sac)
indata = sac
sp.forc = indata[['forc']]
forc.fns = paste0(sp.forc@data[, 'NLDAS_ID'], '.csv')
forc.fns
\end{ExampleCode}
\end{Examples}
\inputencoding{utf8}
\HeaderA{BasicPlot}{Read multiply SHUD model output files and do time-series plot}{BasicPlot}
\keyword{output.}{BasicPlot}
\keyword{read}{BasicPlot}
%
\begin{Description}\relax
Read multiply SHUD model output files and do time-series plot
\end{Description}
%
\begin{Usage}
\begin{verbatim}
BasicPlot(
  varname = c(paste0("eley", c("surf", "unsat", "gw", "snow")), paste0("elev", c("prcp",
    "infil", "rech")), paste0("elev", c("etp", "eta", "etev", "ettr", "etic")),
    paste0("rivq", c("down", "sub", "surf")), paste0("rivy", "stage")),
  xl = NULL,
  sp.riv = NULL,
  rdsfile = NULL,
  plot = TRUE,
  imap = FALSE,
  w.focal = matrix(1/9, 3, 3),
  path = file.path(get("anapath", envir = .shud), "BasicPlot")
)
\end{verbatim}
\end{Usage}
%
\begin{Arguments}
\begin{ldescription}
\item[\code{varname}] vector of output keywords

\item[\code{xl}] list of data returned from loaddata()

\item[\code{sp.riv}] River SpatialLine*

\item[\code{rdsfile}] Save RDS file

\item[\code{plot}] Whether do the time-series plot

\item[\code{imap}] Whether do the raster plot for Element data. Only works for element data

\item[\code{w.focal}] forcal matrix

\item[\code{return}] Whether return the data. Some the results are too huge to load in memoery at once.

\item[\code{iRDS}] Whether save RDS file.
\end{ldescription}
\end{Arguments}
%
\begin{Value}
A list of TimeSeries data.
\end{Value}
\inputencoding{utf8}
\HeaderA{BC}{Generate the default Boundary Condition \code{BC}}{BC}
%
\begin{Description}\relax
Generate the default Boundary Condition
\code{BC}
\end{Description}
%
\begin{Usage}
\begin{verbatim}
BC(x = cbind(0:100, 0), start = as.POSIXct("2000-01-01"))
\end{verbatim}
\end{Usage}
%
\begin{Arguments}
\begin{ldescription}
\item[\code{x}] matrix of BC. c(days, value)

\item[\code{start}] Start time, POSIXct.
\end{ldescription}
\end{Arguments}
%
\begin{Value}
Time-series data of BC
\end{Value}
\inputencoding{utf8}
\HeaderA{build\_river\_network}{Build River Network}{build.Rul.river.Rul.network}
%
\begin{Description}\relax
Constructs a complete river network with all properties calculated,
including stream order, downstream relationships, lengths, slopes,
and hydraulic parameters. This is the main function for river network
preparation in SHUD models.
\end{Description}
%
\begin{Usage}
\begin{verbatim}
build_river_network(
  rivers,
  dem,
  area = NULL,
  river_order = NULL,
  downstream = NULL,
  tolerance = NULL
)
\end{verbatim}
\end{Usage}
%
\begin{Arguments}
\begin{ldescription}
\item[\code{rivers}] sf object with LINESTRING geometry representing river segments

\item[\code{dem}] SpatRaster object with elevation data

\item[\code{area}] Numeric, watershed area for width estimation (optional)

\item[\code{river\_order}] Integer vector of stream orders (optional, will be calculated if NULL)

\item[\code{downstream}] Integer vector of downstream indices (optional, will be calculated if NULL)

\item[\code{tolerance}] Numeric, tolerance for simplification (default: NULL, auto-calculated)
\end{ldescription}
\end{Arguments}
%
\begin{Details}\relax
This function integrates all river processing steps:
\begin{enumerate}

\item{} Calculate stream order (Strahler classification)
\item{} Identify downstream relationships
\item{} Calculate segment lengths
\item{} Extract elevations and calculate slopes
\item{} Generate hydraulic parameters for each river type

\end{enumerate}


The output is ready for use in SHUD model input files.
\end{Details}
%
\begin{Value}
List with components:
\begin{ldescription}
\item[\code{network}] sf object with all river properties
\item[\code{river\_types}] data.frame with hydraulic parameters for each type
\item[\code{points}] data.frame with from/to node coordinates and elevations
\end{ldescription}
\end{Value}
%
\begin{Section}{Migration Note}

This function replaces \code{shud.river()} and uses sf/terra instead of sp/raster.
The output structure is similar but uses modern spatial formats.
\end{Section}
%
\begin{Examples}
\begin{ExampleCode}
## Not run: 
library(sf)
library(terra)

# Load data
rivers <- st_read("rivers.shp")
dem <- rast("dem.tif")

# Build complete river network
river_network <- build_river_network(rivers, dem, area = 1000)

# Access components
network_sf <- river_network$network
river_params <- river_network$river_types
node_info <- river_network$points

## End(Not run)
\end{ExampleCode}
\end{Examples}
\inputencoding{utf8}
\HeaderA{BulkSurfaceResistance}{Bulk Surface Resistance, Eq 4.2.22 in David R Maidment, Handbook of Hydrology \code{BulkSurfaceResistance}}{BulkSurfaceResistance}
%
\begin{Description}\relax
Bulk Surface Resistance, Eq 4.2.22 in David R Maidment, Handbook of Hydrology
\code{BulkSurfaceResistance}
\end{Description}
%
\begin{Usage}
\begin{verbatim}
BulkSurfaceResistance(lai)
\end{verbatim}
\end{Usage}
%
\begin{Arguments}
\begin{ldescription}
\item[\code{lai}] Leaf Area Index \LinkA{m2 m-2}{m2 m.Rdash.2}
\end{ldescription}
\end{Arguments}
%
\begin{Value}
Bulk Surface Resistance \LinkA{s m-1}{s m.Rdash.1}
\end{Value}
\inputencoding{utf8}
\HeaderA{calc\_river\_downstream}{Calculate Downstream Relationships}{calc.Rul.river.Rul.downstream}
%
\begin{Description}\relax
Determines the downstream segment index for each river segment in a network.
A segment's downstream is the segment whose start node matches this segment's
end node.
\end{Description}
%
\begin{Usage}
\begin{verbatim}
calc_river_downstream(rivers, coords = NULL, tolerance = NULL)
\end{verbatim}
\end{Usage}
%
\begin{Arguments}
\begin{ldescription}
\item[\code{rivers}] sf object with LINESTRING geometry representing river segments

\item[\code{coords}] Matrix of unique coordinates (optional, will be calculated if NULL)

\item[\code{tolerance}] Numeric, tolerance for simplification (default: NULL, auto-calculated)
\end{ldescription}
\end{Arguments}
%
\begin{Details}\relax
The function identifies downstream relationships by matching end nodes to
start nodes. Outlets (segments with no downstream) are marked with -1.
If multiple downstream segments are found (unusual), the first is selected
and a message is displayed.
\end{Details}
%
\begin{Value}
Integer vector of downstream indices (-1 for outlets, -3 for no downstream)
\end{Value}
%
\begin{Section}{Migration Note}

This function replaces \code{sp.RiverDown()} and uses sf instead of sp.
\end{Section}
%
\begin{Examples}
\begin{ExampleCode}
## Not run: 
library(sf)
rivers <- st_read("rivers.shp")
downstream <- calc_river_downstream(rivers)
# Identify outlets
outlets <- which(downstream < 0)

## End(Not run)
\end{ExampleCode}
\end{Examples}
\inputencoding{utf8}
\HeaderA{calc\_river\_order}{Calculate River Order (Strahler Classification)}{calc.Rul.river.Rul.order}
%
\begin{Description}\relax
Calculates the Strahler stream order for a river network. This function
uses sf for spatial operations and implements the Strahler algorithm
for stream classification.
\end{Description}
%
\begin{Usage}
\begin{verbatim}
calc_river_order(rivers, tolerance = NULL)
\end{verbatim}
\end{Usage}
%
\begin{Arguments}
\begin{ldescription}
\item[\code{rivers}] sf object with LINESTRING geometry representing river segments

\item[\code{tolerance}] Numeric, tolerance for simplification (default: NULL, auto-calculated)
\end{ldescription}
\end{Arguments}
%
\begin{Details}\relax
The Strahler stream order is calculated as follows:
\begin{itemize}

\item{} First-order streams are headwater streams with no tributaries
\item{} When two streams of order i join, they form a stream of order i+1
\item{} When streams of different orders join, the resulting stream has the
order of the higher-order stream

\end{itemize}


The function automatically simplifies the river network slightly to
handle minor coordinate differences at junctions.
\end{Details}
%
\begin{Value}
Integer vector of stream orders for each river segment
\end{Value}
%
\begin{Section}{Migration Note}

This function replaces \code{sp.RiverOrder()} and uses sf instead of sp.
The algorithm and results are equivalent, but the input must be an sf object.
\end{Section}
%
\begin{Examples}
\begin{ExampleCode}
## Not run: 
library(sf)
# Load river network as sf object
rivers <- st_read("rivers.shp")
# Calculate stream order
order <- calc_river_order(rivers)
# Plot by order
plot(rivers["geometry"], col = order)

## End(Not run)
\end{ExampleCode}
\end{Examples}
\inputencoding{utf8}
\HeaderA{calc\_river\_path}{Calculate River Paths}{calc.Rul.river.Rul.path}
%
\begin{Description}\relax
Dissolves river network into continuous paths from headwaters to outlets
or junctions. Each path represents a continuous flow line.
\end{Description}
%
\begin{Usage}
\begin{verbatim}
calc_river_path(rivers, coords = NULL, downstream = NULL, tolerance = 0)
\end{verbatim}
\end{Usage}
%
\begin{Arguments}
\begin{ldescription}
\item[\code{rivers}] sf object with LINESTRING geometry representing river segments

\item[\code{coords}] Matrix of unique coordinates (optional, will be calculated if NULL)

\item[\code{downstream}] Integer vector of downstream indices (optional, will be calculated if NULL)

\item[\code{tolerance}] Numeric, tolerance for simplification (default: 0, no simplification)
\end{ldescription}
\end{Arguments}
%
\begin{Details}\relax
The function traces upstream from outlets and junctions to create continuous
flow paths. This is useful for simplifying river networks and identifying
main stems vs. tributaries.
\end{Details}
%
\begin{Value}
List with components:
\begin{ldescription}
\item[\code{seg\_ids}] List of segment IDs for each path
\item[\code{point\_ids}] List of point IDs for each path
\item[\code{paths}] sf object with dissolved paths as LINESTRING
\end{ldescription}
\end{Value}
%
\begin{Section}{Migration Note}

This function replaces \code{sp.RiverPath()} and uses sf instead of sp.
The output structure is similar but uses sf objects instead of SpatialLines.
\end{Section}
%
\begin{Examples}
\begin{ExampleCode}
## Not run: 
library(sf)
rivers <- st_read("rivers.shp")
paths <- calc_river_path(rivers)
# Plot dissolved paths
plot(paths$paths)

## End(Not run)
\end{ExampleCode}
\end{Examples}
\inputencoding{utf8}
\HeaderA{calc\_river\_properties}{Calculate River Properties}{calc.Rul.river.Rul.properties}
%
\begin{Description}\relax
Integrated function to calculate multiple river properties including
stream order, downstream relationships, length, and slope. This function
provides a unified interface to avoid redundant calculations.
\end{Description}
%
\begin{Usage}
\begin{verbatim}
calc_river_properties(rivers, dem = NULL, properties = "all", tolerance = NULL)
\end{verbatim}
\end{Usage}
%
\begin{Arguments}
\begin{ldescription}
\item[\code{rivers}] sf object with LINESTRING geometry representing river segments

\item[\code{dem}] SpatRaster object with elevation data (required for slope calculation)

\item[\code{properties}] Character vector specifying which properties to calculate.
Options: "order", "downstream", "length", "slope", "all" (default: "all")

\item[\code{tolerance}] Numeric, tolerance for simplification (default: NULL, auto-calculated)
\end{ldescription}
\end{Arguments}
%
\begin{Details}\relax
This function integrates multiple river property calculations to avoid
redundant processing. Properties are calculated efficiently by reusing
intermediate results (e.g., coordinates, from-to nodes).

Available properties:
\begin{itemize}

\item{} \strong{order}: Strahler stream order
\item{} \strong{downstream}: Index of downstream segment (-1 for outlets)
\item{} \strong{length}: Length of each segment in map units
\item{} \strong{slope}: Bed slope calculated from DEM

\end{itemize}

\end{Details}
%
\begin{Value}
sf object with calculated properties added as columns
\end{Value}
%
\begin{Section}{Migration Note}

This function integrates functionality from multiple legacy functions:
\code{sp.RiverOrder()}, \code{sp.RiverDown()}, \code{RiverLength()},
and \code{RiverSlope()}. It uses sf and terra instead of sp and raster.
\end{Section}
%
\begin{Examples}
\begin{ExampleCode}
## Not run: 
library(sf)
library(terra)

# Load data
rivers <- st_read("rivers.shp")
dem <- rast("dem.tif")

# Calculate all properties
rivers_props <- calc_river_properties(rivers, dem)

# Calculate only specific properties
rivers_order <- calc_river_properties(rivers, properties = "order")

## End(Not run)
\end{ExampleCode}
\end{Examples}
\inputencoding{utf8}
\HeaderA{calc\_river\_width\_from\_area}{Calculate River Width from Watershed Area}{calc.Rul.river.Rul.width.Rul.from.Rul.area}
%
\begin{Description}\relax
Estimates river width for different stream orders based on watershed area
using an empirical relationship.
\end{Description}
%
\begin{Usage}
\begin{verbatim}
calc_river_width_from_area(area, n_types = 10)
\end{verbatim}
\end{Usage}
%
\begin{Arguments}
\begin{ldescription}
\item[\code{area}] Numeric, watershed area (in appropriate units)

\item[\code{n\_types}] Integer, number of river types/orders
\end{ldescription}
\end{Arguments}
%
\begin{Details}\relax
Uses an empirical power-law relationship to estimate channel width
based on drainage area. The relationship accounts for decreasing width
with increasing stream order (higher order = larger rivers).
\end{Details}
%
\begin{Value}
Numeric vector of widths for each type
\end{Value}
%
\begin{Examples}
\begin{ExampleCode}
# Calculate widths for 5 stream orders in a 1000 km² watershed
widths <- calc_river_width_from_area(1000, 5)
\end{ExampleCode}
\end{Examples}
\inputencoding{utf8}
\HeaderA{CanopyResistance}{Canopy resistance \code{CanopyResistance}}{CanopyResistance}
%
\begin{Description}\relax
Canopy resistance
\code{CanopyResistance}
\end{Description}
%
\begin{Usage}
\begin{verbatim}
CanopyResistance(rmin, lai)
\end{verbatim}
\end{Usage}
%
\begin{Arguments}
\begin{ldescription}
\item[\code{rmin}] Minimum Resistance \LinkA{s m-1}{s m.Rdash.1}
\end{ldescription}
\end{Arguments}
%
\begin{Value}
lai Leaf Area Index (LAI) \LinkA{m2 m-2}{m2 m.Rdash.2}
\end{Value}
\inputencoding{utf8}
\HeaderA{check\_file\_exists}{Check if a file exists}{check.Rul.file.Rul.exists}
%
\begin{Description}\relax
Validates that a file path exists and is accessible.
\end{Description}
%
\begin{Usage}
\begin{verbatim}
check_file_exists(path, name = "file", must_be_file = TRUE)
\end{verbatim}
\end{Usage}
%
\begin{Arguments}
\begin{ldescription}
\item[\code{path}] Character string, file path to check

\item[\code{name}] Character string, name of the parameter (for error messages)

\item[\code{must\_be\_file}] Logical, require path to be a file not directory
(default: TRUE)
\end{ldescription}
\end{Arguments}
%
\begin{Value}
Invisible TRUE if valid, otherwise stops with error
\end{Value}
%
\begin{Examples}
\begin{ExampleCode}
## Not run: 
check_file_exists("data/input.shp", "input_file")

## End(Not run)
\end{ExampleCode}
\end{Examples}
\inputencoding{utf8}
\HeaderA{check\_positive}{Check if a numeric value is positive}{check.Rul.positive}
%
\begin{Description}\relax
Validates that a numeric parameter is positive (greater than zero).
\end{Description}
%
\begin{Usage}
\begin{verbatim}
check_positive(x, name = "value", allow_zero = FALSE)
\end{verbatim}
\end{Usage}
%
\begin{Arguments}
\begin{ldescription}
\item[\code{x}] Numeric value to check

\item[\code{name}] Character string, name of the parameter (for error messages)

\item[\code{allow\_zero}] Logical, whether to allow zero (default: FALSE)
\end{ldescription}
\end{Arguments}
%
\begin{Value}
Invisible TRUE if valid, otherwise stops with error
\end{Value}
%
\begin{Examples}
\begin{ExampleCode}
## Not run: 
check_positive(5, "tolerance")
check_positive(0, "count", allow_zero = TRUE)

## End(Not run)
\end{ExampleCode}
\end{Examples}
\inputencoding{utf8}
\HeaderA{check\_spatial\_compatible}{Check if spatial objects are compatible}{check.Rul.spatial.Rul.compatible}
%
\begin{Description}\relax
Validates that two spatial objects have compatible properties
(e.g., same CRS, overlapping extents).
\end{Description}
%
\begin{Usage}
\begin{verbatim}
check_spatial_compatible(
  x,
  y,
  check_crs = TRUE,
  check_overlap = FALSE,
  name_x = "x",
  name_y = "y"
)
\end{verbatim}
\end{Usage}
%
\begin{Arguments}
\begin{ldescription}
\item[\code{x}] First spatial object (SpatRaster, SpatVector, or sf)

\item[\code{y}] Second spatial object (SpatRaster, SpatVector, or sf)

\item[\code{check\_crs}] Logical, check if CRS match (default: TRUE)

\item[\code{check\_overlap}] Logical, check if extents overlap (default: FALSE)

\item[\code{name\_x}] Character string, name of first object (for error messages)

\item[\code{name\_y}] Character string, name of second object (for error messages)
\end{ldescription}
\end{Arguments}
%
\begin{Value}
Invisible TRUE if valid, otherwise stops with error
\end{Value}
%
\begin{Examples}
\begin{ExampleCode}
## Not run: 
if (requireNamespace("terra", quietly = TRUE)) {
  r1 <- terra::rast(ncol=10, nrow=10)
  r2 <- terra::rast(ncol=10, nrow=10)
  terra::crs(r1) <- terra::crs(r2) <- "EPSG:4326"
  check_spatial_compatible(r1, r2)
}

## End(Not run)
\end{ExampleCode}
\end{Examples}
\inputencoding{utf8}
\HeaderA{CommonPoints}{Find the common coordinates in two groups of coordinates.}{CommonPoints}
%
\begin{Description}\relax
Find the common coordinates in two groups of coordinates.
\end{Description}
%
\begin{Usage}
\begin{verbatim}
CommonPoints(x, y)
\end{verbatim}
\end{Usage}
%
\begin{Arguments}
\begin{ldescription}
\item[\code{x, y}] Matrix of coodinates.
\end{ldescription}
\end{Arguments}
%
\begin{Value}
Matrix (m*2). Columns are index of x rows and y rows.
\end{Value}
\inputencoding{utf8}
\HeaderA{compareMaps}{Plot multiple maps \code{compareMaps}}{compareMaps}
%
\begin{Description}\relax
Plot multiple maps
\code{compareMaps}
\end{Description}
%
\begin{Usage}
\begin{verbatim}
compareMaps(r, mfrow, contour = FALSE, ...)
\end{verbatim}
\end{Usage}
%
\begin{Arguments}
\begin{ldescription}
\item[\code{r}] List of raster or SpatialData

\item[\code{mfrow}] mfrow in par()

\item[\code{contour}] Whether plot the contour.

\item[\code{...}] More options in par()
\end{ldescription}
\end{Arguments}
%
\begin{Examples}
\begin{ExampleCode}
library(raster)
data(volcano)
r <- raster(volcano)
extent(r) <- c(0, 610, 0, 870)
r1 = sin(r / 100)
r2 = cos(r / 100)
compareMaps(list(r, r1), mfrow = c(1, 2))
compareMaps(list(r, r1, r2, r1 + r2), mfrow = c(2, 2))
compareMaps(list(r, r1, r2, r1 + r2), mfrow = c(2, 2), contour = TRUE)
\end{ExampleCode}
\end{Examples}
\inputencoding{utf8}
\HeaderA{compatible\_crs}{Check if two CRS are compatible}{compatible.Rul.crs}
%
\begin{Description}\relax
Determines if two coordinate reference systems are compatible.
Handles different CRS object types (character, CRS, crs, etc.).
\end{Description}
%
\begin{Usage}
\begin{verbatim}
compatible_crs(crs1, crs2)
\end{verbatim}
\end{Usage}
%
\begin{Arguments}
\begin{ldescription}
\item[\code{crs1}] First CRS (character, CRS object, or crs object)

\item[\code{crs2}] Second CRS (character, CRS object, or crs object)
\end{ldescription}
\end{Arguments}
%
\begin{Value}
Logical, TRUE if CRS are compatible
\end{Value}
%
\begin{Examples}
\begin{ExampleCode}
## Not run: 
compatible_crs("EPSG:4326", "EPSG:4326")
compatible_crs("+proj=longlat", "+proj=longlat +datum=WGS84")

## End(Not run)
\end{ExampleCode}
\end{Examples}
\inputencoding{utf8}
\HeaderA{correctRiverSlope}{Convert the bed slope of river, to avoid negative/zero slope \code{correctRiverSlope}}{correctRiverSlope}
%
\begin{Description}\relax
Convert the bed slope of river, to avoid negative/zero slope
\code{correctRiverSlope}
\end{Description}
%
\begin{Usage}
\begin{verbatim}
correctRiverSlope(pr, minSlope = 1e-05, maxrun = 500)
\end{verbatim}
\end{Usage}
%
\begin{Arguments}
\begin{ldescription}
\item[\code{pr}] SHUD.RIVER class

\item[\code{minSlope}] Minimum slope

\item[\code{maxrun}] Maximum number to run the loops
\end{ldescription}
\end{Arguments}
%
\begin{Value}
SpatialLinesDataFrame
\end{Value}
\inputencoding{utf8}
\HeaderA{count}{Count the frequency of numbers \code{count}}{count}
%
\begin{Description}\relax
Count the frequency of numbers
\code{count}
\end{Description}
%
\begin{Usage}
\begin{verbatim}
count(x, n = NULL)
\end{verbatim}
\end{Usage}
%
\begin{Arguments}
\begin{ldescription}
\item[\code{x}] target value

\item[\code{n}] specific frequency
\end{ldescription}
\end{Arguments}
%
\begin{Value}
count of specific frequency
\end{Value}
%
\begin{Examples}
\begin{ExampleCode}
x=c(round(rnorm(100, 2)+1), 0,0)
count(x)
count(x, sum(x==0))
\end{ExampleCode}
\end{Examples}
\inputencoding{utf8}
\HeaderA{crs.Albers}{Build an Albers Equal Area projection}{crs.Albers}
%
\begin{Description}\relax
Creates an Albers Equal Area Conic projection based on spatial data extent.
This projection is suitable for areas with east-west extent and preserves
area measurements.
\end{Description}
%
\begin{Usage}
\begin{verbatim}
crs.Albers(spx = NULL, ext = NULL)
\end{verbatim}
\end{Usage}
%
\begin{Arguments}
\begin{ldescription}
\item[\code{spx}] Spatial data object (sf, SpatVector, or SpatRaster). If NULL,
must provide ext parameter.

\item[\code{ext}] Numeric vector of extent: c(xmin, xmax, ymin, ymax) in
longitude/latitude (EPSG:4326). If NULL, extracted from spx.
\end{ldescription}
\end{Arguments}
%
\begin{Value}
Character string with PROJ.4 CRS definition for Albers projection
\end{Value}
%
\begin{Examples}
\begin{ExampleCode}
# Create Albers projection from extent
crs_albers <- crs.Albers(ext = c(-120, -110, 35, 45))

# Use with spatial data
if (requireNamespace("sf", quietly = TRUE)) {
  pts <- sf::st_as_sf(
    data.frame(x = c(-115, -112), y = c(40, 42)),
    coords = c("x", "y"),
    crs = 4326
  )
  crs_albers <- crs.Albers(spx = pts)
}
\end{ExampleCode}
\end{Examples}
\inputencoding{utf8}
\HeaderA{crs.Lambert}{Build a Lambert Equal Area projection}{crs.Lambert}
%
\begin{Description}\relax
Creates a Lambert Azimuthal Equal Area projection based on spatial data
extent. This projection is suitable for areas with similar north-south and
east-west extent and preserves area measurements.
\end{Description}
%
\begin{Usage}
\begin{verbatim}
crs.Lambert(spx = NULL, ext = NULL)
\end{verbatim}
\end{Usage}
%
\begin{Arguments}
\begin{ldescription}
\item[\code{spx}] Spatial data object (sf, SpatVector, or SpatRaster). If NULL,
must provide ext parameter.

\item[\code{ext}] Numeric vector of extent: c(xmin, xmax, ymin, ymax) in
longitude/latitude (EPSG:4326). If NULL, extracted from spx.
\end{ldescription}
\end{Arguments}
%
\begin{Value}
Character string with PROJ.4 CRS definition for Lambert projection
\end{Value}
%
\begin{Examples}
\begin{ExampleCode}
# Create Lambert projection from extent
crs_lambert <- crs.Lambert(ext = c(-120, -110, 35, 45))

# Use with spatial data
if (requireNamespace("sf", quietly = TRUE)) {
  pts <- sf::st_as_sf(
    data.frame(x = c(-115, -112), y = c(40, 42)),
    coords = c("x", "y"),
    crs = 4326
  )
  crs_lambert <- crs.Lambert(spx = pts)
}
\end{ExampleCode}
\end{Examples}
\inputencoding{utf8}
\HeaderA{crs.long2utm}{Get UTM projection from longitude and latitude}{crs.long2utm}
%
\begin{Description}\relax
Creates a UTM (Universal Transverse Mercator) projection CRS string based
on longitude and latitude coordinates.
\end{Description}
%
\begin{Usage}
\begin{verbatim}
crs.long2utm(lon, lat = 30)
\end{verbatim}
\end{Usage}
%
\begin{Arguments}
\begin{ldescription}
\item[\code{lon}] Numeric, longitude in degrees (-180 to 180)

\item[\code{lat}] Numeric, latitude in degrees (-90 to 90). Default is 30 (Northern
Hemisphere). Use negative values for Southern Hemisphere.
\end{ldescription}
\end{Arguments}
%
\begin{Value}
Character string with PROJ.4 CRS definition for UTM projection.
If lon is a vector, returns a list of CRS strings.
\end{Value}
%
\begin{Examples}
\begin{ExampleCode}
# Single location (Northern Hemisphere)
crs_utm <- crs.long2utm(-75, 40)

# Southern Hemisphere
crs_utm_south <- crs.long2utm(150, -30)

# Multiple longitudes
lons <- seq(-180, 179, 6) + 1
crs_list <- crs.long2utm(lons)
\end{ExampleCode}
\end{Examples}
\inputencoding{utf8}
\HeaderA{crs.long2utmZone}{Get UTM zone from longitude}{crs.long2utmZone}
%
\begin{Description}\relax
Calculates the UTM zone number from longitude coordinate(s).
\end{Description}
%
\begin{Usage}
\begin{verbatim}
crs.long2utmZone(lon)
\end{verbatim}
\end{Usage}
%
\begin{Arguments}
\begin{ldescription}
\item[\code{lon}] Numeric vector of longitude values in degrees (-180 to 180)
\end{ldescription}
\end{Arguments}
%
\begin{Value}
Integer vector of UTM zone numbers (1-60)
\end{Value}
%
\begin{Examples}
\begin{ExampleCode}
# Single longitude
crs.long2utmZone(-75)  # Returns 18

# Multiple longitudes
lons <- seq(-180, 179, 6) + 1
zones <- crs.long2utmZone(lons)
\end{ExampleCode}
\end{Examples}
\inputencoding{utf8}
\HeaderA{datafilter.riv}{Filter the river data with a filter \code{datafilter.riv}}{datafilter.riv}
%
\begin{Description}\relax
Filter the river data with a filter
\code{datafilter.riv}
\end{Description}
%
\begin{Usage}
\begin{verbatim}
datafilter.riv(x, filter = NULL, plot = TRUE)
\end{verbatim}
\end{Usage}
%
\begin{Arguments}
\begin{ldescription}
\item[\code{x}] Input data

\item[\code{filter}] Data filter

\item[\code{plot}] Whether plot the data
\end{ldescription}
\end{Arguments}
%
\begin{Value}
Matrix information, c('ID', 'Vmin', 'Vmax', 'Filter')
\end{Value}
\inputencoding{utf8}
\HeaderA{days\_in\_year}{Number of days in years.}{days.Rul.in.Rul.year}
%
\begin{Description}\relax
Number of days in years.
\end{Description}
%
\begin{Usage}
\begin{verbatim}
days_in_year(years)
\end{verbatim}
\end{Usage}
%
\begin{Arguments}
\begin{ldescription}
\item[\code{years}] Years
\end{ldescription}
\end{Arguments}
%
\begin{Value}
Number of days between first and last years.
\end{Value}
%
\begin{Examples}
\begin{ExampleCode}
days_in_year(1980)
days_in_year(1983:2019)
\end{ExampleCode}
\end{Examples}
\inputencoding{utf8}
\HeaderA{DeltaS}{Calculate the Change of Storage. \code{DeltaS}}{DeltaS}
%
\begin{Description}\relax
Calculate the Change of Storage.
\code{DeltaS}
\end{Description}
%
\begin{Usage}
\begin{verbatim}
DeltaS(x, x0 = x[t1, ], t1 = 1, t2 = nrow(x))
\end{verbatim}
\end{Usage}
%
\begin{Arguments}
\begin{ldescription}
\item[\code{x}] Time-Series Matrix

\item[\code{x0}] Intial condition or the values of first time-step

\item[\code{t1}] Time-step one, default = 1

\item[\code{t2}] Time-step two, default = nrow(x)
\end{ldescription}
\end{Arguments}
%
\begin{Value}
A vector
\end{Value}
\inputencoding{utf8}
\HeaderA{Eudist}{Eudilic Distance \code{Eudist}}{Eudist}
%
\begin{Description}\relax
Eudilic Distance
\code{Eudist}
\end{Description}
%
\begin{Usage}
\begin{verbatim}
Eudist(p1, p2)
\end{verbatim}
\end{Usage}
%
\begin{Arguments}
\begin{ldescription}
\item[\code{p1}] coordinates of point 1

\item[\code{p2}] coordinates of point 2
\end{ldescription}
\end{Arguments}
%
\begin{Value}
Distance between \code{p1} and \code{p2}
\end{Value}
%
\begin{Examples}
\begin{ExampleCode}
p1 = c(1,1)
p2 = c(2,2)
Eudist(p1,p2)
\end{ExampleCode}
\end{Examples}
\inputencoding{utf8}
\HeaderA{extractCoords}{Extract Coordinates of SpatialLines or SpatialPolygons \code{extractCoords}}{extractCoords}
%
\begin{Description}\relax
Extract Coordinates of SpatialLines or SpatialPolygons
\code{extractCoords}
\end{Description}
%
\begin{Usage}
\begin{verbatim}
extractCoords(x, unique = TRUE, aslist = FALSE)
\end{verbatim}
\end{Usage}
%
\begin{Arguments}
\begin{ldescription}
\item[\code{x}] SpatialLines or SpatialPolygons

\item[\code{unique}] Whether return unique coordinates or duplicated coordinates

\item[\code{aslist}] Whether return as list of coordinates
\end{ldescription}
\end{Arguments}
%
\begin{Value}
coordinates of points, m x 2.
\end{Value}
\inputencoding{utf8}
\HeaderA{extractRaster}{Extract values on Raster map. The line is a straight line between (0,1). \code{extractRaster}}{extractRaster}
%
\begin{Description}\relax
Extract values on Raster map. The line is a straight line between (0,1).
\code{extractRaster}
\end{Description}
%
\begin{Usage}
\begin{verbatim}
extractRaster(r, xy = NULL, ext = raster::extent(r), plot = TRUE)
\end{verbatim}
\end{Usage}
%
\begin{Arguments}
\begin{ldescription}
\item[\code{r}] Raster

\item[\code{xy}] coordinates of the line, dim = (Npoints, 2); x and y must be in \LinkA{0, 1}{0, 1}

\item[\code{ext}] extension of value xy.

\item[\code{plot}] Whether plot result.
\end{ldescription}
\end{Arguments}
%
\begin{Examples}
\begin{ExampleCode}
library(raster)
\end{ExampleCode}
\end{Examples}
\inputencoding{utf8}
\HeaderA{fdc}{Default Comparison of the observation and simulation \code{fdc}}{fdc}
%
\begin{Description}\relax
Default Comparison of the observation and simulation
\code{fdc}
\end{Description}
%
\begin{Usage}
\begin{verbatim}
fdc(val, xlab = "Exceedance (%)", ylab = "value")
\end{verbatim}
\end{Usage}
%
\begin{Arguments}
\begin{ldescription}
\item[\code{val}] Matrix or data.frame, number of column = 2

\item[\code{xlab}] xlab

\item[\code{ylab}] ylab
\end{ldescription}
\end{Arguments}
%
\begin{Value}
ggplot handler
\end{Value}
%
\begin{Examples}
\begin{ExampleCode}
obs = rnorm(100)
sim = obs + rnorm(100) / 2
fdc(cbind(obs, sim))
\end{ExampleCode}
\end{Examples}
\inputencoding{utf8}
\HeaderA{fishnet}{Generatue fishnet \code{fishnet} Generate fishnet by the coordinates}{fishnet}
%
\begin{Description}\relax
Generatue fishnet
\code{fishnet} Generate fishnet by the coordinates
\end{Description}
%
\begin{Usage}
\begin{verbatim}
fishnet(xx, yy, crs = sp::CRS("+init=epsg:4326"), type = "polygon")
\end{verbatim}
\end{Usage}
%
\begin{Arguments}
\begin{ldescription}
\item[\code{xx}] x coordinates

\item[\code{yy}] y coordinates

\item[\code{crs}] projections parameters, defaul = epsg:4326

\item[\code{type}] option = 'polygon', 'points', 'line'
\end{ldescription}
\end{Arguments}
%
\begin{Value}
spatildata (.shp)
\end{Value}
%
\begin{Examples}
\begin{ExampleCode}
library(raster)
ext = c(0, 8, 2, 10)
dx = 2; dy = 4
xx = seq(ext[1], ext[2], by = dx)
yy = seq(ext[3], ext[4], by = dy)
sp1 = fishnet(xx = ext[1:2], yy = ext[3:4])
sp2 = fishnet(xx = xx + .5 * dx, yy = yy + 0.5 * dy)
sp3 = fishnet(xx = xx, yy = yy, type = 'point')
plot(sp1, axes = TRUE, xlim = c(-1, 1) * dx + ext[1:2], ylim = c(-1, 1) * dy + ext[3:4])
plot(sp2, axes = TRUE, add = TRUE, border = 2)
plot(sp3, axes = TRUE, add = TRUE, col = 3, pch = 20)
grid()
\end{ExampleCode}
\end{Examples}
\inputencoding{utf8}
\HeaderA{ForcingCoverage}{Generate the coverage map for forcing sites. \code{ForcingCoverage}}{ForcingCoverage}
%
\begin{Description}\relax
Generate the coverage map for forcing sites.
\code{ForcingCoverage}
\end{Description}
%
\begin{Usage}
\begin{verbatim}
ForcingCoverage(
  sp.meteoSite = NULL,
  filenames = paste0(sp.meteoSite@data$ID, ".csv"),
  pcs,
  gcs = sp::CRS("+init=epsg:4326"),
  dem,
  wbd,
  enlarge = 10000
)
\end{verbatim}
\end{Usage}
%
\begin{Arguments}
\begin{ldescription}
\item[\code{sp.meteoSite}] ShapePoints* in PCS

\item[\code{pcs}] Projected Coordinate System

\item[\code{gcs}] Geographic Coordinate System

\item[\code{dem}] DEM raster

\item[\code{wbd}] watershed boundary

\item[\code{enlarge}] enlarge factor for the boundary.
\end{ldescription}
\end{Arguments}
%
\begin{Value}
ShapePolygon*
\end{Value}
\inputencoding{utf8}
\HeaderA{FromToNode}{return the FROM and TO nodes index of the SpatialLines \code{FromToNode}}{FromToNode}
%
\begin{Description}\relax
return the FROM and TO nodes index of the SpatialLines
\code{FromToNode}
\end{Description}
%
\begin{Usage}
\begin{verbatim}
FromToNode(sp, coord = extractCoords(sp, unique = TRUE), simplify = TRUE)
\end{verbatim}
\end{Usage}
%
\begin{Arguments}
\begin{ldescription}
\item[\code{sp}] SpatialLines*

\item[\code{coord}] Coordinate of vertex in \code{sp}

\item[\code{simplify}] whether simplify the spatialLines
\end{ldescription}
\end{Arguments}
%
\begin{Value}
FROM and TO nodes index of the SpatialLines
\end{Value}
\inputencoding{utf8}
\HeaderA{generate\_river\_types}{Generate River Type Parameters}{generate.Rul.river.Rul.types}
%
\begin{Description}\relax
Creates a data frame of hydraulic parameters for different river types/orders.
This is used to assign width, depth, and other hydraulic properties based
on stream order or other classification.
\end{Description}
%
\begin{Usage}
\begin{verbatim}
generate_river_types(
  n,
  width = 2 * (1:n),
  depth = 5 + 0.5 * (1:n),
  manning = 0.04
)
\end{verbatim}
\end{Usage}
%
\begin{Arguments}
\begin{ldescription}
\item[\code{n}] Integer, number of river types

\item[\code{width}] Numeric vector of widths for each type (default: 2 * (1:n))

\item[\code{depth}] Numeric vector of depths (default: 5.0 + 0.5 * (1:n))

\item[\code{manning}] Numeric, Manning's roughness coefficient (default: 0.04)
\end{ldescription}
\end{Arguments}
%
\begin{Details}\relax
The function generates default hydraulic parameters that scale with river
type/order. Users can customize width and depth, while other parameters
use reasonable defaults.

Parameters included:
\begin{itemize}

\item{} Index: River type index
\item{} Depth: Channel depth (m)
\item{} BankSlope: Bank slope (default: 0)
\item{} Width: Channel width (m)
\item{} Sinuosity: Channel sinuosity (default: 1.0)
\item{} Manning: Manning's n (s/m\textasciicircum{}(1/3))
\item{} Cwr: Coefficient for width-discharge relationship (default: 0.6)
\item{} KsatH: Horizontal hydraulic conductivity (default: 0.1)
\item{} BedThick: Bed sediment thickness (default: 0.1)

\end{itemize}

\end{Details}
%
\begin{Value}
data.frame with hydraulic parameters for each river type
\end{Value}
%
\begin{Section}{Migration Note}

This function replaces \code{RiverType()} with the same interface.
\end{Section}
%
\begin{Examples}
\begin{ExampleCode}
# Generate parameters for 5 river types
river_types <- generate_river_types(5)

# Custom widths based on watershed area
widths <- c(2, 5, 10, 20, 40)
river_types <- generate_river_types(5, width = widths)
\end{ExampleCode}
\end{Examples}
\inputencoding{utf8}
\HeaderA{get\_default\_variables}{Get default output variables}{get.Rul.default.Rul.variables}
\keyword{output-io}{get\_default\_variables}
%
\begin{Description}\relax
Returns a vector of commonly used SHUD output variables.
\end{Description}
%
\begin{Usage}
\begin{verbatim}
get_default_variables()
\end{verbatim}
\end{Usage}
%
\begin{Value}
Character vector of variable names
\end{Value}
%
\begin{SeeAlso}\relax
Other output-io: 
\code{\LinkA{load\_output\_data}{load.Rul.output.Rul.data}()},
\code{\LinkA{read\_output\_header}{read.Rul.output.Rul.header}()},
\code{\LinkA{read\_output}{read.Rul.output}()}
\end{SeeAlso}
\inputencoding{utf8}
\HeaderA{get\_nc\_info}{Get NetCDF file information}{get.Rul.nc.Rul.info}
\keyword{netcdf-io}{get\_nc\_info}
%
\begin{Description}\relax
Extracts basic information about NetCDF file structure.
\end{Description}
%
\begin{Usage}
\begin{verbatim}
get_nc_info(ncid)
\end{verbatim}
\end{Usage}
%
\begin{Arguments}
\begin{ldescription}
\item[\code{ncid}] NetCDF connection object
\end{ldescription}
\end{Arguments}
%
\begin{Value}
List with file information
\end{Value}
%
\begin{SeeAlso}\relax
Other netcdf-io: 
\code{\LinkA{nc\_to\_raster}{nc.Rul.to.Rul.raster}()},
\code{\LinkA{open\_nc\_file}{open.Rul.nc.Rul.file}()},
\code{\LinkA{read\_nc\_data}{read.Rul.nc.Rul.data}()},
\code{\LinkA{read\_nc\_time}{read.Rul.nc.Rul.time}()}
\end{SeeAlso}
\inputencoding{utf8}
\HeaderA{get\_river\_outlets}{Get River Outlets}{get.Rul.river.Rul.outlets}
%
\begin{Description}\relax
Identifies outlet segments in a river network (segments with no downstream).
\end{Description}
%
\begin{Usage}
\begin{verbatim}
get_river_outlets(rivers)
\end{verbatim}
\end{Usage}
%
\begin{Arguments}
\begin{ldescription}
\item[\code{rivers}] sf object with river network, or integer vector of downstream indices
\end{ldescription}
\end{Arguments}
%
\begin{Details}\relax
Outlets are identified as segments where the downstream index is negative
(typically -1).
\end{Details}
%
\begin{Value}
Integer vector of outlet segment indices
\end{Value}
%
\begin{Section}{Migration Note}

This function replaces \code{getOutlets()} with a simpler interface.
\end{Section}
%
\begin{Examples}
\begin{ExampleCode}
## Not run: 
library(sf)
rivers <- st_read("rivers.shp")
downstream <- calc_river_downstream(rivers)
outlets <- get_river_outlets(downstream)

## End(Not run)
\end{ExampleCode}
\end{Examples}
\inputencoding{utf8}
\HeaderA{getAquiferDepth}{Calculate the AquiferDepth of mesh cells \code{getAquiferDepth}}{getAquiferDepth}
%
\begin{Description}\relax
Calculate the AquiferDepth of mesh cells
\code{getAquiferDepth}
\end{Description}
%
\begin{Usage}
\begin{verbatim}
getAquiferDepth(pm = readmesh())
\end{verbatim}
\end{Usage}
%
\begin{Arguments}
\begin{ldescription}
\item[\code{pm}] \code{shud.mesh}
\end{ldescription}
\end{Arguments}
%
\begin{Value}
aquifer thickness
\end{Value}
\inputencoding{utf8}
\HeaderA{getArea}{Calculate the area of the cells \code{getArea}}{getArea}
%
\begin{Description}\relax
Calculate the area of the cells
\code{getArea}
\end{Description}
%
\begin{Usage}
\begin{verbatim}
getArea(pm = readmesh())
\end{verbatim}
\end{Usage}
%
\begin{Arguments}
\begin{ldescription}
\item[\code{pm}] \code{shud.mesh}
\end{ldescription}
\end{Arguments}
%
\begin{Value}
Area of cells
\end{Value}
\inputencoding{utf8}
\HeaderA{getCentroid}{Get the centroids of the cells \code{getCentroid}}{getCentroid}
%
\begin{Description}\relax
Get the centroids of the cells
\code{getCentroid}
\end{Description}
%
\begin{Usage}
\begin{verbatim}
getCentroid(pm = readmesh())
\end{verbatim}
\end{Usage}
%
\begin{Arguments}
\begin{ldescription}
\item[\code{pm}] \code{shud.mesh}
\end{ldescription}
\end{Arguments}
%
\begin{Value}
centroids of \code{shud.mesh}, ncell x 4 c('x','y','AqD', 'zmax')
\end{Value}
\inputencoding{utf8}
\HeaderA{getElevation}{Calculate the elevation of mesh cells \code{getElevation}}{getElevation}
%
\begin{Description}\relax
Calculate the elevation of mesh cells
\code{getElevation}
\end{Description}
%
\begin{Usage}
\begin{verbatim}
getElevation(pm = readmesh())
\end{verbatim}
\end{Usage}
%
\begin{Arguments}
\begin{ldescription}
\item[\code{pm}] \code{shud.mesh}
\end{ldescription}
\end{Arguments}
%
\begin{Value}
elevation
\end{Value}
\inputencoding{utf8}
\HeaderA{getLakeEleID}{Get the Lake Element Index \code{getLakeEleID}}{getLakeEleID}
%
\begin{Description}\relax
Get the Lake Element Index
\code{getLakeEleID}
\end{Description}
%
\begin{Usage}
\begin{verbatim}
getLakeEleID(pa = readatt())
\end{verbatim}
\end{Usage}
%
\begin{Arguments}
\begin{ldescription}
\item[\code{pa}] \code{shud.att}
\end{ldescription}
\end{Arguments}
%
\begin{Value}
index of lake element index
\end{Value}
\inputencoding{utf8}
\HeaderA{getOutlets}{Get outlet ID(s) \code{getOutlets}}{getOutlets}
%
\begin{Description}\relax
Get outlet ID(s)
\code{getOutlets}
\end{Description}
%
\begin{Usage}
\begin{verbatim}
getOutlets(pr = readriv())
\end{verbatim}
\end{Usage}
%
\begin{Arguments}
\begin{ldescription}
\item[\code{pr}] SHUD.RIVER class
\end{ldescription}
\end{Arguments}
%
\begin{Value}
Index of outlets, numeric
\end{Value}
\inputencoding{utf8}
\HeaderA{getRiverNodes}{Get the From/To nodes of the river. \code{getRiverNodes}}{getRiverNodes}
%
\begin{Description}\relax
Get the From/To nodes of the river.
\code{getRiverNodes}
\end{Description}
%
\begin{Usage}
\begin{verbatim}
getRiverNodes(spr = readriv.sp())
\end{verbatim}
\end{Usage}
%
\begin{Arguments}
\begin{ldescription}
\item[\code{spr}] SpatialLine* of river streams.
\end{ldescription}
\end{Arguments}
%
\begin{Value}
a list, c(points, FT\_ID)
\end{Value}
\inputencoding{utf8}
\HeaderA{getVertex}{Calculate the Vertex of the cells \code{getVertex}}{getVertex}
%
\begin{Description}\relax
Calculate the Vertex of the cells
\code{getVertex}
\end{Description}
%
\begin{Usage}
\begin{verbatim}
getVertex(pm = readmesh())
\end{verbatim}
\end{Usage}
%
\begin{Arguments}
\begin{ldescription}
\item[\code{pm}] \code{shud.mesh}
\end{ldescription}
\end{Arguments}
%
\begin{Value}
Vertex of \code{shud.mesh}, dim = c(ncell, vertex, 4), 3-dimension = c('x','y','AqD', 'zmax')
\end{Value}
\inputencoding{utf8}
\HeaderA{grid.subset}{Find the subset of a extent in a grid. \code{grid.subset}}{grid.subset}
%
\begin{Description}\relax
Find the subset of a extent in a grid.
\code{grid.subset}
\end{Description}
%
\begin{Usage}
\begin{verbatim}
grid.subset(
  ext,
  res,
  ext.sub,
  dx = matrix(res, 2, 1)[1],
  dy = matrix(res, 2, 1)[2],
  x = seq(ext[1] + dx/2, ext[2] - dx/2, by = dx),
  y = seq(ext[3] + dy/2, ext[4] - dy/2, by = dy)
)
\end{verbatim}
\end{Usage}
%
\begin{Arguments}
\begin{ldescription}
\item[\code{ext}] txtent of the grid

\item[\code{res}] resolution of the grid

\item[\code{ext.sub}] the extent of subset

\item[\code{x}] x coordinates of the grids

\item[\code{y}] y coordinates of the grids
\end{ldescription}
\end{Arguments}
%
\begin{Value}
list, list(xid, yid, x, y)
\end{Value}
\inputencoding{utf8}
\HeaderA{highlight\_id}{Highlight elements with id. \code{highlight\_id}}{highlight.Rul.id}
%
\begin{Description}\relax
Highlight elements with id.
\code{highlight\_id}
\end{Description}
%
\begin{Usage}
\begin{verbatim}
highlight_id(EleID = NULL, RivID = NULL)
\end{verbatim}
\end{Usage}
%
\begin{Arguments}
\begin{ldescription}
\item[\code{EleID}] Index of element

\item[\code{RivID}] Index of river reaches
\end{ldescription}
\end{Arguments}
\inputencoding{utf8}
\HeaderA{hydrograph}{Plot hydrograph \code{hydrograph}}{hydrograph}
%
\begin{Description}\relax
Plot hydrograph
\code{hydrograph}
\end{Description}
%
\begin{Usage}
\begin{verbatim}
hydrograph(
  x,
  legend.position = "bottom",
  unit = rep("", ncol(x)),
  col = c(3, 4),
  heights = c(3, 7),
  bg = TRUE,
  ylabs = NULL
)
\end{verbatim}
\end{Usage}
%
\begin{Arguments}
\begin{ldescription}
\item[\code{x}] time-seres matrix. The first column will be plot on top subfigure.

\item[\code{legend.position}] Location to put the legend for discharge.

\item[\code{unit}] Unit of the variables.

\item[\code{col}] colors of each variable.

\item[\code{heights}] Heights of top (rainfall) figure and bottom (discharge) figure

\item[\code{bg}] Whether plot the background rectangles.

\item[\code{ylabs}] ylab for top and bottom plot.
\end{ldescription}
\end{Arguments}
\inputencoding{utf8}
\HeaderA{is\_modern\_river}{Check if SHUD.RIVER Uses Modern Format}{is.Rul.modern.Rul.river}
%
\begin{Description}\relax
Determines if a SHUD.RIVER object uses the modern sf-based format
or the legacy data.frame-only format.
\end{Description}
%
\begin{Usage}
\begin{verbatim}
is_modern_river(shud_river)
\end{verbatim}
\end{Usage}
%
\begin{Arguments}
\begin{ldescription}
\item[\code{shud\_river}] SHUD.RIVER S4 object
\end{ldescription}
\end{Arguments}
%
\begin{Value}
Logical, TRUE if modern format
\end{Value}
\inputencoding{utf8}
\HeaderA{LAI2hc}{Calculate vegetation height from Leaf Area Index (LAI), Eq 4.2.23 in David R Maidment, Handbook of Hydrology \code{LAI2hc}}{LAI2hc}
%
\begin{Description}\relax
Calculate vegetation height from Leaf Area Index (LAI), Eq 4.2.23 in David R Maidment, Handbook of Hydrology
\code{LAI2hc}
\end{Description}
%
\begin{Usage}
\begin{verbatim}
LAI2hc(lai)
\end{verbatim}
\end{Usage}
%
\begin{Arguments}
\begin{ldescription}
\item[\code{lai}] Leaf Area Index (LAI) \LinkA{m2 m-2}{m2 m.Rdash.2}
\end{ldescription}
\end{Arguments}
%
\begin{Value}
Height of crop \LinkA{m}{m}
\end{Value}
%
\begin{Examples}
\begin{ExampleCode}
lai = seq(0, 10, length.out=100)
hc = LAI2hc(lai)
# lai=5.5+log(hc)*1.5
plot(lai, hc)
\end{ExampleCode}
\end{Examples}
\inputencoding{utf8}
\HeaderA{LaiRf.GLC}{Generate the default LAI and Roughness Length for USGS GLC \code{LaiRf.GLC}}{LaiRf.GLC}
%
\begin{Description}\relax
Generate the default LAI and Roughness Length for USGS GLC
\code{LaiRf.GLC}
\end{Description}
%
\begin{Usage}
\begin{verbatim}
LaiRf.GLC(lc = NULL, years = 2000 + 1:2, if.daily = FALSE)
\end{verbatim}
\end{Usage}
%
\begin{Arguments}
\begin{ldescription}
\item[\code{lc}] land classe codes

\item[\code{years}] numeric years.

\item[\code{if.daily}] whether daily output
\end{ldescription}
\end{Arguments}
%
\begin{Value}
Default LaiRf.GLC parameters, a list \$LAI, \$RL
\end{Value}
%
\begin{Examples}
\begin{ExampleCode}
lc = 1:14
lr = LaiRf.GLC(lc, years = 2000:2001)
par(mfrow = c(2, 1))
col = 1:length(lc)
plot(lr$LAI, col = col, main = 'LAI');
legend('top', paste0(lc), col = col, lwd = 1)
plot(lr$RL, col = col, main = 'Roughness Length');
legend('top', paste0(lc), col = col, lwd = 1)
\end{ExampleCode}
\end{Examples}
\inputencoding{utf8}
\HeaderA{LaiRf.NLCD}{Generate the default LAI and Roughness Length \code{LaiRf.NLCD}}{LaiRf.NLCD}
%
\begin{Description}\relax
Generate the default LAI and Roughness Length
\code{LaiRf.NLCD}
\end{Description}
%
\begin{Usage}
\begin{verbatim}
LaiRf.NLCD(lc, years = 2000)
\end{verbatim}
\end{Usage}
%
\begin{Arguments}
\begin{ldescription}
\item[\code{lc}] land classe codes

\item[\code{years}] numeric years.

\item[\code{class}] The classification of the Landuse of NLCD
\end{ldescription}
\end{Arguments}
%
\begin{Value}
Default LaiRf.NLCD parameters, a list \$LAI, \$RL
\end{Value}
%
\begin{Examples}
\begin{ExampleCode}
lc = c(43, 23, 81, 11)
lr = LaiRf.NLCD(lc, years = 2000:2001)
par(mfrow = c(2, 1))
col = 1:length(lc)
plot(lr$LAI, col = col, main = 'LAI');
legend('top', paste0(lc), col = col, lwd = 1)
plot(lr$RL, col = col, main = 'Roughness Length');
legend('top', paste0(lc), col = col, lwd = 1)
\end{ExampleCode}
\end{Examples}
\inputencoding{utf8}
\HeaderA{LatentHeat}{Latent Heat, Eq 4.2.1 in David R Maidment, Handbook of Hydrology \code{shud.ic}}{LatentHeat}
%
\begin{Description}\relax
Latent Heat, Eq 4.2.1 in David R Maidment, Handbook of Hydrology
\code{shud.ic}
\end{Description}
%
\begin{Usage}
\begin{verbatim}
LatentHeat(Temp)
\end{verbatim}
\end{Usage}
%
\begin{Arguments}
\begin{ldescription}
\item[\code{Temp}] Air temperature, \LinkA{C}{C}
\end{ldescription}
\end{Arguments}
%
\begin{Value}
Latent heat, lambda in \LinkA{MJ/kg}{MJ/kg}
\end{Value}
\inputencoding{utf8}
\HeaderA{lc.GLC}{The default land cover parameters from USGS GLC \code{lc.GLC}}{lc.GLC}
%
\begin{Description}\relax
The default land cover parameters from USGS GLC
\code{lc.GLC}
\end{Description}
%
\begin{Usage}
\begin{verbatim}
lc.GLC()
\end{verbatim}
\end{Usage}
%
\begin{Value}
Default land cover parameters
\end{Value}
\inputencoding{utf8}
\HeaderA{lc.NLCD}{Generate the default land cover parameters of NLCD classes. \code{PTF.NLCD}}{lc.NLCD}
%
\begin{Description}\relax
Generate the default land cover parameters of NLCD classes.
\code{PTF.NLCD}
\end{Description}
%
\begin{Usage}
\begin{verbatim}
lc.NLCD(lc)
\end{verbatim}
\end{Usage}
%
\begin{Arguments}
\begin{ldescription}
\item[\code{lc}] NLCD (2001 and later version) land use code.
\end{ldescription}
\end{Arguments}
%
\begin{Value}
Default land cover parameters of NLCD classes.
\end{Value}
\inputencoding{utf8}
\HeaderA{lc.UMD}{The default land cover parameters from UMD \code{lc.UMD}}{lc.UMD}
%
\begin{Description}\relax
The default land cover parameters from UMD
\code{lc.UMD}
\end{Description}
%
\begin{Usage}
\begin{verbatim}
lc.UMD()
\end{verbatim}
\end{Usage}
%
\begin{Value}
Default land cover parameters
\end{Value}
\inputencoding{utf8}
\HeaderA{lc\_EQ}{convert equation from Univeristy of Marryland classes to NLCD classes \code{lc\_EQ}}{lc.Rul.EQ}
%
\begin{Description}\relax
convert equation from Univeristy of Marryland classes to NLCD classes
\code{lc\_EQ}
\end{Description}
%
\begin{Usage}
\begin{verbatim}
lc_EQ(dtab, lc)
\end{verbatim}
\end{Usage}
%
\begin{Arguments}
\begin{ldescription}
\item[\code{dtab}] conversion table.

\item[\code{lc}] land classe codes in NLCD classes.
\end{ldescription}
\end{Arguments}
%
\begin{Value}
Converted values.
\end{Value}
%
\begin{References}\relax
Lele Shu(2017), Gopal Bhatt(2012)
\end{References}
\inputencoding{utf8}
\HeaderA{LineFit}{Linefit the observation and simulation \code{LineFit}}{LineFit}
%
\begin{Description}\relax
Linefit the observation and simulation
\code{LineFit}
\end{Description}
%
\begin{Usage}
\begin{verbatim}
LineFit(x, y = NULL, xlab = "Observation", ylab = "Simulation", ...)
\end{verbatim}
\end{Usage}
%
\begin{Arguments}
\begin{ldescription}
\item[\code{x}] Observation data or matrix/data.frame consist of observation and simulation

\item[\code{y}] Simulation data

\item[\code{xlab}] xlab, default is 'Observation'

\item[\code{ylab}] ylab, default is 'Simulation'

\item[\code{...}] more options for stat\_poly\_eq()
\end{ldescription}
\end{Arguments}
%
\begin{Value}
Triangle mesh of model domain.
\end{Value}
%
\begin{Examples}
\begin{ExampleCode}
obs <- rnorm(100)
sim <- obs + rnorm(100) / 2
p <- LineFit(cbind(obs, sim))
p
if(require(ggpmisc)){
formula = y ~ x
p + ggpmisc::stat_poly_eq(aes(label = paste(..eq.label.., ..rr.label.., sep = "~~~")), 
label.x.npc = "right", label.y.npc = 0.15,
formula = formula, parse = TRUE, size = 3)
}
\end{ExampleCode}
\end{Examples}
\inputencoding{utf8}
\HeaderA{load\_output\_data}{Load multiple SHUD output files}{load.Rul.output.Rul.data}
\keyword{output-io}{load\_output\_data}
%
\begin{Description}\relax
Loads multiple output variables and returns them as a named list.
Supports both modern and legacy formats.
\end{Description}
%
\begin{Usage}
\begin{verbatim}
load_output_data(
  variables = get_default_variables(),
  format = c("modern", "legacy"),
  save_rds = NULL,
  version = get_model_version()
)
\end{verbatim}
\end{Usage}
%
\begin{Arguments}
\begin{ldescription}
\item[\code{variables}] Character vector. Output variable keywords to load

\item[\code{format}] Character. Output format ("modern" or "legacy")

\item[\code{save\_rds}] Character. Path to save RDS file (optional)

\item[\code{version}] Numeric. Output file version
\end{ldescription}
\end{Arguments}
%
\begin{Value}
Named list of xts objects
\end{Value}
%
\begin{SeeAlso}\relax
Other output-io: 
\code{\LinkA{get\_default\_variables}{get.Rul.default.Rul.variables}()},
\code{\LinkA{read\_output\_header}{read.Rul.output.Rul.header}()},
\code{\LinkA{read\_output}{read.Rul.output}()}
\end{SeeAlso}
%
\begin{Examples}
\begin{ExampleCode}
## Not run: 
# Load common variables
vars <- c("eleygw", "eleysurf", "rivqdown")
data_list <- load_output_data(vars)

# Use modern format
data_list <- load_output_data(vars, format = "modern")

## End(Not run)
\end{ExampleCode}
\end{Examples}
\inputencoding{utf8}
\HeaderA{loaddata}{Read multiply SHUD model output files}{loaddata}
\keyword{output.}{loaddata}
\keyword{read}{loaddata}
%
\begin{Description}\relax
Read multiply SHUD model output files
\end{Description}
%
\begin{Usage}
\begin{verbatim}
loaddata(
  varname = c(paste0("eley", c("surf", "unsat", "gw", "snow")), paste0("elev", c("prcp",
    "infil", "rech")), paste0("elev", c("etp", "eta", "etev", "ettr", "etic")),
    paste0("rivq", c("down", "sub", "surf")), paste0("rivy", "stage")),
  rdsfile = NULL,
  ver = get("version", envir = .shud)
)
\end{verbatim}
\end{Usage}
%
\begin{Arguments}
\begin{ldescription}
\item[\code{varname}] vector of output keywords

\item[\code{rdsfile}] Save RDS file. NULL=do not save rds file.
\end{ldescription}
\end{Arguments}
%
\begin{Value}
A list of TimeSeries data.
\end{Value}
\inputencoding{utf8}
\HeaderA{map2d}{Raster plot the mesh data \code{map2d}}{map2d}
%
\begin{Description}\relax
Raster plot the mesh data
\code{map2d}
\end{Description}
%
\begin{Usage}
\begin{verbatim}
map2d(x = getElevation(), sp.riv = NULL)
\end{verbatim}
\end{Usage}
%
\begin{Arguments}
\begin{ldescription}
\item[\code{x}] Vector or Matrix of mesh data. Length of nrow must be equal the number of cells

\item[\code{sp.riv}] The SpatialLines of river network. sp.riv = NULL means plot raster only.
\end{ldescription}
\end{Arguments}
%
\begin{Value}
Raster of the x (vector) or the last row of x (matrix)
\end{Value}
\inputencoding{utf8}
\HeaderA{MeltFactor}{Generate the default melt factor \code{MeltFactor}}{MeltFactor}
%
\begin{Description}\relax
Generate the default melt factor
\code{MeltFactor}
\end{Description}
%
\begin{Usage}
\begin{verbatim}
MeltFactor(
  years,
  mf = c(0.001308019, 0.001633298, 0.002131198, 0.002632776, 0.003031171, 0.003197325,
    0.003095839, 0.00274524, 0.002260213, 0.001759481, 0.001373646, 0.001202083)
)
\end{verbatim}
\end{Usage}
%
\begin{Arguments}
\begin{ldescription}
\item[\code{years}] Years.

\item[\code{mf}] Default monthly melting factors.
\end{ldescription}
\end{Arguments}
%
\begin{Value}
Multiple years monthly melting factor.
\end{Value}
%
\begin{Examples}
\begin{ExampleCode}
years = 2000:2001
x = MeltFactor(years=years)
plot(x)
\end{ExampleCode}
\end{Examples}
\inputencoding{utf8}
\HeaderA{mesh\_to\_raster}{Convert mesh data to raster}{mesh.Rul.to.Rul.raster}
%
\begin{Description}\relax
Converts mesh (unstructured grid) data to a regular raster grid using
spatial interpolation. This function uses terra and sf for all spatial
operations and supports multiple interpolation methods.
\end{Description}
%
\begin{Usage}
\begin{verbatim}
mesh_to_raster(
  data,
  mesh = NULL,
  template = NULL,
  method = c("idw", "linear", "nearest"),
  resolution = NULL,
  crs = NULL,
  stack = FALSE
)
\end{verbatim}
\end{Usage}
%
\begin{Arguments}
\begin{ldescription}
\item[\code{data}] Numeric vector or matrix of values to rasterize.
If matrix, each row represents a time step and each column a mesh cell.
Length/ncol must equal the number of mesh cells.

\item[\code{mesh}] Mesh object (optional). If NULL, uses default mesh from
\code{readmesh()}.

\item[\code{template}] SpatRaster template defining output grid (optional).
If NULL, creates template from mesh extent.

\item[\code{method}] Character, interpolation method. One of:
\begin{itemize}

\item{} "idw" - Inverse Distance Weighting (default)
\item{} "linear" - Linear interpolation
\item{} "nearest" - Nearest neighbor

\end{itemize}


\item[\code{resolution}] Numeric, output raster resolution in map units (optional).
Only used if template is NULL.

\item[\code{crs}] Character or CRS object, coordinate reference system (optional).

\item[\code{stack}] Logical, if TRUE and data is a matrix, returns a multi-layer
SpatRaster with one layer per row (default: FALSE).
\end{ldescription}
\end{Arguments}
%
\begin{Value}
SpatRaster object. If stack=TRUE and data is a matrix, returns
a multi-layer raster.
\end{Value}
%
\begin{Section}{Migration Note}

This function replaces \code{MeshData2Raster()} and uses terra instead of
raster. Key differences:
\begin{itemize}

\item{} Returns SpatRaster instead of RasterLayer
\item{} Does not accept raster/sp objects - use terra/sf only
\item{} Method names simplified: "ide" -> "idw"
\item{} Removed plot parameter - use terra::plot() on result

\end{itemize}

\end{Section}
%
\begin{Examples}
\begin{ExampleCode}
## Not run: 
# Basic usage with vector data
mesh <- readmesh()
elevation <- getElevation()
r <- mesh_to_raster(elevation, mesh = mesh)
terra::plot(r)

# With custom resolution
r <- mesh_to_raster(elevation, mesh = mesh, resolution = 100)

# Time series data (matrix)
timeseries_data <- matrix(rnorm(1000 * 10), nrow = 10, ncol = 1000)
r_stack <- mesh_to_raster(timeseries_data, mesh = mesh, stack = TRUE)

## End(Not run)
\end{ExampleCode}
\end{Examples}
\inputencoding{utf8}
\HeaderA{MeshAtt}{Summary of the attributions on Mesh \code{MeshAtt}}{MeshAtt}
%
\begin{Description}\relax
Summary of the attributions on Mesh
\code{MeshAtt}
\end{Description}
%
\begin{Usage}
\begin{verbatim}
MeshAtt(
  pm = readmesh(),
  att = readatt(),
  soil = readsoil(),
  geol = readgeol(),
  lc = readlc()
)
\end{verbatim}
\end{Usage}
%
\begin{Arguments}
\begin{ldescription}
\item[\code{pm}] SHUD.MESH, .sp.mesh

\item[\code{att}] .sp.att of input

\item[\code{soil}] .para.soil of input

\item[\code{geol}] .para.geol of input

\item[\code{lc}] .para.lc of input
\end{ldescription}
\end{Arguments}
%
\begin{Value}
Attributes of each element.
\end{Value}
\inputencoding{utf8}
\HeaderA{meshSinks}{Get the Sink cells in SHUD mesh \code{meshSinks}}{meshSinks}
%
\begin{Description}\relax
Get the Sink cells in SHUD mesh
\code{meshSinks}
\end{Description}
%
\begin{Usage}
\begin{verbatim}
meshSinks(pm = readmesh(), details = FALSE)
\end{verbatim}
\end{Usage}
%
\begin{Arguments}
\begin{ldescription}
\item[\code{pm}] SHUD mesh

\item[\code{details}] Whether return details of sinks infomration.
\end{ldescription}
\end{Arguments}
%
\begin{Value}
Vector of TRUE/FALSE or data.frame
\end{Value}
\inputencoding{utf8}
\HeaderA{ModelInfo}{Summary of the basic information of model input data. \code{ModelInfo}}{ModelInfo}
%
\begin{Description}\relax
Summary of the basic information of model input data.
\code{ModelInfo}
\end{Description}
%
\begin{Usage}
\begin{verbatim}
ModelInfo(
  path = shud.filein()["outpath"],
  spr = readriv.sp(),
  crs.pcs = raster::crs(spr)
)
\end{verbatim}
\end{Usage}
%
\begin{Arguments}
\begin{ldescription}
\item[\code{path}] Where to export the basic model infomation.

\item[\code{spr}] river shapefile

\item[\code{crs.pcs}] The crs4string that defines the (x,y) in the model.
\end{ldescription}
\end{Arguments}
%
\begin{Value}
Basic model infomation, figures and tables
\end{Value}
\inputencoding{utf8}
\HeaderA{nc\_to\_raster}{Convert NetCDF data to raster/terra format}{nc.Rul.to.Rul.raster}
\keyword{netcdf-io}{nc\_to\_raster}
%
\begin{Description}\relax
Converts NetCDF data arrays to raster (legacy) or terra (modern) format.
Handles both single layers and multi-layer datasets.
\end{Description}
%
\begin{Usage}
\begin{verbatim}
nc_to_raster(
  nc_data = NULL,
  x = NULL,
  y = NULL,
  data_array = NULL,
  resolution = NULL,
  flip = TRUE,
  format = c("modern", "legacy")
)
\end{verbatim}
\end{Usage}
%
\begin{Arguments}
\begin{ldescription}
\item[\code{nc\_data}] List returned from read\_nc\_data()

\item[\code{x}] Numeric vector of x coordinates (alternative input)

\item[\code{y}] Numeric vector of y coordinates (alternative input)

\item[\code{data\_array}] Array of data values (alternative input)

\item[\code{resolution}] Numeric vector c(dx, dy) for grid resolution

\item[\code{flip}] Logical. Whether to flip the array vertically

\item[\code{format}] Character. Output format ("modern" or "legacy")
\end{ldescription}
\end{Arguments}
%
\begin{Value}
SpatRaster (modern) or RasterStack/RasterLayer (legacy)
\end{Value}
%
\begin{SeeAlso}\relax
Other netcdf-io: 
\code{\LinkA{get\_nc\_info}{get.Rul.nc.Rul.info}()},
\code{\LinkA{open\_nc\_file}{open.Rul.nc.Rul.file}()},
\code{\LinkA{read\_nc\_data}{read.Rul.nc.Rul.data}()},
\code{\LinkA{read\_nc\_time}{read.Rul.nc.Rul.time}()}
\end{SeeAlso}
%
\begin{Examples}
\begin{ExampleCode}
## Not run: 
library(ncdf4)
nc <- nc_open("climate_data.nc")
nc_data <- read_nc_data(nc)

# Convert to terra format (modern)
raster_data <- nc_to_raster(nc_data, format = "modern")

# Convert to raster format (legacy)
raster_data <- nc_to_raster(nc_data, format = "legacy")

nc_close(nc)

## End(Not run)
\end{ExampleCode}
\end{Examples}
\inputencoding{utf8}
\HeaderA{netcdf-io}{NetCDF I/O Functions}{netcdf.Rdash.io}
%
\begin{Description}\relax
This module provides functions for reading NetCDF files with support for
both legacy raster and modern terra formats. Functions maintain ncdf4
interface compatibility while adding terra integration.
\end{Description}
\inputencoding{utf8}
\HeaderA{NLCD.colors}{Return NLCD colors. \code{NLCD.colors}}{NLCD.colors}
\keyword{Landuse,}{NLCD.colors}
\keyword{NLCD,}{NLCD.colors}
\keyword{colormap}{NLCD.colors}
%
\begin{Description}\relax
Return NLCD colors.
\code{NLCD.colors}
\end{Description}
%
\begin{Usage}
\begin{verbatim}
NLCD.colors(x, def.col = "#EEEEEE")
\end{verbatim}
\end{Usage}
%
\begin{Arguments}
\begin{ldescription}
\item[\code{x, }] which is the integer number defined in NLCD 2001,2006, and 2011.  Any value out of (0,100) will be assigned as black.

\item[\code{def.col}] Default color, if the landuse code is missing.
\end{ldescription}
\end{Arguments}
%
\begin{Value}
HEX color strings
\end{Value}
%
\begin{Source}\relax
See the link of NLCD: \url{http://www.mrlc.gov/nlcd11_leg.php}
\end{Source}
\inputencoding{utf8}
\HeaderA{NLCD.names}{Return NLCD class names. \code{NLCD.names}}{NLCD.names}
\keyword{Landuse}{NLCD.names}
\keyword{NLCD,}{NLCD.names}
%
\begin{Description}\relax
Return NLCD class names.
\code{NLCD.names}
\end{Description}
%
\begin{Usage}
\begin{verbatim}
NLCD.names(x, def.name = "n/a")
\end{verbatim}
\end{Usage}
%
\begin{Arguments}
\begin{ldescription}
\item[\code{x, }] which is the integer number defined in NLCD 2001,2006, and 2011.  Any value out of (0,100) will be assigned as black.

\item[\code{def.name}] Default name.
\end{ldescription}
\end{Arguments}
%
\begin{Value}
Characters
\end{Value}
\inputencoding{utf8}
\HeaderA{NodeIDList}{return the Index of nodes in SpatialData \code{NodeIDList}}{NodeIDList}
%
\begin{Description}\relax
return the Index of nodes in SpatialData
\code{NodeIDList}
\end{Description}
%
\begin{Usage}
\begin{verbatim}
NodeIDList(sp, coord = extractCoords(sp, unique = TRUE))
\end{verbatim}
\end{Usage}
%
\begin{Arguments}
\begin{ldescription}
\item[\code{sp}] SpatialLines*

\item[\code{coord}] Coordinate of vertex in \code{sp}
\end{ldescription}
\end{Arguments}
%
\begin{Value}
list of Index of each SpatialData, line or polygon
\end{Value}
\inputencoding{utf8}
\HeaderA{open\_nc\_file}{Open NetCDF file with error handling}{open.Rul.nc.Rul.file}
\keyword{netcdf-io}{open\_nc\_file}
%
\begin{Description}\relax
Wrapper around ncdf4::nc\_open with improved error handling.
\end{Description}
%
\begin{Usage}
\begin{verbatim}
open_nc_file(file)
\end{verbatim}
\end{Usage}
%
\begin{Arguments}
\begin{ldescription}
\item[\code{file}] Character. Path to NetCDF file
\end{ldescription}
\end{Arguments}
%
\begin{Value}
NetCDF connection object
\end{Value}
%
\begin{SeeAlso}\relax
Other netcdf-io: 
\code{\LinkA{get\_nc\_info}{get.Rul.nc.Rul.info}()},
\code{\LinkA{nc\_to\_raster}{nc.Rul.to.Rul.raster}()},
\code{\LinkA{read\_nc\_data}{read.Rul.nc.Rul.data}()},
\code{\LinkA{read\_nc\_time}{read.Rul.nc.Rul.time}()}
\end{SeeAlso}
\inputencoding{utf8}
\HeaderA{output-io}{Model Output I/O Functions}{output.Rdash.io}
%
\begin{Description}\relax
This module provides functions for reading SHUD model output files.
Functions support both legacy and modern formats with terra integration.
\end{Description}
\inputencoding{utf8}
\HeaderA{PET\_PM}{Potential Evapotranspiration with Pennmann-Monteith equation. \code{PET\_PM}}{PET.Rul.PM}
%
\begin{Description}\relax
Potential Evapotranspiration with Pennmann-Monteith equation.
\code{PET\_PM}
\end{Description}
%
\begin{Usage}
\begin{verbatim}
PET_PM(
  Wind,
  Temp,
  RH,
  RadNet,
  Press,
  WindHeight = 10,
  Veg_height = 0.12,
  albedo = 0.23,
  Res_surf = 70,
  Elevation = NULL
)
\end{verbatim}
\end{Usage}
%
\begin{Arguments}
\begin{ldescription}
\item[\code{Wind}] Windspeed \LinkA{m s-1}{m s.Rdash.1}

\item[\code{Temp}] Air temperature, \LinkA{C}{C}

\item[\code{RH}] Relative Humidity \LinkA{-}{.Rdash.} 0\textasciitilde{}1.

\item[\code{RadNet}] Net Solar Radiation \LinkA{w m-2}{w m.Rdash.2}

\item[\code{Press}] Pressure \LinkA{kPa}{kPa}

\item[\code{WindHeight}] Height that wind was measured \LinkA{m}{m}

\item[\code{Veg\_height}] Height of vegetation \LinkA{m}{m}

\item[\code{albedo}] Albedo \LinkA{-}{.Rdash.}

\item[\code{Res\_surf}] Aerodynamic surface resistance \LinkA{s m-1}{s m.Rdash.1}

\item[\code{Elevation}] Elevation \LinkA{m}{m}
\end{ldescription}
\end{Arguments}
%
\begin{Value}
PET at \LinkA{mm m-2 s-1}{mm m.Rdash.2 s.Rdash.1} or \LinkA{kg m-2 s-1}{kg m.Rdash.2 s.Rdash.1}
\end{Value}
\inputencoding{utf8}
\HeaderA{plot\_animate}{Plot animation maps \code{plot.animate}}{plot.Rul.animate}
%
\begin{Description}\relax
Plot animation maps
\code{plot.animate}
\end{Description}
%
\begin{Usage}
\begin{verbatim}
plot_animate(x, id = NULL, nmap = 10, rlist = NULL)
\end{verbatim}
\end{Usage}
%
\begin{Arguments}
\begin{ldescription}
\item[\code{x}] Time-series data.

\item[\code{id}] Which row(s) to plot in the animation.

\item[\code{nmap}] Number of maps to be plot.

\item[\code{rlist}] RasterStack of the maps.
\end{ldescription}
\end{Arguments}
\inputencoding{utf8}
\HeaderA{plot\_sp}{Plot SpatialPolygons maps \code{plot.sp}}{plot.Rul.sp}
%
\begin{Description}\relax
Plot SpatialPolygons maps
\code{plot.sp}
\end{Description}
%
\begin{Usage}
\begin{verbatim}
plot_sp(x, zcol, col.fun = topo.colors, ncol = 20, ...)
\end{verbatim}
\end{Usage}
%
\begin{Arguments}
\begin{ldescription}
\item[\code{x}] Spatial*

\item[\code{zcol}] The column to plot

\item[\code{col.fun}] Functions of coloring.

\item[\code{ncol}] Number of colors in color function.

\item[\code{...}] more options in raster::plot()
\end{ldescription}
\end{Arguments}
\inputencoding{utf8}
\HeaderA{plot\_tsd}{Plot xts data.. \code{plot\_tsd}}{plot.Rul.tsd}
%
\begin{Description}\relax
Plot xts data..
\code{plot\_tsd}
\end{Description}
%
\begin{Usage}
\begin{verbatim}
plot_tsd(x, time.col = "year")
\end{verbatim}
\end{Usage}
%
\begin{Arguments}
\begin{ldescription}
\item[\code{x}] xts data

\item[\code{time.col}] Time interval of banded background. Default is 'year'
\end{ldescription}
\end{Arguments}
%
\begin{Examples}
\begin{ExampleCode}
library(xts)
nday = 1000
xd=as.POSIXct(as.Date('2000-01-01')+ 1:nday )
x=as.xts(sin(1:1000 / 100), order.by=xd)
plot_tsd(x)
\end{ExampleCode}
\end{Examples}
\inputencoding{utf8}
\HeaderA{png.control}{Prepare png figure out. \code{png.control}}{png.control}
%
\begin{Description}\relax
Prepare png figure out.
\code{png.control}
\end{Description}
%
\begin{Usage}
\begin{verbatim}
png.control(
  fn = "figure.png",
  file = file.path(path, fn),
  path = get("anapath", envir = .shud),
  ratio = 9/11,
  ht = 11,
  wd = ht * ratio,
  units = "in",
  res = 200
)
\end{verbatim}
\end{Usage}
%
\begin{Arguments}
\begin{ldescription}
\item[\code{fn}] Filename

\item[\code{file}] Full path of file

\item[\code{path}] The path only. Default = anapath in .shud environment.

\item[\code{ratio}] Ratio of width to height.

\item[\code{ht}] Height.

\item[\code{wd}] Width.

\item[\code{units}] Units of width and height.

\item[\code{res}] Resolution along height and width.
\end{ldescription}
\end{Arguments}
\inputencoding{utf8}
\HeaderA{PointInDistance}{Find the points in distance less than tol. \code{PointInDistance}}{PointInDistance}
%
\begin{Description}\relax
Find the points in distance less than tol.
\code{PointInDistance}
\end{Description}
%
\begin{Usage}
\begin{verbatim}
PointInDistance(pt, tol)
\end{verbatim}
\end{Usage}
%
\begin{Arguments}
\begin{ldescription}
\item[\code{pt}] 2-column coordinates (x,y).

\item[\code{tol}] Tolerance
\end{ldescription}
\end{Arguments}
%
\begin{Value}
Index of points that within tol
\end{Value}
\inputencoding{utf8}
\HeaderA{polygonArea}{Calculate the Polygon Area;}{polygonArea}
%
\begin{Description}\relax
Calculate the Polygon Area;
\end{Description}
%
\begin{Usage}
\begin{verbatim}
polygonArea(X, Y)
\end{verbatim}
\end{Usage}
%
\begin{Arguments}
\begin{ldescription}
\item[\code{X}] required

\item[\code{Y}] required
\end{ldescription}
\end{Arguments}
%
\begin{Value}
Area
\end{Value}
\inputencoding{utf8}
\HeaderA{PressureElevation}{Air pressure at the elevation, Allen(1998) Eq(7) \code{PressureElevation}}{PressureElevation}
%
\begin{Description}\relax
Air pressure at the elevation, Allen(1998) Eq(7)
\code{PressureElevation}
\end{Description}
%
\begin{Usage}
\begin{verbatim}
PressureElevation(z)
\end{verbatim}
\end{Usage}
%
\begin{Arguments}
\begin{ldescription}
\item[\code{z}] Elevation \LinkA{m}{m}
\end{ldescription}
\end{Arguments}
%
\begin{Value}
Air pressure at the elevation \LinkA{kPa}{kPa}
\end{Value}
\inputencoding{utf8}
\HeaderA{project\_coords}{Project coordinates to a new CRS}{project.Rul.coords}
%
\begin{Description}\relax
Unified interface for projecting spatial coordinates or objects to a new
coordinate reference system. This function provides a consistent interface
for coordinate transformation across different spatial object types.
\end{Description}
%
\begin{Usage}
\begin{verbatim}
project_coords(x, from_crs = NULL, to_crs, coords = c("x", "y"))
\end{verbatim}
\end{Usage}
%
\begin{Arguments}
\begin{ldescription}
\item[\code{x}] Spatial object (sf, SpatVector, SpatRaster) or numeric matrix/
data.frame with coordinates (columns: x, y or lon, lat)

\item[\code{from\_crs}] Source CRS. If NULL, extracted from x (for spatial objects)
or assumed to be EPSG:4326 (for coordinate matrices)

\item[\code{to\_crs}] Target CRS (character string, EPSG code, or PROJ.4 string)

\item[\code{coords}] Character vector of length 2 specifying coordinate column
names for matrix/data.frame input. Default: c("x", "y")
\end{ldescription}
\end{Arguments}
%
\begin{Value}
Transformed spatial object or coordinate matrix/data.frame
\end{Value}
%
\begin{Examples}
\begin{ExampleCode}
# Project sf object
if (requireNamespace("sf", quietly = TRUE)) {
  pts <- sf::st_as_sf(
    data.frame(x = c(-75, -76), y = c(40, 41)),
    coords = c("x", "y"),
    crs = 4326
  )
  pts_utm <- project_coords(pts, to_crs = crs.long2utm(-75, 40))
}

# Project coordinate matrix
coords <- matrix(c(-75, 40, -76, 41), ncol = 2, byrow = TRUE)
coords_utm <- project_coords(coords, from_crs = "EPSG:4326",
                              to_crs = crs.long2utm(-75, 40))
\end{ExampleCode}
\end{Examples}
\inputencoding{utf8}
\HeaderA{ProjectCoordinate}{Re-project coordinates betwen GCS and PCS \code{ProjectCoordinate}}{ProjectCoordinate}
%
\begin{Description}\relax
Re-project coordinates betwen GCS and PCS
\code{ProjectCoordinate}
\end{Description}
%
\begin{Usage}
\begin{verbatim}
ProjectCoordinate(x, proj4string, P2G = TRUE)
\end{verbatim}
\end{Usage}
%
\begin{Arguments}
\begin{ldescription}
\item[\code{x}] 2-column matrix of coordinates.

\item[\code{proj4string}] proj4string

\item[\code{P2G}] if TRUE, a cartographic projection into lat/long, otherwise projects from lat/long into a cartographic projection.
\end{ldescription}
\end{Arguments}
%
\begin{Value}
Basic model infomation, figures and tables
\end{Value}
\inputencoding{utf8}
\HeaderA{projection}{Coordinate Projection Functions}{projection}
%
\begin{Description}\relax
Functions for creating and working with coordinate reference systems (CRS)
using modern spatial libraries (sf).
\end{Description}
\inputencoding{utf8}
\HeaderA{PsychrometricConstant}{Psychrometric Constant, Eq 4.2.1 in David R Maidment, Handbook of Hydrology \code{PsychrometricConstant}}{PsychrometricConstant}
%
\begin{Description}\relax
Psychrometric Constant, Eq 4.2.1 in David R Maidment, Handbook of Hydrology
\code{PsychrometricConstant}
\end{Description}
%
\begin{Usage}
\begin{verbatim}
PsychrometricConstant(Pressure, lambda)
\end{verbatim}
\end{Usage}
%
\begin{Arguments}
\begin{ldescription}
\item[\code{Pressure}] Air pressure \LinkA{kPa}{kPa}

\item[\code{lambda}] Latent heat of vaporized water. \LinkA{MJ/kg}{MJ/kg}
\end{ldescription}
\end{Arguments}
%
\begin{Value}
γ psychrometric constant \LinkA{kPa °C-1}{kPa °C.Rdash.1}
\end{Value}
\inputencoding{utf8}
\HeaderA{PTF}{Pedotransfer function to generate soil/geol parameter from soil texture \code{PTF}}{PTF}
%
\begin{Description}\relax
Pedotransfer function to generate soil/geol parameter from soil texture
\code{PTF}
\end{Description}
%
\begin{Usage}
\begin{verbatim}
PTF(
  x = t(matrix(c(33, 33, 2, 1.4), ncol = 5, nrow = 4)),
  topsoil = TRUE,
  xmin = c(0.1, 0.1, 0.8, 1.2)
)
\end{verbatim}
\end{Usage}
%
\begin{Arguments}
\begin{ldescription}
\item[\code{x}] data.frame or matrix, column = c(Silt\_perc, Clay\_perc, OrganicMatter \_perc, BulkDensity g/cm3)

\item[\code{topsoil}] TRUE/FALSE, default = TRUE

\item[\code{xmin}] Lower limitation of the value of c(Silt\_perc, Clay\_perc, OrganicMatter \_perc, BulkDensity g/cm3)
\end{ldescription}
\end{Arguments}
%
\begin{Value}
Hydraulic parameters, matrix
\end{Value}
%
\begin{Examples}
\begin{ExampleCode}
x = cbind(10, 40, 1:20, 1.5)
y = PTF(x)
apply(y[, -1], 2, summary)
plot(x[, 3], y[, 2])
\end{ExampleCode}
\end{Examples}
\inputencoding{utf8}
\HeaderA{PTF.geol}{Pedotransfer function to generate soil parameter \code{PTF.geol}}{PTF.geol}
%
\begin{Description}\relax
Pedotransfer function to generate soil parameter
\code{PTF.geol}
\end{Description}
%
\begin{Usage}
\begin{verbatim}
PTF.geol(
  x = t(matrix(c(33, 33, 2, 1.4), ncol = 5, nrow = 4)),
  topsoil = FALSE,
  rm.outlier = FALSE
)
\end{verbatim}
\end{Usage}
%
\begin{Arguments}
\begin{ldescription}
\item[\code{x}] data.frame or matrix, column = c(Silt\_perc, Clay\_perc, OrganicMatter \_perc, BulkDensity g/cm3)

\item[\code{topsoil}] TRUE/FALSE, default = FALSE

\item[\code{rm.outlier}] Whether replace the outliers in each column with the mean value of the column.
\end{ldescription}
\end{Arguments}
%
\begin{Value}
Hydraulic parameters, matrix
\end{Value}
%
\begin{Examples}
\begin{ExampleCode}
x = cbind(10,40, 1:20, 1.5)
y = PTF.geol(x)
apply(y[,-1], 2, summary)
plot(x[,3], y[,2])
\end{ExampleCode}
\end{Examples}
\inputencoding{utf8}
\HeaderA{PTF.soil}{Pedotransfer function to generate soil parameter \code{PTF.soil}}{PTF.soil}
%
\begin{Description}\relax
Pedotransfer function to generate soil parameter
\code{PTF.soil}
\end{Description}
%
\begin{Usage}
\begin{verbatim}
PTF.soil(
  x = t(matrix(c(33, 33, 2, 1.4), ncol = 5, nrow = 4)),
  topsoil = TRUE,
  rm.outlier = FALSE
)
\end{verbatim}
\end{Usage}
%
\begin{Arguments}
\begin{ldescription}
\item[\code{x}] data.frame or matrix, column = c(Silt\_perc, Clay\_perc, OrganicMatter \_perc, BulkDensity g/cm3)

\item[\code{topsoil}] TRUE/FALSE, default = TRUE

\item[\code{rm.outlier}] Whether replace the outliers in each column with the mean value of the column.
\end{ldescription}
\end{Arguments}
%
\begin{Value}
Hydraulic parameters, matrix
\end{Value}
%
\begin{Examples}
\begin{ExampleCode}
x=cbind(10,40, 1:20, 1.5)
y=PTF.soil(x)
apply(y[,-1], 2, summary)
plot(x[,3], y[,2])
\end{ExampleCode}
\end{Examples}
\inputencoding{utf8}
\HeaderA{quick\_model}{Quick Model Building with Default Parameters}{quick.Rul.model}
%
\begin{Description}\relax
Simplified interface for building SHUD models with sensible defaults.
This function is designed for quick prototyping and common use cases.
\end{Description}
%
\begin{Usage}
\begin{verbatim}
quick_model(
  project_name,
  domain,
  dem,
  rivers = NULL,
  landcover = NULL,
  output_dir = getwd(),
  years = 2000:2010,
  verbose = TRUE
)
\end{verbatim}
\end{Usage}
%
\begin{Arguments}
\begin{ldescription}
\item[\code{project\_name}] Character, project name for output files

\item[\code{domain}] sf object with POLYGON geometry defining watershed boundary

\item[\code{dem}] SpatRaster object with elevation data

\item[\code{rivers}] sf object with LINESTRING geometry for river network (optional)

\item[\code{landcover}] SpatRaster object with land cover data (optional)

\item[\code{output\_dir}] Character, output directory path (default: current directory)

\item[\code{years}] Integer vector, years for time series (default: 2000:2010)

\item[\code{verbose}] Logical, display progress messages (default: TRUE)
\end{ldescription}
\end{Arguments}
%
\begin{Details}\relax
This function uses default parameters optimized for typical watersheds:
\begin{itemize}

\item{} Maximum triangle area: 200,000 m² (20 ha)
\item{} Minimum triangle angle: 30 degrees
\item{} Boundary simplification: 200 m
\item{} Aquifer depth: 10 m

\end{itemize}


For more control over parameters, use \code{\LinkA{auto\_build\_model}{auto.Rul.build.Rul.model}}.
\end{Details}
%
\begin{Value}
List with components:
\begin{ldescription}
\item[\code{mesh}] SHUD.MESH object
\item[\code{river}] SHUD.RIVER object (if rivers provided)
\item[\code{files}] Named vector of output file paths
\end{ldescription}
\end{Value}
%
\begin{Section}{Migration Note}

This is a new function with no legacy equivalent. It provides a streamlined
interface for users who want to quickly build models without specifying
detailed options.
\end{Section}
%
\begin{Examples}
\begin{ExampleCode}
## Not run: 
library(sf)
library(terra)

# Load minimal required data
domain <- st_read("watershed.shp")
dem <- rast("dem.tif")

# Quick model build
model <- quick_model(
  project_name = "test_model",
  domain = domain,
  dem = dem
)

# With rivers and land cover
rivers <- st_read("rivers.shp")
lc <- rast("landcover.tif")

model <- quick_model(
  project_name = "full_model",
  domain = domain,
  dem = dem,
  rivers = rivers,
  landcover = lc,
  output_dir = "quick_model"
)

## End(Not run)
\end{ExampleCode}
\end{Examples}
\inputencoding{utf8}
\HeaderA{QvsO}{Default Comparison of the observation and simulation \code{QvsO}}{QvsO}
%
\begin{Description}\relax
Default Comparison of the observation and simulation
\code{QvsO}
\end{Description}
%
\begin{Usage}
\begin{verbatim}
QvsO(qq, sim = qq[, 2], obs = qq[, 1], ...)
\end{verbatim}
\end{Usage}
%
\begin{Arguments}
\begin{ldescription}
\item[\code{qq}] Matrix or data.frame. Column 1 is observation, while colum 2 is simulation.

\item[\code{sim}] Simulation time-seris data;

\item[\code{obs}] Observation time-seris data;

\item[\code{...}] more options for plot()
\end{ldescription}
\end{Arguments}
%
\begin{Examples}
\begin{ExampleCode}
obs = rnorm(100)
sim = obs + rnorm(100) / 2
QvsO(cbind(obs, sim))
\end{ExampleCode}
\end{Examples}
\inputencoding{utf8}
\HeaderA{rcpp\_hello\_world}{Simple function using Rcpp}{rcpp.Rul.hello.Rul.world}
%
\begin{Description}\relax
Simple function using Rcpp
\end{Description}
%
\begin{Usage}
\begin{verbatim}
rcpp_hello_world()	
\end{verbatim}
\end{Usage}
%
\begin{Examples}
\begin{ExampleCode}
## Not run: 
rcpp_hello_world()

## End(Not run)
\end{ExampleCode}
\end{Examples}
\inputencoding{utf8}
\HeaderA{read.df}{Read the matrix or data.frame file which has (nrow, ncol) in header \code{read.df}}{read.df}
%
\begin{Description}\relax
Read the matrix or data.frame file which has (nrow, ncol) in header
\code{read.df}
\end{Description}
%
\begin{Usage}
\begin{verbatim}
read.df(file, text = readLines(file), sep = "\t")
\end{verbatim}
\end{Usage}
%
\begin{Arguments}
\begin{ldescription}
\item[\code{file}] full path of file

\item[\code{text}] text return from readLines(file)

\item[\code{sep}] seperator of the table.
\end{ldescription}
\end{Arguments}
%
\begin{Value}
a list of matrix
\end{Value}
\inputencoding{utf8}
\HeaderA{read.tsd}{Read the time-series file, which has c(nrow, ncol, timestring) in header \code{read.tsd}}{read.tsd}
%
\begin{Description}\relax
Read the time-series file, which has c(nrow, ncol, timestring) in header
\code{read.tsd}
\end{Description}
%
\begin{Usage}
\begin{verbatim}
read.tsd(file)
\end{verbatim}
\end{Usage}
%
\begin{Arguments}
\begin{ldescription}
\item[\code{file}] full path of file
\end{ldescription}
\end{Arguments}
%
\begin{Value}
xts
\end{Value}
\inputencoding{utf8}
\HeaderA{read\_att}{Read SHUD attributes file}{read.Rul.att}
\keyword{shud-io}{read\_att}
%
\begin{Description}\relax
Reads a SHUD attributes file (.att) and returns attribute data.
\end{Description}
%
\begin{Usage}
\begin{verbatim}
read_att(file = shud.filein()["md.att"])
\end{verbatim}
\end{Usage}
%
\begin{Arguments}
\begin{ldescription}
\item[\code{file}] Character. Full path to the .att file
\end{ldescription}
\end{Arguments}
%
\begin{Value}
Matrix containing attribute data
\end{Value}
%
\begin{SeeAlso}\relax
Other shud-io: 
\code{\LinkA{read\_calib}{read.Rul.calib}()},
\code{\LinkA{read\_config}{read.Rul.config}()},
\code{\LinkA{read\_df}{read.Rul.df}()},
\code{\LinkA{read\_forc\_csv}{read.Rul.forc.Rul.csv}()},
\code{\LinkA{read\_forc\_fn}{read.Rul.forc.Rul.fn}()},
\code{\LinkA{read\_geol}{read.Rul.geol}()},
\code{\LinkA{read\_ic}{read.Rul.ic}()},
\code{\LinkA{read\_lc}{read.Rul.lc}()},
\code{\LinkA{read\_mesh}{read.Rul.mesh}()},
\code{\LinkA{read\_para}{read.Rul.para}()},
\code{\LinkA{read\_river\_sp}{read.Rul.river.Rul.sp}()},
\code{\LinkA{read\_river}{read.Rul.river}()},
\code{\LinkA{read\_rivseg}{read.Rul.rivseg}()},
\code{\LinkA{read\_soil}{read.Rul.soil}()},
\code{\LinkA{write\_config}{write.Rul.config}()},
\code{\LinkA{write\_df}{write.Rul.df}()},
\code{\LinkA{write\_forc}{write.Rul.forc}()},
\code{\LinkA{write\_ic}{write.Rul.ic}()},
\code{\LinkA{write\_mesh}{write.Rul.mesh}()},
\code{\LinkA{write\_river}{write.Rul.river}()}
\end{SeeAlso}
%
\begin{Examples}
\begin{ExampleCode}
## Not run: 
att <- read_att("model.att")

## End(Not run)
\end{ExampleCode}
\end{Examples}
\inputencoding{utf8}
\HeaderA{read\_calib}{Read SHUD calibration file}{read.Rul.calib}
\keyword{shud-io}{read\_calib}
%
\begin{Description}\relax
Reads a SHUD calibration file (.calib) and returns calibration data.
\end{Description}
%
\begin{Usage}
\begin{verbatim}
read_calib(file = shud.filein()["md.calib"])
\end{verbatim}
\end{Usage}
%
\begin{Arguments}
\begin{ldescription}
\item[\code{file}] Character. Full path to the .calib file
\end{ldescription}
\end{Arguments}
%
\begin{Value}
Data frame containing calibration data
\end{Value}
%
\begin{SeeAlso}\relax
Other shud-io: 
\code{\LinkA{read\_att}{read.Rul.att}()},
\code{\LinkA{read\_config}{read.Rul.config}()},
\code{\LinkA{read\_df}{read.Rul.df}()},
\code{\LinkA{read\_forc\_csv}{read.Rul.forc.Rul.csv}()},
\code{\LinkA{read\_forc\_fn}{read.Rul.forc.Rul.fn}()},
\code{\LinkA{read\_geol}{read.Rul.geol}()},
\code{\LinkA{read\_ic}{read.Rul.ic}()},
\code{\LinkA{read\_lc}{read.Rul.lc}()},
\code{\LinkA{read\_mesh}{read.Rul.mesh}()},
\code{\LinkA{read\_para}{read.Rul.para}()},
\code{\LinkA{read\_river\_sp}{read.Rul.river.Rul.sp}()},
\code{\LinkA{read\_river}{read.Rul.river}()},
\code{\LinkA{read\_rivseg}{read.Rul.rivseg}()},
\code{\LinkA{read\_soil}{read.Rul.soil}()},
\code{\LinkA{write\_config}{write.Rul.config}()},
\code{\LinkA{write\_df}{write.Rul.df}()},
\code{\LinkA{write\_forc}{write.Rul.forc}()},
\code{\LinkA{write\_ic}{write.Rul.ic}()},
\code{\LinkA{write\_mesh}{write.Rul.mesh}()},
\code{\LinkA{write\_river}{write.Rul.river}()}
\end{SeeAlso}
%
\begin{Examples}
\begin{ExampleCode}
## Not run: 
calib <- read_calib("model.calib")

## End(Not run)
\end{ExampleCode}
\end{Examples}
\inputencoding{utf8}
\HeaderA{read\_config}{Read SHUD configuration file}{read.Rul.config}
\keyword{shud-io}{read\_config}
%
\begin{Description}\relax
Reads a SHUD configuration file (.para or .calib) and returns configuration data.
\end{Description}
%
\begin{Usage}
\begin{verbatim}
read_config(file = shud.filein()["md.para"])
\end{verbatim}
\end{Usage}
%
\begin{Arguments}
\begin{ldescription}
\item[\code{file}] Character. Full path to the configuration file
\end{ldescription}
\end{Arguments}
%
\begin{Value}
Data frame containing configuration data
\end{Value}
%
\begin{SeeAlso}\relax
Other shud-io: 
\code{\LinkA{read\_att}{read.Rul.att}()},
\code{\LinkA{read\_calib}{read.Rul.calib}()},
\code{\LinkA{read\_df}{read.Rul.df}()},
\code{\LinkA{read\_forc\_csv}{read.Rul.forc.Rul.csv}()},
\code{\LinkA{read\_forc\_fn}{read.Rul.forc.Rul.fn}()},
\code{\LinkA{read\_geol}{read.Rul.geol}()},
\code{\LinkA{read\_ic}{read.Rul.ic}()},
\code{\LinkA{read\_lc}{read.Rul.lc}()},
\code{\LinkA{read\_mesh}{read.Rul.mesh}()},
\code{\LinkA{read\_para}{read.Rul.para}()},
\code{\LinkA{read\_river\_sp}{read.Rul.river.Rul.sp}()},
\code{\LinkA{read\_river}{read.Rul.river}()},
\code{\LinkA{read\_rivseg}{read.Rul.rivseg}()},
\code{\LinkA{read\_soil}{read.Rul.soil}()},
\code{\LinkA{write\_config}{write.Rul.config}()},
\code{\LinkA{write\_df}{write.Rul.df}()},
\code{\LinkA{write\_forc}{write.Rul.forc}()},
\code{\LinkA{write\_ic}{write.Rul.ic}()},
\code{\LinkA{write\_mesh}{write.Rul.mesh}()},
\code{\LinkA{write\_river}{write.Rul.river}()}
\end{SeeAlso}
%
\begin{Examples}
\begin{ExampleCode}
## Not run: 
config <- read_config("model.para")

## End(Not run)
\end{ExampleCode}
\end{Examples}
\inputencoding{utf8}
\HeaderA{read\_df}{Read matrix or data frame file with header}{read.Rul.df}
\keyword{shud-io}{read\_df}
%
\begin{Description}\relax
Reads a file containing matrix or data frame data with (nrow, ncol) in header.
This is a core utility function used by other SHUD file readers.
\end{Description}
%
\begin{Usage}
\begin{verbatim}
read_df(file, text = readLines(file), sep = "\t")
\end{verbatim}
\end{Usage}
%
\begin{Arguments}
\begin{ldescription}
\item[\code{file}] Character. Full path to file

\item[\code{text}] Character vector. Text content from readLines(file). If provided, file is ignored

\item[\code{sep}] Character. Separator for the table (default: tab)
\end{ldescription}
\end{Arguments}
%
\begin{Value}
List of matrices/data frames
\end{Value}
%
\begin{SeeAlso}\relax
Other shud-io: 
\code{\LinkA{read\_att}{read.Rul.att}()},
\code{\LinkA{read\_calib}{read.Rul.calib}()},
\code{\LinkA{read\_config}{read.Rul.config}()},
\code{\LinkA{read\_forc\_csv}{read.Rul.forc.Rul.csv}()},
\code{\LinkA{read\_forc\_fn}{read.Rul.forc.Rul.fn}()},
\code{\LinkA{read\_geol}{read.Rul.geol}()},
\code{\LinkA{read\_ic}{read.Rul.ic}()},
\code{\LinkA{read\_lc}{read.Rul.lc}()},
\code{\LinkA{read\_mesh}{read.Rul.mesh}()},
\code{\LinkA{read\_para}{read.Rul.para}()},
\code{\LinkA{read\_river\_sp}{read.Rul.river.Rul.sp}()},
\code{\LinkA{read\_river}{read.Rul.river}()},
\code{\LinkA{read\_rivseg}{read.Rul.rivseg}()},
\code{\LinkA{read\_soil}{read.Rul.soil}()},
\code{\LinkA{write\_config}{write.Rul.config}()},
\code{\LinkA{write\_df}{write.Rul.df}()},
\code{\LinkA{write\_forc}{write.Rul.forc}()},
\code{\LinkA{write\_ic}{write.Rul.ic}()},
\code{\LinkA{write\_mesh}{write.Rul.mesh}()},
\code{\LinkA{write\_river}{write.Rul.river}()}
\end{SeeAlso}
%
\begin{Examples}
\begin{ExampleCode}
## Not run: 
data <- read_df("data.txt")

## End(Not run)
\end{ExampleCode}
\end{Examples}
\inputencoding{utf8}
\HeaderA{read\_forc\_csv}{Read SHUD forcing CSV files}{read.Rul.forc.Rul.csv}
\keyword{shud-io}{read\_forc\_csv}
%
\begin{Description}\relax
Reads the CSV files referenced in a SHUD forcing file (.forc).
\end{Description}
%
\begin{Usage}
\begin{verbatim}
read_forc_csv(file = shud.filein()["md.forc"], id = NULL)
\end{verbatim}
\end{Usage}
%
\begin{Arguments}
\begin{ldescription}
\item[\code{file}] Character. Full path to the .forc file

\item[\code{id}] Numeric vector. Index of forcing sites to read. If NULL, returns average of all sites
\end{ldescription}
\end{Arguments}
%
\begin{Value}
Forcing data as xts object or list of xts objects
\end{Value}
%
\begin{SeeAlso}\relax
Other shud-io: 
\code{\LinkA{read\_att}{read.Rul.att}()},
\code{\LinkA{read\_calib}{read.Rul.calib}()},
\code{\LinkA{read\_config}{read.Rul.config}()},
\code{\LinkA{read\_df}{read.Rul.df}()},
\code{\LinkA{read\_forc\_fn}{read.Rul.forc.Rul.fn}()},
\code{\LinkA{read\_geol}{read.Rul.geol}()},
\code{\LinkA{read\_ic}{read.Rul.ic}()},
\code{\LinkA{read\_lc}{read.Rul.lc}()},
\code{\LinkA{read\_mesh}{read.Rul.mesh}()},
\code{\LinkA{read\_para}{read.Rul.para}()},
\code{\LinkA{read\_river\_sp}{read.Rul.river.Rul.sp}()},
\code{\LinkA{read\_river}{read.Rul.river}()},
\code{\LinkA{read\_rivseg}{read.Rul.rivseg}()},
\code{\LinkA{read\_soil}{read.Rul.soil}()},
\code{\LinkA{write\_config}{write.Rul.config}()},
\code{\LinkA{write\_df}{write.Rul.df}()},
\code{\LinkA{write\_forc}{write.Rul.forc}()},
\code{\LinkA{write\_ic}{write.Rul.ic}()},
\code{\LinkA{write\_mesh}{write.Rul.mesh}()},
\code{\LinkA{write\_river}{write.Rul.river}()}
\end{SeeAlso}
%
\begin{Examples}
\begin{ExampleCode}
## Not run: 
forc_data <- read_forc_csv("model.forc", id = 1:3)

## End(Not run)
\end{ExampleCode}
\end{Examples}
\inputencoding{utf8}
\HeaderA{read\_forc\_fn}{Read SHUD forcing file list}{read.Rul.forc.Rul.fn}
\keyword{shud-io}{read\_forc\_fn}
%
\begin{Description}\relax
Reads a SHUD forcing file (.forc) and returns forcing site information.
\end{Description}
%
\begin{Usage}
\begin{verbatim}
read_forc_fn(file = shud.filein()["md.forc"])
\end{verbatim}
\end{Usage}
%
\begin{Arguments}
\begin{ldescription}
\item[\code{file}] Character. Full path to the .forc file
\end{ldescription}
\end{Arguments}
%
\begin{Value}
List containing start time and forcing sites data frame
\end{Value}
%
\begin{SeeAlso}\relax
Other shud-io: 
\code{\LinkA{read\_att}{read.Rul.att}()},
\code{\LinkA{read\_calib}{read.Rul.calib}()},
\code{\LinkA{read\_config}{read.Rul.config}()},
\code{\LinkA{read\_df}{read.Rul.df}()},
\code{\LinkA{read\_forc\_csv}{read.Rul.forc.Rul.csv}()},
\code{\LinkA{read\_geol}{read.Rul.geol}()},
\code{\LinkA{read\_ic}{read.Rul.ic}()},
\code{\LinkA{read\_lc}{read.Rul.lc}()},
\code{\LinkA{read\_mesh}{read.Rul.mesh}()},
\code{\LinkA{read\_para}{read.Rul.para}()},
\code{\LinkA{read\_river\_sp}{read.Rul.river.Rul.sp}()},
\code{\LinkA{read\_river}{read.Rul.river}()},
\code{\LinkA{read\_rivseg}{read.Rul.rivseg}()},
\code{\LinkA{read\_soil}{read.Rul.soil}()},
\code{\LinkA{write\_config}{write.Rul.config}()},
\code{\LinkA{write\_df}{write.Rul.df}()},
\code{\LinkA{write\_forc}{write.Rul.forc}()},
\code{\LinkA{write\_ic}{write.Rul.ic}()},
\code{\LinkA{write\_mesh}{write.Rul.mesh}()},
\code{\LinkA{write\_river}{write.Rul.river}()}
\end{SeeAlso}
%
\begin{Examples}
\begin{ExampleCode}
## Not run: 
forc <- read_forc_fn("model.forc")

## End(Not run)
\end{ExampleCode}
\end{Examples}
\inputencoding{utf8}
\HeaderA{read\_geol}{Read SHUD geology file}{read.Rul.geol}
\keyword{shud-io}{read\_geol}
%
\begin{Description}\relax
Reads a SHUD geology file (.geol) and returns geology data.
\end{Description}
%
\begin{Usage}
\begin{verbatim}
read_geol(file = shud.filein()["md.geol"])
\end{verbatim}
\end{Usage}
%
\begin{Arguments}
\begin{ldescription}
\item[\code{file}] Character. Full path to the .geol file
\end{ldescription}
\end{Arguments}
%
\begin{Value}
Matrix containing geology data
\end{Value}
%
\begin{SeeAlso}\relax
Other shud-io: 
\code{\LinkA{read\_att}{read.Rul.att}()},
\code{\LinkA{read\_calib}{read.Rul.calib}()},
\code{\LinkA{read\_config}{read.Rul.config}()},
\code{\LinkA{read\_df}{read.Rul.df}()},
\code{\LinkA{read\_forc\_csv}{read.Rul.forc.Rul.csv}()},
\code{\LinkA{read\_forc\_fn}{read.Rul.forc.Rul.fn}()},
\code{\LinkA{read\_ic}{read.Rul.ic}()},
\code{\LinkA{read\_lc}{read.Rul.lc}()},
\code{\LinkA{read\_mesh}{read.Rul.mesh}()},
\code{\LinkA{read\_para}{read.Rul.para}()},
\code{\LinkA{read\_river\_sp}{read.Rul.river.Rul.sp}()},
\code{\LinkA{read\_river}{read.Rul.river}()},
\code{\LinkA{read\_rivseg}{read.Rul.rivseg}()},
\code{\LinkA{read\_soil}{read.Rul.soil}()},
\code{\LinkA{write\_config}{write.Rul.config}()},
\code{\LinkA{write\_df}{write.Rul.df}()},
\code{\LinkA{write\_forc}{write.Rul.forc}()},
\code{\LinkA{write\_ic}{write.Rul.ic}()},
\code{\LinkA{write\_mesh}{write.Rul.mesh}()},
\code{\LinkA{write\_river}{write.Rul.river}()}
\end{SeeAlso}
%
\begin{Examples}
\begin{ExampleCode}
## Not run: 
geol <- read_geol("model.geol")

## End(Not run)
\end{ExampleCode}
\end{Examples}
\inputencoding{utf8}
\HeaderA{read\_ic}{Read SHUD initial conditions file}{read.Rul.ic}
\keyword{shud-io}{read\_ic}
%
\begin{Description}\relax
Reads a SHUD initial conditions file (.ic) and returns initial condition data.
\end{Description}
%
\begin{Usage}
\begin{verbatim}
read_ic(file = shud.filein()["md.ic"])
\end{verbatim}
\end{Usage}
%
\begin{Arguments}
\begin{ldescription}
\item[\code{file}] Character. Full path to the .ic file
\end{ldescription}
\end{Arguments}
%
\begin{Value}
List containing initial condition data
\end{Value}
%
\begin{SeeAlso}\relax
Other shud-io: 
\code{\LinkA{read\_att}{read.Rul.att}()},
\code{\LinkA{read\_calib}{read.Rul.calib}()},
\code{\LinkA{read\_config}{read.Rul.config}()},
\code{\LinkA{read\_df}{read.Rul.df}()},
\code{\LinkA{read\_forc\_csv}{read.Rul.forc.Rul.csv}()},
\code{\LinkA{read\_forc\_fn}{read.Rul.forc.Rul.fn}()},
\code{\LinkA{read\_geol}{read.Rul.geol}()},
\code{\LinkA{read\_lc}{read.Rul.lc}()},
\code{\LinkA{read\_mesh}{read.Rul.mesh}()},
\code{\LinkA{read\_para}{read.Rul.para}()},
\code{\LinkA{read\_river\_sp}{read.Rul.river.Rul.sp}()},
\code{\LinkA{read\_river}{read.Rul.river}()},
\code{\LinkA{read\_rivseg}{read.Rul.rivseg}()},
\code{\LinkA{read\_soil}{read.Rul.soil}()},
\code{\LinkA{write\_config}{write.Rul.config}()},
\code{\LinkA{write\_df}{write.Rul.df}()},
\code{\LinkA{write\_forc}{write.Rul.forc}()},
\code{\LinkA{write\_ic}{write.Rul.ic}()},
\code{\LinkA{write\_mesh}{write.Rul.mesh}()},
\code{\LinkA{write\_river}{write.Rul.river}()}
\end{SeeAlso}
%
\begin{Examples}
\begin{ExampleCode}
## Not run: 
ic <- read_ic("model.ic")

## End(Not run)
\end{ExampleCode}
\end{Examples}
\inputencoding{utf8}
\HeaderA{read\_lai}{Read LAI time series file}{read.Rul.lai}
\keyword{timeseries-io}{read\_lai}
%
\begin{Description}\relax
Reads a LAI (Leaf Area Index) time series file (.ts.lai) and returns
properly named time series components.
\end{Description}
%
\begin{Usage}
\begin{verbatim}
read_lai(file = shud.filein()["ts.lai"])
\end{verbatim}
\end{Usage}
%
\begin{Arguments}
\begin{ldescription}
\item[\code{file}] Character. Full path to LAI time series file
\end{ldescription}
\end{Arguments}
%
\begin{Value}
List of time series data with LAI and RL (Roughness Length) components
\end{Value}
%
\begin{SeeAlso}\relax
Other timeseries-io: 
\code{\LinkA{read\_tsd}{read.Rul.tsd}()},
\code{\LinkA{ts\_to\_daily}{ts.Rul.to.Rul.daily}()},
\code{\LinkA{ts\_to\_df}{ts.Rul.to.Rul.df}()},
\code{\LinkA{write\_tsd}{write.Rul.tsd}()},
\code{\LinkA{write\_xts}{write.Rul.xts}()}
\end{SeeAlso}
%
\begin{Examples}
\begin{ExampleCode}
## Not run: 
lai_data <- read_lai("model.ts.lai")
plot(lai_data$LAI)
plot(lai_data$RL)

## End(Not run)
\end{ExampleCode}
\end{Examples}
\inputencoding{utf8}
\HeaderA{read\_lc}{Read SHUD land cover file}{read.Rul.lc}
\keyword{shud-io}{read\_lc}
%
\begin{Description}\relax
Reads a SHUD land cover file (.lc) and returns land cover data.
\end{Description}
%
\begin{Usage}
\begin{verbatim}
read_lc(file = shud.filein()["md.lc"])
\end{verbatim}
\end{Usage}
%
\begin{Arguments}
\begin{ldescription}
\item[\code{file}] Character. Full path to the .lc file
\end{ldescription}
\end{Arguments}
%
\begin{Value}
Matrix containing land cover data
\end{Value}
%
\begin{SeeAlso}\relax
Other shud-io: 
\code{\LinkA{read\_att}{read.Rul.att}()},
\code{\LinkA{read\_calib}{read.Rul.calib}()},
\code{\LinkA{read\_config}{read.Rul.config}()},
\code{\LinkA{read\_df}{read.Rul.df}()},
\code{\LinkA{read\_forc\_csv}{read.Rul.forc.Rul.csv}()},
\code{\LinkA{read\_forc\_fn}{read.Rul.forc.Rul.fn}()},
\code{\LinkA{read\_geol}{read.Rul.geol}()},
\code{\LinkA{read\_ic}{read.Rul.ic}()},
\code{\LinkA{read\_mesh}{read.Rul.mesh}()},
\code{\LinkA{read\_para}{read.Rul.para}()},
\code{\LinkA{read\_river\_sp}{read.Rul.river.Rul.sp}()},
\code{\LinkA{read\_river}{read.Rul.river}()},
\code{\LinkA{read\_rivseg}{read.Rul.rivseg}()},
\code{\LinkA{read\_soil}{read.Rul.soil}()},
\code{\LinkA{write\_config}{write.Rul.config}()},
\code{\LinkA{write\_df}{write.Rul.df}()},
\code{\LinkA{write\_forc}{write.Rul.forc}()},
\code{\LinkA{write\_ic}{write.Rul.ic}()},
\code{\LinkA{write\_mesh}{write.Rul.mesh}()},
\code{\LinkA{write\_river}{write.Rul.river}()}
\end{SeeAlso}
%
\begin{Examples}
\begin{ExampleCode}
## Not run: 
lc <- read_lc("model.lc")

## End(Not run)
\end{ExampleCode}
\end{Examples}
\inputencoding{utf8}
\HeaderA{read\_mesh}{Read SHUD mesh file}{read.Rul.mesh}
\keyword{shud-io}{read\_mesh}
%
\begin{Description}\relax
Reads a SHUD mesh file (.mesh) and returns a SHUD.MESH object.
\end{Description}
%
\begin{Usage}
\begin{verbatim}
read_mesh(file = shud.filein()["md.mesh"])
\end{verbatim}
\end{Usage}
%
\begin{Arguments}
\begin{ldescription}
\item[\code{file}] Character. Full path to the .mesh file
\end{ldescription}
\end{Arguments}
%
\begin{Value}
SHUD.MESH object containing mesh and point data
\end{Value}
%
\begin{SeeAlso}\relax
Other shud-io: 
\code{\LinkA{read\_att}{read.Rul.att}()},
\code{\LinkA{read\_calib}{read.Rul.calib}()},
\code{\LinkA{read\_config}{read.Rul.config}()},
\code{\LinkA{read\_df}{read.Rul.df}()},
\code{\LinkA{read\_forc\_csv}{read.Rul.forc.Rul.csv}()},
\code{\LinkA{read\_forc\_fn}{read.Rul.forc.Rul.fn}()},
\code{\LinkA{read\_geol}{read.Rul.geol}()},
\code{\LinkA{read\_ic}{read.Rul.ic}()},
\code{\LinkA{read\_lc}{read.Rul.lc}()},
\code{\LinkA{read\_para}{read.Rul.para}()},
\code{\LinkA{read\_river\_sp}{read.Rul.river.Rul.sp}()},
\code{\LinkA{read\_river}{read.Rul.river}()},
\code{\LinkA{read\_rivseg}{read.Rul.rivseg}()},
\code{\LinkA{read\_soil}{read.Rul.soil}()},
\code{\LinkA{write\_config}{write.Rul.config}()},
\code{\LinkA{write\_df}{write.Rul.df}()},
\code{\LinkA{write\_forc}{write.Rul.forc}()},
\code{\LinkA{write\_ic}{write.Rul.ic}()},
\code{\LinkA{write\_mesh}{write.Rul.mesh}()},
\code{\LinkA{write\_river}{write.Rul.river}()}
\end{SeeAlso}
%
\begin{Examples}
\begin{ExampleCode}
## Not run: 
mesh <- read_mesh("model.mesh")

## End(Not run)
\end{ExampleCode}
\end{Examples}
\inputencoding{utf8}
\HeaderA{read\_nc\_data}{Read NetCDF data with spatial subsetting}{read.Rul.nc.Rul.data}
\keyword{netcdf-io}{read\_nc\_data}
%
\begin{Description}\relax
Reads NetCDF data with optional spatial and variable subsetting.
Supports both legacy raster and modern terra output formats.
\end{Description}
%
\begin{Usage}
\begin{verbatim}
read_nc_data(
  ncid,
  variables = NULL,
  extent = NULL,
  format = c("modern", "legacy")
)
\end{verbatim}
\end{Usage}
%
\begin{Arguments}
\begin{ldescription}
\item[\code{ncid}] NetCDF connection object from ncdf4::nc\_open()

\item[\code{variables}] Character vector of variable names to read (NULL = all)

\item[\code{extent}] Numeric vector c(xmin, xmax, ymin, ymax) for spatial subset

\item[\code{format}] Character. Output format ("modern" uses terra, "legacy" uses raster)
\end{ldescription}
\end{Arguments}
%
\begin{Value}
List containing coordinates, data array, and time information
\end{Value}
%
\begin{SeeAlso}\relax
Other netcdf-io: 
\code{\LinkA{get\_nc\_info}{get.Rul.nc.Rul.info}()},
\code{\LinkA{nc\_to\_raster}{nc.Rul.to.Rul.raster}()},
\code{\LinkA{open\_nc\_file}{open.Rul.nc.Rul.file}()},
\code{\LinkA{read\_nc\_time}{read.Rul.nc.Rul.time}()}
\end{SeeAlso}
%
\begin{Examples}
\begin{ExampleCode}
## Not run: 
library(ncdf4)
nc <- nc_open("climate_data.nc")

# Read all variables
data <- read_nc_data(nc)

# Read specific variables with spatial subset
extent <- c(-180, -120, 30, 60)  # Western North America
data <- read_nc_data(nc, variables = c("temp", "precip"), extent = extent)

# Use modern terra format
data <- read_nc_data(nc, format = "modern")

nc_close(nc)

## End(Not run)
\end{ExampleCode}
\end{Examples}
\inputencoding{utf8}
\HeaderA{read\_nc\_time}{Read NetCDF time dimension}{read.Rul.nc.Rul.time}
\keyword{netcdf-io}{read\_nc\_time}
%
\begin{Description}\relax
Extracts time information from NetCDF files and converts to POSIXct objects.
Supports various time formats and units commonly used in climate data.
\end{Description}
%
\begin{Usage}
\begin{verbatim}
read_nc_time(ncid)
\end{verbatim}
\end{Usage}
%
\begin{Arguments}
\begin{ldescription}
\item[\code{ncid}] NetCDF connection object from ncdf4::nc\_open()
\end{ldescription}
\end{Arguments}
%
\begin{Value}
POSIXct vector of time stamps
\end{Value}
%
\begin{SeeAlso}\relax
Other netcdf-io: 
\code{\LinkA{get\_nc\_info}{get.Rul.nc.Rul.info}()},
\code{\LinkA{nc\_to\_raster}{nc.Rul.to.Rul.raster}()},
\code{\LinkA{open\_nc\_file}{open.Rul.nc.Rul.file}()},
\code{\LinkA{read\_nc\_data}{read.Rul.nc.Rul.data}()}
\end{SeeAlso}
%
\begin{Examples}
\begin{ExampleCode}
## Not run: 
library(ncdf4)
nc <- nc_open("climate_data.nc")
times <- read_nc_time(nc)
nc_close(nc)

## End(Not run)
\end{ExampleCode}
\end{Examples}
\inputencoding{utf8}
\HeaderA{read\_output}{Read SHUD model output file}{read.Rul.output}
\keyword{output-io}{read\_output}
%
\begin{Description}\relax
Reads binary output files from SHUD model with support for both legacy
and modern formats. Returns time series data as xts objects.
\end{Description}
%
\begin{Usage}
\begin{verbatim}
read_output(
  keyword = NULL,
  file = NULL,
  path = get_output_path(),
  version = get_model_version(),
  format = c("modern", "legacy"),
  ascii = FALSE
)
\end{verbatim}
\end{Usage}
%
\begin{Arguments}
\begin{ldescription}
\item[\code{keyword}] Character. Keyword of the output file (e.g., "eleygw", "rivqdown")

\item[\code{file}] Character. Full path to output file (optional if keyword provided)

\item[\code{path}] Character. Path to output directory (default: from environment)

\item[\code{version}] Numeric. Output file version (1.0 or 2.0+)

\item[\code{format}] Character. Output format ("modern" uses terra, "legacy" uses raster)

\item[\code{ascii}] Logical. If TRUE, convert binary to ASCII format (.csv)
\end{ldescription}
\end{Arguments}
%
\begin{Value}
xts object containing time series data
\end{Value}
%
\begin{SeeAlso}\relax
Other output-io: 
\code{\LinkA{get\_default\_variables}{get.Rul.default.Rul.variables}()},
\code{\LinkA{load\_output\_data}{load.Rul.output.Rul.data}()},
\code{\LinkA{read\_output\_header}{read.Rul.output.Rul.header}()}
\end{SeeAlso}
%
\begin{Examples}
\begin{ExampleCode}
## Not run: 
# Read groundwater elevation data
gw_data <- read_output("eleygw")

# Read with specific file path
gw_data <- read_output(file = "model.eleygw.dat")

# Use modern format with terra
gw_data <- read_output("eleygw", format = "modern")

## End(Not run)
\end{ExampleCode}
\end{Examples}
\inputencoding{utf8}
\HeaderA{read\_output\_header}{Read output file header information}{read.Rul.output.Rul.header}
\keyword{output-io}{read\_output\_header}
%
\begin{Description}\relax
Reads header information from SHUD output files for inspection.
\end{Description}
%
\begin{Usage}
\begin{verbatim}
read_output_header(keyword = NULL, file = NULL, path = get_output_path())
\end{verbatim}
\end{Usage}
%
\begin{Arguments}
\begin{ldescription}
\item[\code{keyword}] Character. Keyword of the output file

\item[\code{file}] Character. Full path to output file (optional)

\item[\code{path}] Character. Path to output directory
\end{ldescription}
\end{Arguments}
%
\begin{Value}
Character string containing header information
\end{Value}
%
\begin{SeeAlso}\relax
Other output-io: 
\code{\LinkA{get\_default\_variables}{get.Rul.default.Rul.variables}()},
\code{\LinkA{load\_output\_data}{load.Rul.output.Rul.data}()},
\code{\LinkA{read\_output}{read.Rul.output}()}
\end{SeeAlso}
%
\begin{Examples}
\begin{ExampleCode}
## Not run: 
header_info <- read_output_header("eleygw")

## End(Not run)
\end{ExampleCode}
\end{Examples}
\inputencoding{utf8}
\HeaderA{read\_para}{Read SHUD parameters file}{read.Rul.para}
\keyword{shud-io}{read\_para}
%
\begin{Description}\relax
Reads a SHUD parameters file (.para) and returns parameter data.
\end{Description}
%
\begin{Usage}
\begin{verbatim}
read_para(file = shud.filein()["md.para"])
\end{verbatim}
\end{Usage}
%
\begin{Arguments}
\begin{ldescription}
\item[\code{file}] Character. Full path to the .para file
\end{ldescription}
\end{Arguments}
%
\begin{Value}
Data frame containing parameter data
\end{Value}
%
\begin{SeeAlso}\relax
Other shud-io: 
\code{\LinkA{read\_att}{read.Rul.att}()},
\code{\LinkA{read\_calib}{read.Rul.calib}()},
\code{\LinkA{read\_config}{read.Rul.config}()},
\code{\LinkA{read\_df}{read.Rul.df}()},
\code{\LinkA{read\_forc\_csv}{read.Rul.forc.Rul.csv}()},
\code{\LinkA{read\_forc\_fn}{read.Rul.forc.Rul.fn}()},
\code{\LinkA{read\_geol}{read.Rul.geol}()},
\code{\LinkA{read\_ic}{read.Rul.ic}()},
\code{\LinkA{read\_lc}{read.Rul.lc}()},
\code{\LinkA{read\_mesh}{read.Rul.mesh}()},
\code{\LinkA{read\_river\_sp}{read.Rul.river.Rul.sp}()},
\code{\LinkA{read\_river}{read.Rul.river}()},
\code{\LinkA{read\_rivseg}{read.Rul.rivseg}()},
\code{\LinkA{read\_soil}{read.Rul.soil}()},
\code{\LinkA{write\_config}{write.Rul.config}()},
\code{\LinkA{write\_df}{write.Rul.df}()},
\code{\LinkA{write\_forc}{write.Rul.forc}()},
\code{\LinkA{write\_ic}{write.Rul.ic}()},
\code{\LinkA{write\_mesh}{write.Rul.mesh}()},
\code{\LinkA{write\_river}{write.Rul.river}()}
\end{SeeAlso}
%
\begin{Examples}
\begin{ExampleCode}
## Not run: 
para <- read_para("model.para")

## End(Not run)
\end{ExampleCode}
\end{Examples}
\inputencoding{utf8}
\HeaderA{read\_river}{Read SHUD river file}{read.Rul.river}
\keyword{shud-io}{read\_river}
%
\begin{Description}\relax
Reads a SHUD river file (.riv) and returns a SHUD.RIVER object.
\end{Description}
%
\begin{Usage}
\begin{verbatim}
read_river(file = shud.filein()["md.riv"])
\end{verbatim}
\end{Usage}
%
\begin{Arguments}
\begin{ldescription}
\item[\code{file}] Character. Full path to the .riv file
\end{ldescription}
\end{Arguments}
%
\begin{Value}
SHUD.RIVER object containing river network data
\end{Value}
%
\begin{SeeAlso}\relax
Other shud-io: 
\code{\LinkA{read\_att}{read.Rul.att}()},
\code{\LinkA{read\_calib}{read.Rul.calib}()},
\code{\LinkA{read\_config}{read.Rul.config}()},
\code{\LinkA{read\_df}{read.Rul.df}()},
\code{\LinkA{read\_forc\_csv}{read.Rul.forc.Rul.csv}()},
\code{\LinkA{read\_forc\_fn}{read.Rul.forc.Rul.fn}()},
\code{\LinkA{read\_geol}{read.Rul.geol}()},
\code{\LinkA{read\_ic}{read.Rul.ic}()},
\code{\LinkA{read\_lc}{read.Rul.lc}()},
\code{\LinkA{read\_mesh}{read.Rul.mesh}()},
\code{\LinkA{read\_para}{read.Rul.para}()},
\code{\LinkA{read\_river\_sp}{read.Rul.river.Rul.sp}()},
\code{\LinkA{read\_rivseg}{read.Rul.rivseg}()},
\code{\LinkA{read\_soil}{read.Rul.soil}()},
\code{\LinkA{write\_config}{write.Rul.config}()},
\code{\LinkA{write\_df}{write.Rul.df}()},
\code{\LinkA{write\_forc}{write.Rul.forc}()},
\code{\LinkA{write\_ic}{write.Rul.ic}()},
\code{\LinkA{write\_mesh}{write.Rul.mesh}()},
\code{\LinkA{write\_river}{write.Rul.river}()}
\end{SeeAlso}
%
\begin{Examples}
\begin{ExampleCode}
## Not run: 
river <- read_river("model.riv")

## End(Not run)
\end{ExampleCode}
\end{Examples}
\inputencoding{utf8}
\HeaderA{read\_river\_sp}{Read river shapefile}{read.Rul.river.Rul.sp}
\keyword{shud-io}{read\_river\_sp}
%
\begin{Description}\relax
Reads a river shapefile and returns a SpatialLines object.
Note: This function uses legacy spatial libraries and will be updated
in future versions to use sf.
\end{Description}
%
\begin{Usage}
\begin{verbatim}
read_river_sp(file = file.path(shud.filein()["inpath"], "gis", "river.shp"))
\end{verbatim}
\end{Usage}
%
\begin{Arguments}
\begin{ldescription}
\item[\code{file}] Character. Full path to river shapefile
\end{ldescription}
\end{Arguments}
%
\begin{Value}
SpatialLines object
\end{Value}
%
\begin{SeeAlso}\relax
Other shud-io: 
\code{\LinkA{read\_att}{read.Rul.att}()},
\code{\LinkA{read\_calib}{read.Rul.calib}()},
\code{\LinkA{read\_config}{read.Rul.config}()},
\code{\LinkA{read\_df}{read.Rul.df}()},
\code{\LinkA{read\_forc\_csv}{read.Rul.forc.Rul.csv}()},
\code{\LinkA{read\_forc\_fn}{read.Rul.forc.Rul.fn}()},
\code{\LinkA{read\_geol}{read.Rul.geol}()},
\code{\LinkA{read\_ic}{read.Rul.ic}()},
\code{\LinkA{read\_lc}{read.Rul.lc}()},
\code{\LinkA{read\_mesh}{read.Rul.mesh}()},
\code{\LinkA{read\_para}{read.Rul.para}()},
\code{\LinkA{read\_river}{read.Rul.river}()},
\code{\LinkA{read\_rivseg}{read.Rul.rivseg}()},
\code{\LinkA{read\_soil}{read.Rul.soil}()},
\code{\LinkA{write\_config}{write.Rul.config}()},
\code{\LinkA{write\_df}{write.Rul.df}()},
\code{\LinkA{write\_forc}{write.Rul.forc}()},
\code{\LinkA{write\_ic}{write.Rul.ic}()},
\code{\LinkA{write\_mesh}{write.Rul.mesh}()},
\code{\LinkA{write\_river}{write.Rul.river}()}
\end{SeeAlso}
%
\begin{Examples}
\begin{ExampleCode}
## Not run: 
river_sp <- read_river_sp("rivers.shp")

## End(Not run)
\end{ExampleCode}
\end{Examples}
\inputencoding{utf8}
\HeaderA{read\_rivseg}{Read SHUD river segments file}{read.Rul.rivseg}
\keyword{shud-io}{read\_rivseg}
%
\begin{Description}\relax
Reads a SHUD river segments file (.rivseg) and returns segment data.
\end{Description}
%
\begin{Usage}
\begin{verbatim}
read_rivseg(file = shud.filein()["md.rivseg"])
\end{verbatim}
\end{Usage}
%
\begin{Arguments}
\begin{ldescription}
\item[\code{file}] Character. Full path to the .rivseg file
\end{ldescription}
\end{Arguments}
%
\begin{Value}
Matrix containing river segment data
\end{Value}
%
\begin{SeeAlso}\relax
Other shud-io: 
\code{\LinkA{read\_att}{read.Rul.att}()},
\code{\LinkA{read\_calib}{read.Rul.calib}()},
\code{\LinkA{read\_config}{read.Rul.config}()},
\code{\LinkA{read\_df}{read.Rul.df}()},
\code{\LinkA{read\_forc\_csv}{read.Rul.forc.Rul.csv}()},
\code{\LinkA{read\_forc\_fn}{read.Rul.forc.Rul.fn}()},
\code{\LinkA{read\_geol}{read.Rul.geol}()},
\code{\LinkA{read\_ic}{read.Rul.ic}()},
\code{\LinkA{read\_lc}{read.Rul.lc}()},
\code{\LinkA{read\_mesh}{read.Rul.mesh}()},
\code{\LinkA{read\_para}{read.Rul.para}()},
\code{\LinkA{read\_river\_sp}{read.Rul.river.Rul.sp}()},
\code{\LinkA{read\_river}{read.Rul.river}()},
\code{\LinkA{read\_soil}{read.Rul.soil}()},
\code{\LinkA{write\_config}{write.Rul.config}()},
\code{\LinkA{write\_df}{write.Rul.df}()},
\code{\LinkA{write\_forc}{write.Rul.forc}()},
\code{\LinkA{write\_ic}{write.Rul.ic}()},
\code{\LinkA{write\_mesh}{write.Rul.mesh}()},
\code{\LinkA{write\_river}{write.Rul.river}()}
\end{SeeAlso}
%
\begin{Examples}
\begin{ExampleCode}
## Not run: 
rivseg <- read_rivseg("model.rivseg")

## End(Not run)
\end{ExampleCode}
\end{Examples}
\inputencoding{utf8}
\HeaderA{read\_soil}{Read SHUD soil file}{read.Rul.soil}
\keyword{shud-io}{read\_soil}
%
\begin{Description}\relax
Reads a SHUD soil file (.soil) and returns soil data.
\end{Description}
%
\begin{Usage}
\begin{verbatim}
read_soil(file = shud.filein()["md.soil"])
\end{verbatim}
\end{Usage}
%
\begin{Arguments}
\begin{ldescription}
\item[\code{file}] Character. Full path to the .soil file
\end{ldescription}
\end{Arguments}
%
\begin{Value}
Matrix containing soil data
\end{Value}
%
\begin{SeeAlso}\relax
Other shud-io: 
\code{\LinkA{read\_att}{read.Rul.att}()},
\code{\LinkA{read\_calib}{read.Rul.calib}()},
\code{\LinkA{read\_config}{read.Rul.config}()},
\code{\LinkA{read\_df}{read.Rul.df}()},
\code{\LinkA{read\_forc\_csv}{read.Rul.forc.Rul.csv}()},
\code{\LinkA{read\_forc\_fn}{read.Rul.forc.Rul.fn}()},
\code{\LinkA{read\_geol}{read.Rul.geol}()},
\code{\LinkA{read\_ic}{read.Rul.ic}()},
\code{\LinkA{read\_lc}{read.Rul.lc}()},
\code{\LinkA{read\_mesh}{read.Rul.mesh}()},
\code{\LinkA{read\_para}{read.Rul.para}()},
\code{\LinkA{read\_river\_sp}{read.Rul.river.Rul.sp}()},
\code{\LinkA{read\_river}{read.Rul.river}()},
\code{\LinkA{read\_rivseg}{read.Rul.rivseg}()},
\code{\LinkA{write\_config}{write.Rul.config}()},
\code{\LinkA{write\_df}{write.Rul.df}()},
\code{\LinkA{write\_forc}{write.Rul.forc}()},
\code{\LinkA{write\_ic}{write.Rul.ic}()},
\code{\LinkA{write\_mesh}{write.Rul.mesh}()},
\code{\LinkA{write\_river}{write.Rul.river}()}
\end{SeeAlso}
%
\begin{Examples}
\begin{ExampleCode}
## Not run: 
soil <- read_soil("model.soil")

## End(Not run)
\end{ExampleCode}
\end{Examples}
\inputencoding{utf8}
\HeaderA{read\_tsd}{Read time series data file}{read.Rul.tsd}
\keyword{timeseries-io}{read\_tsd}
%
\begin{Description}\relax
Reads a time series file with header containing (nrow, ncol, timestring).
This function handles SHUD-specific time series format with automatic
time parsing and xts object creation.
\end{Description}
%
\begin{Usage}
\begin{verbatim}
read_tsd(file)
\end{verbatim}
\end{Usage}
%
\begin{Arguments}
\begin{ldescription}
\item[\code{file}] Character. Full path to time series file
\end{ldescription}
\end{Arguments}
%
\begin{Value}
List of xts objects containing time series data
\end{Value}
%
\begin{SeeAlso}\relax
Other timeseries-io: 
\code{\LinkA{read\_lai}{read.Rul.lai}()},
\code{\LinkA{ts\_to\_daily}{ts.Rul.to.Rul.daily}()},
\code{\LinkA{ts\_to\_df}{ts.Rul.to.Rul.df}()},
\code{\LinkA{write\_tsd}{write.Rul.tsd}()},
\code{\LinkA{write\_xts}{write.Rul.xts}()}
\end{SeeAlso}
%
\begin{Examples}
\begin{ExampleCode}
## Not run: 
ts_data <- read_tsd("timeseries.tsd")
lai_data <- read_lai("model.ts.lai")

## End(Not run)
\end{ExampleCode}
\end{Examples}
\inputencoding{utf8}
\HeaderA{readatt}{Read the .att file \code{readmesh}}{readatt}
%
\begin{Description}\relax
Read the .att file
\code{readmesh}
\end{Description}
%
\begin{Usage}
\begin{verbatim}
readatt(file = shud.filein()["md.att"])
\end{verbatim}
\end{Usage}
%
\begin{Arguments}
\begin{ldescription}
\item[\code{file}] full path of file
\end{ldescription}
\end{Arguments}
%
\begin{Value}
matrix
\end{Value}
\inputencoding{utf8}
\HeaderA{readcalib}{Read the .calib file \code{readcalib}}{readcalib}
%
\begin{Description}\relax
Read the .calib file
\code{readcalib}
\end{Description}
%
\begin{Usage}
\begin{verbatim}
readcalib(file = shud.filein()["md.calib"])
\end{verbatim}
\end{Usage}
%
\begin{Arguments}
\begin{ldescription}
\item[\code{file}] full path of file
\end{ldescription}
\end{Arguments}
%
\begin{Value}
.calib
\end{Value}
\inputencoding{utf8}
\HeaderA{readconfig}{Read the configuration (.para, .calib) file \code{readconfig}}{readconfig}
%
\begin{Description}\relax
Read the configuration (.para, .calib) file
\code{readconfig}
\end{Description}
%
\begin{Usage}
\begin{verbatim}
readconfig(file = shud.filein()["md.para"])
\end{verbatim}
\end{Usage}
%
\begin{Arguments}
\begin{ldescription}
\item[\code{file}] full path of file
\end{ldescription}
\end{Arguments}
%
\begin{Value}
.para or .calib
\end{Value}
\inputencoding{utf8}
\HeaderA{readforc.csv}{Read the .csv files in .forc file \code{readforc.csv}}{readforc.csv}
%
\begin{Description}\relax
Read the .csv files in .forc file
\code{readforc.csv}
\end{Description}
%
\begin{Usage}
\begin{verbatim}
readforc.csv(file = shud.filein()["md.forc"], id = NULL)
\end{verbatim}
\end{Usage}
%
\begin{Arguments}
\begin{ldescription}
\item[\code{file}] full path of file

\item[\code{id}] Index of the forcing sites.  default = NULL, which return average values of all sites.
\end{ldescription}
\end{Arguments}
%
\begin{Value}
forcing data, list.
\end{Value}
\inputencoding{utf8}
\HeaderA{readforc.fn}{Read the .forc file \code{readlc}}{readforc.fn}
%
\begin{Description}\relax
Read the .forc file
\code{readlc}
\end{Description}
%
\begin{Usage}
\begin{verbatim}
readforc.fn(file = shud.filein()["md.forc"])
\end{verbatim}
\end{Usage}
%
\begin{Arguments}
\begin{ldescription}
\item[\code{file}] full path of file
\end{ldescription}
\end{Arguments}
%
\begin{Value}
data.frame of forcing sites.
\end{Value}
\inputencoding{utf8}
\HeaderA{readgeol}{Read the .geol file \code{readgeol}}{readgeol}
%
\begin{Description}\relax
Read the .geol file
\code{readgeol}
\end{Description}
%
\begin{Usage}
\begin{verbatim}
readgeol(file = shud.filein()["md.geol"])
\end{verbatim}
\end{Usage}
%
\begin{Arguments}
\begin{ldescription}
\item[\code{file}] full path of file
\end{ldescription}
\end{Arguments}
%
\begin{Value}
.geol
\end{Value}
\inputencoding{utf8}
\HeaderA{readic}{Read the .ic file \code{readic}}{readic}
%
\begin{Description}\relax
Read the .ic file
\code{readic}
\end{Description}
%
\begin{Usage}
\begin{verbatim}
readic(file = shud.filein()["md.ic"])
\end{verbatim}
\end{Usage}
%
\begin{Arguments}
\begin{ldescription}
\item[\code{file}] full path of file
\end{ldescription}
\end{Arguments}
%
\begin{Value}
.ic
\end{Value}
\inputencoding{utf8}
\HeaderA{readlai}{Read the .ts.lai Time-Series file \code{readlai}}{readlai}
%
\begin{Description}\relax
Read the .ts.lai Time-Series file
\code{readlai}
\end{Description}
%
\begin{Usage}
\begin{verbatim}
readlai(file = shud.filein()["ts.lai"])
\end{verbatim}
\end{Usage}
%
\begin{Arguments}
\begin{ldescription}
\item[\code{file}] full path of file
\end{ldescription}
\end{Arguments}
%
\begin{Value}
list of Time-Series data
\end{Value}
\inputencoding{utf8}
\HeaderA{readlc}{Read the .lc file \code{readlc}}{readlc}
%
\begin{Description}\relax
Read the .lc file
\code{readlc}
\end{Description}
%
\begin{Usage}
\begin{verbatim}
readlc(file = shud.filein()["md.lc"])
\end{verbatim}
\end{Usage}
%
\begin{Arguments}
\begin{ldescription}
\item[\code{file}] full path of file
\end{ldescription}
\end{Arguments}
%
\begin{Value}
.lc
\end{Value}
\inputencoding{utf8}
\HeaderA{readmesh}{Read the .mesh file \code{readmesh}}{readmesh}
%
\begin{Description}\relax
Read the .mesh file
\code{readmesh}
\end{Description}
%
\begin{Usage}
\begin{verbatim}
readmesh(file = shud.filein()["md.mesh"])
\end{verbatim}
\end{Usage}
%
\begin{Arguments}
\begin{ldescription}
\item[\code{file}] full path of .mesh file
\end{ldescription}
\end{Arguments}
%
\begin{Value}
Mesh domain
\end{Value}
\inputencoding{utf8}
\HeaderA{readnc}{Read the time tage from NC file.}{readnc}
%
\begin{Description}\relax
Read the time tage from NC file.
\end{Description}
%
\begin{Usage}
\begin{verbatim}
readnc(ncid, varid = NULL, ext = NULL)
\end{verbatim}
\end{Usage}
%
\begin{Arguments}
\begin{ldescription}
\item[\code{ncid}] the ncid from nc\_open()

\item[\code{varid}] The varid.

\item[\code{ext}] box coordinates of the subset.
\end{ldescription}
\end{Arguments}
%
\begin{Value}
a POSIXct date-time objects
\end{Value}
\inputencoding{utf8}
\HeaderA{readnc.time}{Read the time tage from NC file.}{readnc.time}
%
\begin{Description}\relax
Read the time tage from NC file.
\end{Description}
%
\begin{Usage}
\begin{verbatim}
readnc.time(ncid)
\end{verbatim}
\end{Usage}
%
\begin{Arguments}
\begin{ldescription}
\item[\code{ncid}] the ncid from nc\_open()
\end{ldescription}
\end{Arguments}
%
\begin{Value}
a POSIXct date-time objects
\end{Value}
\inputencoding{utf8}
\HeaderA{readout}{Read output file from SHUD model}{readout}
\keyword{Could}{readout}
\keyword{and}{readout}
\keyword{be}{readout}
\keyword{for}{readout}
\keyword{format.}{readout}
\keyword{mesh}{readout}
\keyword{output.}{readout}
\keyword{read}{readout}
\keyword{reading}{readout}
\keyword{river}{readout}
\keyword{used}{readout}
%
\begin{Description}\relax
Read output file from SHUD model
\end{Description}
%
\begin{Usage}
\begin{verbatim}
readout(
  keyword,
  path = get("outpath", envir = .shud),
  ASCII = FALSE,
  ver = get("version", envir = .shud),
  file = file.path(path, paste0(get("PRJNAME", envir = .shud), ".", keyword, ".dat"))
)
\end{verbatim}
\end{Usage}
%
\begin{Arguments}
\begin{ldescription}
\item[\code{keyword}] keyword of the output file. File name format is: projectname.keyword.dat

\item[\code{path}] path of the output file

\item[\code{ASCII}] If TRUE, convert the binary file to ASCII format (.csv)

\item[\code{ver}] Version of output file. Ver = 1.0, 2.0

\item[\code{file}] Full path of output file.
\end{ldescription}
\end{Arguments}
%
\begin{Value}
A Time-Series data. This list require support of xts packages.
\end{Value}
\inputencoding{utf8}
\HeaderA{readout.header}{Read header information of output file}{readout.header}
\keyword{Could}{readout.header}
\keyword{and}{readout.header}
\keyword{be}{readout.header}
\keyword{for}{readout.header}
\keyword{format.}{readout.header}
\keyword{mesh}{readout.header}
\keyword{output.}{readout.header}
\keyword{read}{readout.header}
\keyword{reading}{readout.header}
\keyword{river}{readout.header}
\keyword{used}{readout.header}
%
\begin{Description}\relax
Read header information of output file
\end{Description}
%
\begin{Usage}
\begin{verbatim}
readout.header(
  keyword,
  path = get("outpath", envir = .shud),
  file = file.path(path, paste0(get("PRJNAME", envir = .shud), ".", keyword, ".dat"))
)
\end{verbatim}
\end{Usage}
%
\begin{Arguments}
\begin{ldescription}
\item[\code{keyword}] keyword of the output file. File name format is: projectname.keyword.dat

\item[\code{path}] path of the output file

\item[\code{file}] Full path of output file.
\end{ldescription}
\end{Arguments}
%
\begin{Value}
Header information
\end{Value}
\inputencoding{utf8}
\HeaderA{readpara}{Read the .para file \code{readpara}}{readpara}
%
\begin{Description}\relax
Read the .para file
\code{readpara}
\end{Description}
%
\begin{Usage}
\begin{verbatim}
readpara(file = shud.filein()["md.para"])
\end{verbatim}
\end{Usage}
%
\begin{Arguments}
\begin{ldescription}
\item[\code{file}] full path of file
\end{ldescription}
\end{Arguments}
%
\begin{Value}
.para
\end{Value}
\inputencoding{utf8}
\HeaderA{readriv}{Read the .riv file \code{readriv}}{readriv}
%
\begin{Description}\relax
Read the .riv file
\code{readriv}
\end{Description}
%
\begin{Usage}
\begin{verbatim}
readriv(file = shud.filein()["md.riv"])
\end{verbatim}
\end{Usage}
%
\begin{Arguments}
\begin{ldescription}
\item[\code{file}] full path of .riv file
\end{ldescription}
\end{Arguments}
%
\begin{Value}
SHUD.RIVER
\end{Value}
\inputencoding{utf8}
\HeaderA{readriv.sp}{Read the river shapefile \code{readriv.sp}}{readriv.sp}
%
\begin{Description}\relax
Read the river shapefile
\code{readriv.sp}
\end{Description}
%
\begin{Usage}
\begin{verbatim}
readriv.sp(file = file.path(shud.filein()["inpath"], "gis", "river.shp"))
\end{verbatim}
\end{Usage}
%
\begin{Arguments}
\begin{ldescription}
\item[\code{file}] full path of spatialLines file
\end{ldescription}
\end{Arguments}
%
\begin{Value}
spatialLines file
\end{Value}
\inputencoding{utf8}
\HeaderA{readrivseg}{Read the .rivseg file \code{readrivseg}}{readrivseg}
%
\begin{Description}\relax
Read the .rivseg file
\code{readrivseg}
\end{Description}
%
\begin{Usage}
\begin{verbatim}
readrivseg(file = shud.filein()["md.rivseg"])
\end{verbatim}
\end{Usage}
%
\begin{Arguments}
\begin{ldescription}
\item[\code{file}] full path of .rivseg file
\end{ldescription}
\end{Arguments}
%
\begin{Value}
matrix
\end{Value}
\inputencoding{utf8}
\HeaderA{readsoil}{Read the .soil file \code{readsoil}}{readsoil}
%
\begin{Description}\relax
Read the .soil file
\code{readsoil}
\end{Description}
%
\begin{Usage}
\begin{verbatim}
readsoil(file = shud.filein()["md.soil"])
\end{verbatim}
\end{Usage}
%
\begin{Arguments}
\begin{ldescription}
\item[\code{file}] full path of file
\end{ldescription}
\end{Arguments}
%
\begin{Value}
.soil
\end{Value}
\inputencoding{utf8}
\HeaderA{removeholes}{Remove the holes in polygons \code{removeholes}}{removeholes}
%
\begin{Description}\relax
Remove the holes in polygons
\code{removeholes}
\end{Description}
%
\begin{Usage}
\begin{verbatim}
removeholes(sp)
\end{verbatim}
\end{Usage}
%
\begin{Arguments}
\begin{ldescription}
\item[\code{sp}] SpatialPolygons or SpatialPolygonDataFrame
\end{ldescription}
\end{Arguments}
%
\begin{Value}
Raster map
\end{Value}
%
\begin{Examples}
\begin{ExampleCode}
library(sp)
p.out = Polygon(cbind(c(4,4,6,7,4),c(5,3,2,5,5))  )
p.hole = Polygon(cbind(c(5,6,6,5,5),c(4,4,3,3,4) ), hole = TRUE)
sp <- SpatialPolygons(list(Polygons(list(p.out, p.hole), "1")))
s = removeholes(sp)
par(mfrow = c(1, 2))
plot(sp)
plot(s)
\end{ExampleCode}
\end{Examples}
\inputencoding{utf8}
\HeaderA{river\_network}{River Network Functions}{river.Rul.network}
%
\begin{Description}\relax
Functions for building and analyzing river networks, including
integrated property calculations.
\end{Description}
\inputencoding{utf8}
\HeaderA{river\_processing}{River Processing Functions}{river.Rul.processing}
%
\begin{Description}\relax
Functions for processing river networks including order calculation,
path analysis, and topology determination using modern sf library.
\end{Description}
\inputencoding{utf8}
\HeaderA{RiverAtt}{Summary of the attributions on Rivers \code{RiverAtt}}{RiverAtt}
%
\begin{Description}\relax
Summary of the attributions on Rivers
\code{RiverAtt}
\end{Description}
%
\begin{Usage}
\begin{verbatim}
RiverAtt(riv = readriv())
\end{verbatim}
\end{Usage}
%
\begin{Arguments}
\begin{ldescription}
\item[\code{riv}] SHUD.RIVER
\end{ldescription}
\end{Arguments}
%
\begin{Value}
Attributes of each river reach
\end{Value}
\inputencoding{utf8}
\HeaderA{RiverLength}{Calculate the river length \code{RiverLength}}{RiverLength}
%
\begin{Description}\relax
Calculate the river length
\code{RiverLength}
\end{Description}
%
\begin{Usage}
\begin{verbatim}
RiverLength(pr)
\end{verbatim}
\end{Usage}
%
\begin{Arguments}
\begin{ldescription}
\item[\code{pr}] SHUD.RIVER class
\end{ldescription}
\end{Arguments}
%
\begin{Value}
River length, vector
\end{Value}
\inputencoding{utf8}
\HeaderA{RiverSlope}{Calculate the river slope \code{RiverSlope}}{RiverSlope}
%
\begin{Description}\relax
Calculate the river slope
\code{RiverSlope}
\end{Description}
%
\begin{Usage}
\begin{verbatim}
RiverSlope(pr)
\end{verbatim}
\end{Usage}
%
\begin{Arguments}
\begin{ldescription}
\item[\code{pr}] SHUD.RIVER class
\end{ldescription}
\end{Arguments}
%
\begin{Value}
River slopes, vector
\end{Value}
\inputencoding{utf8}
\HeaderA{RiverType}{parameters for river types \code{RiverType}}{RiverType}
%
\begin{Description}\relax
parameters for river types
\code{RiverType}
\end{Description}
%
\begin{Usage}
\begin{verbatim}
RiverType(n, width = 2 * (1:n))
\end{verbatim}
\end{Usage}
%
\begin{Arguments}
\begin{ldescription}
\item[\code{n}] number of types

\item[\code{width}] Width of rivers
\end{ldescription}
\end{Arguments}
%
\begin{Value}
data.frame of parameters
\end{Value}
\inputencoding{utf8}
\HeaderA{rmDuplicatedLines}{Remove the duplicated lines, which means the FROM and TO points are identical. \code{rmDuplicatedLines}}{rmDuplicatedLines}
%
\begin{Description}\relax
Remove the duplicated lines, which means the FROM and TO points are identical.
\code{rmDuplicatedLines}
\end{Description}
%
\begin{Usage}
\begin{verbatim}
rmDuplicatedLines(x, ...)
\end{verbatim}
\end{Usage}
%
\begin{Arguments}
\begin{ldescription}
\item[\code{x}] ShapeLine*

\item[\code{...}] More options in duplicated()
\end{ldescription}
\end{Arguments}
%
\begin{Value}
ShapeLine* without duplicated lines.
\end{Value}
\inputencoding{utf8}
\HeaderA{rowMatch}{Identify whether every row of x is same as m;}{rowMatch}
%
\begin{Description}\relax
Identify whether every row of x is same as m;
\end{Description}
%
\begin{Usage}
\begin{verbatim}
rowMatch(x, m)
\end{verbatim}
\end{Usage}
%
\begin{Arguments}
\begin{ldescription}
\item[\code{x}] Row vector

\item[\code{m}] Matrix
\end{ldescription}
\end{Arguments}
%
\begin{Value}
TRUE/FALSE of matching
\end{Value}
%
\begin{Examples}
\begin{ExampleCode}
rowMatch(c(1, 1), cbind(1:10, 1:10))
\end{ExampleCode}
\end{Examples}
\inputencoding{utf8}
\HeaderA{sac}{This is data of one of sub-catchmetn in Sacramento Watershed  to be included in my package}{sac}
\keyword{data}{sac}
%
\begin{Description}\relax
This is data of one of sub-catchmetn in Sacramento Watershed  to be included in my package
\end{Description}
\inputencoding{utf8}
\HeaderA{sh}{This is data to be included in my package}{sh}
\keyword{data}{sh}
%
\begin{Description}\relax
This is data to be included in my package
\end{Description}
\inputencoding{utf8}
\HeaderA{SharedPoints}{Find Shared points in SpatialLine* \code{SharedPoints}}{SharedPoints}
%
\begin{Description}\relax
Find Shared points in SpatialLine*
\code{SharedPoints}
\end{Description}
%
\begin{Usage}
\begin{verbatim}
SharedPoints(sp, plot = FALSE)
\end{verbatim}
\end{Usage}
%
\begin{Arguments}
\begin{ldescription}
\item[\code{sp}] SpatialLines

\item[\code{plot}] Whether to plot the results for debugging
\end{ldescription}
\end{Arguments}
%
\begin{Value}
Topologic relationship between components of SpatialLines.
\end{Value}
\inputencoding{utf8}
\HeaderA{SHUD River-class}{Class of SHUD.RIVER}{SHUD River.Rdash.class}
\aliasA{SHUD.RIVER}{SHUD River-class}{SHUD.RIVER}
%
\begin{Description}\relax
S4 class representing a river network for SHUD models.
The class stores river segment properties, hydraulic parameters,
and node information. Now supports both legacy data.frame format
and modern sf spatial format.
\end{Description}
%
\begin{Value}
Class of SHUD.RIVER
\end{Value}
%
\begin{Section}{Slots}

\begin{description}

\item[\code{river}] data.frame with river segment properties

\item[\code{rivertype}] data.frame with hydraulic parameters for each river type

\item[\code{point}] data.frame with from/to node coordinates and elevations

\item[\code{network}] sf object with river network geometry (optional, for modern format)

\item[\code{crs}] character string with coordinate reference system (optional)

\end{description}
\end{Section}
%
\begin{Section}{River Data Frame}

The river slot contains:
\begin{itemize}

\item{} Index: Segment ID
\item{} Down: Downstream segment index (-1 for outlets)
\item{} Type: River type/order
\item{} Slope: Bed slope
\item{} Length: Segment length
\item{} BC: Boundary condition flag

\end{itemize}

\end{Section}
%
\begin{Section}{River Type Data Frame}

The rivertype slot contains hydraulic parameters:
\begin{itemize}

\item{} Index: Type index
\item{} Depth: Channel depth (m)
\item{} BankSlope: Bank slope
\item{} Width: Channel width (m)
\item{} Sinuosity: Channel sinuosity
\item{} Manning: Manning's roughness coefficient
\item{} Cwr: Width-discharge coefficient
\item{} KsatH: Horizontal hydraulic conductivity
\item{} BedThick: Bed sediment thickness

\end{itemize}

\end{Section}
%
\begin{Section}{Point Data Frame}

The point slot contains node information:
\begin{itemize}

\item{} From.x, From.y, From.z: Start node coordinates and elevation
\item{} To.x, To.y, To.z: End node coordinates and elevation

\end{itemize}

\end{Section}
%
\begin{Section}{Modern Format}

When created with modern functions (e.g., \code{build\_river\_network()}),
the network slot contains an sf object with LINESTRING geometry and
all river properties as attributes. The crs slot stores the coordinate
reference system.
\end{Section}
\inputencoding{utf8}
\HeaderA{shud-io}{SHUD File I/O Functions}{shud.Rdash.io}
%
\begin{Description}\relax
This module provides modern functions for reading and writing SHUD model files.
All functions use snake\_case naming convention and provide comprehensive
error handling and validation.
\end{Description}
\inputencoding{utf8}
\HeaderA{shud.att}{Calculate mesh attributes from raster data}{shud.att}
%
\begin{Description}\relax
Extracts attribute values from raster layers to mesh triangle centroids.
Modernized implementation using terra for raster operations with improved
performance, with automatic raster to terra conversion.
\end{Description}
%
\begin{Usage}
\begin{verbatim}
shud.att(
  tri,
  r.soil = NULL,
  r.geol = NULL,
  r.lc = NULL,
  r.forc = NULL,
  r.mf = NULL,
  r.BC = NULL,
  r.SS = NULL,
  sp.lake = NULL
)

shud.att(
  tri,
  r.soil = NULL,
  r.geol = NULL,
  r.lc = NULL,
  r.forc = NULL,
  r.mf = NULL,
  r.BC = NULL,
  r.SS = NULL,
  sp.lake = NULL
)
\end{verbatim}
\end{Usage}
%
\begin{Arguments}
\begin{ldescription}
\item[\code{tri}] Triangles defination

\item[\code{r.soil}] raster of soil classification

\item[\code{r.geol}] raster of geology layer

\item[\code{r.lc}] raster of land cover, LAI and Roughness Length

\item[\code{r.forc}] raster of forcing data

\item[\code{r.mf}] raster of melt factor

\item[\code{r.BC}] raster of boundary condition

\item[\code{r.SS}] raster of Source/Sink

\item[\code{sp.lake}] SpatialPolygon of lake layers.
\end{ldescription}
\end{Arguments}
%
\begin{Details}\relax
This function has been modernized to use terra for raster operations,
providing 2-5x better performance. Legacy raster and sp objects are
automatically converted to terra/sf formats.
\end{Details}
%
\begin{Value}
data.frame with attribute indices for each mesh triangle

data.frame of SHUD .att
\end{Value}
%
\begin{Examples}
\begin{ExampleCode}
## Not run: 
library(sf)
library(terra)

# Generate mesh
boundary <- sf::st_as_sf(
  sf::st_sfc(sf::st_polygon(list(matrix(c(0,0, 10,0, 10,10, 0,10, 0,0),
                                         ncol=2, byrow=TRUE))))
)
tri <- shud.triangle(wb = boundary, q = 30, a = 2)

# Create attribute rasters
soil <- terra::rast(ncol=10, nrow=10, xmin=0, xmax=10, ymin=0, ymax=10)
terra::values(soil) <- sample(1:5, 100, replace=TRUE)

# Calculate attributes
att <- shud.att(tri, r.soil = soil)

## End(Not run)
\end{ExampleCode}
\end{Examples}
\inputencoding{utf8}
\HeaderA{shud.calib}{Generate the default model calibration \code{shud.calib}}{shud.calib}
%
\begin{Description}\relax
Generate the default model calibration
\code{shud.calib}
\end{Description}
%
\begin{Usage}
\begin{verbatim}
shud.calib()
\end{verbatim}
\end{Usage}
%
\begin{Value}
Default calibration values for model
\end{Value}
\inputencoding{utf8}
\HeaderA{shud.env}{Setup .shud environment. \code{shud.env}}{shud.env}
%
\begin{Description}\relax
Setup .shud environment.
\code{shud.env}
\end{Description}
%
\begin{Usage}
\begin{verbatim}
shud.env(
  prjname,
  inpath = file.path("input", prjname),
  outpath = file.path("output", paste0(prjname, ".out")),
  anapath = file.path(outpath, "postprocessing"),
  version = 2
)
\end{verbatim}
\end{Usage}
%
\begin{Arguments}
\begin{ldescription}
\item[\code{prjname}] character of SHUD model project name.

\item[\code{inpath}] path of SHUD model input directory

\item[\code{outpath}] path of SHUD model output directory

\item[\code{anapath}] character of SHUD model project name.

\item[\code{version}] version of SHUD model.
\end{ldescription}
\end{Arguments}
\inputencoding{utf8}
\HeaderA{shud.filein}{Prepare Input file of SHUD model \code{shud.filein}}{shud.filein}
%
\begin{Description}\relax
Prepare Input file of SHUD model
\code{shud.filein}
\end{Description}
%
\begin{Usage}
\begin{verbatim}
shud.filein(
  projname = get("PRJNAME", envir = .shud),
  inpath = get("inpath", envir = .shud),
  outpath = get("outpath", envir = .shud)
)
\end{verbatim}
\end{Usage}
%
\begin{Arguments}
\begin{ldescription}
\item[\code{projname}] Character, project name, default= PRJNAME which is a global variable.

\item[\code{inpath}] SHUD model input directory, default = inpath which is global variable

\item[\code{outpath}] SHUD model output directory, default = outpath which is global variable
\end{ldescription}
\end{Arguments}
%
\begin{Value}
Character of full path of input files for SHUD model
\end{Value}
\inputencoding{utf8}
\HeaderA{shud.ic}{extract Coordinates of SpatialLines or  SpatialPolygons \code{shud.ic}}{shud.ic}
%
\begin{Description}\relax
extract Coordinates of SpatialLines or  SpatialPolygons
\code{shud.ic}
\end{Description}
%
\begin{Usage}
\begin{verbatim}
shud.ic(
  ncell,
  nriv,
  AqD = 10,
  stage = 0.1,
  p1 = 0.4,
  p2 = p1,
  lakestage = NULL
)
\end{verbatim}
\end{Usage}
%
\begin{Arguments}
\begin{ldescription}
\item[\code{ncell}] Number of triangles in SHUD.MESH domain

\item[\code{nriv}] Number of river lines

\item[\code{AqD}] Default aquifer depth

\item[\code{stage}] Default river stage

\item[\code{p1}] Fraction of unsat zone, \LinkA{0,1}{0,1}

\item[\code{p2}] Fraction of Groundwater, \LinkA{0,1}{0,1}

\item[\code{lakestage}] Lake Stages. Null=No lake.
\end{ldescription}
\end{Arguments}
%
\begin{Value}
List of initial condition
\end{Value}
\inputencoding{utf8}
\HeaderA{shud.mask}{Generate the raster mask of Mesh domain \code{shud.mask}}{shud.mask}
%
\begin{Description}\relax
Generate the raster mask of Mesh domain
\code{shud.mask}
\end{Description}
%
\begin{Usage}
\begin{verbatim}
shud.mask(
  pm = readmesh(),
  proj = NULL,
  rr = get("MASK", envir = .shud),
  n = 10000,
  cellsize = NULL
)
\end{verbatim}
\end{Usage}
%
\begin{Arguments}
\begin{ldescription}
\item[\code{pm}] \code{shud.mesh}

\item[\code{proj}] Projection parameter

\item[\code{rr}] Default mask in .shud environment

\item[\code{n}] Number of grid

\item[\code{cellsize}] Resolution, defaul = NULL, resolution = extent / ngrids
\end{ldescription}
\end{Arguments}
%
\begin{Value}
Raster map
\end{Value}
\inputencoding{utf8}
\HeaderA{shud.mesh}{Generate SHUD mesh domain from triangulation}{shud.mesh}
%
\begin{Description}\relax
Creates a SHUD.MESH S4 object from a triangulation, extracting elevation
and aquifer depth from raster data. Modernized implementation using terra
for raster operations, with automatic raster to terra conversion.
\end{Description}
%
\begin{Usage}
\begin{verbatim}
shud.mesh(tri, dem, AqDepth = 10, r.aq = dem * 0 + AqDepth)

shud.mesh(tri, dem, AqDepth = 10, r.aq = dem * 0 + AqDepth)
\end{verbatim}
\end{Usage}
%
\begin{Arguments}
\begin{ldescription}
\item[\code{tri}] Triangles defination

\item[\code{dem}] Elevation. Projection of the DEM raster must be same as watershed boundary.

\item[\code{AqDepth}] Aquifer depth, numeric.

\item[\code{r.aq}] Aquifer Thickness. Raster object
\end{ldescription}
\end{Arguments}
%
\begin{Details}\relax
This function has been modernized to use terra for raster operations,
providing better performance. Legacy raster objects are automatically
converted to terra SpatRaster format.
\end{Details}
%
\begin{Value}
SHUD.MESH S4 object with mesh and point data

Triangle mesh of model domain.
\end{Value}
%
\begin{Examples}
\begin{ExampleCode}
## Not run: 
library(sf)
library(terra)

# Create boundary and generate mesh
boundary <- sf::st_as_sf(
  sf::st_sfc(sf::st_polygon(list(matrix(c(0,0, 10,0, 10,10, 0,10, 0,0),
                                         ncol=2, byrow=TRUE))))
)
tri <- shud.triangle(wb = boundary, q = 30, a = 2)

# Create DEM
dem <- terra::rast(ncol=10, nrow=10, xmin=0, xmax=10, ymin=0, ymax=10)
terra::values(dem) <- runif(100, 100, 200)

# Generate mesh domain
mesh <- shud.mesh(tri, dem = dem, AqDepth = 10)

## End(Not run)
library(raster)
data(sac)
wbd = sac[['wbd']]
dem = sac[['dem']]
a.max = 1e6 * 1;
q.min = 33;
tol.riv = 200
tol.wb = 200
tol.len = 500
AqDepth = 10

wbbuf = rgeos::gBuffer(wbd, width = 2000)
dem = raster::crop(dem, wbbuf)

wb.dis = rgeos::gUnionCascaded(wbd)
wb.simp = rgeos::gSimplify(wb.dis, tol = tol.wb, topologyPreserve = TRUE)


tri = shud.triangle(wb = wb.simp, q = q.min, a = a.max)
plot(tri, asp = 1)

# generate SHUD .mesh
pm = shud.mesh(tri, dem = dem, AqDepth = AqDepth)
sm = sp.mesh2Shape(pm)
raster::plot(sm)
\end{ExampleCode}
\end{Examples}
\inputencoding{utf8}
\HeaderA{shud.para}{Generate the default model configuration \code{shud.para}}{shud.para}
%
\begin{Description}\relax
Generate the default model configuration
\code{shud.para}
\end{Description}
%
\begin{Usage}
\begin{verbatim}
shud.para(nday = 10)
\end{verbatim}
\end{Usage}
%
\begin{Arguments}
\begin{ldescription}
\item[\code{nday}] Simulation periods (days)
\end{ldescription}
\end{Arguments}
%
\begin{Value}
model configureation, Vector
\end{Value}
\inputencoding{utf8}
\HeaderA{shud.river}{calculate river order, downstream, slope and length \code{shud.river}}{shud.river}
%
\begin{Description}\relax
calculate river order, downstream, slope and length
\code{shud.river}
\end{Description}
%
\begin{Usage}
\begin{verbatim}
shud.river(sl, dem, rivord = NULL, rivdown = NULL, AREA = NULL)
\end{verbatim}
\end{Usage}
%
\begin{Arguments}
\begin{ldescription}
\item[\code{sl}] SpatialLines*

\item[\code{dem}] Raster of elevation

\item[\code{rivord}] Order of each river reach

\item[\code{rivdown}] Downstream Index of each river reach.

\item[\code{AREA}] Area of the watershed, for estimating the width/depth of river.
\end{ldescription}
\end{Arguments}
%
\begin{Value}
SHUD.RIVER
\end{Value}
\inputencoding{utf8}
\HeaderA{shud.rivseg}{Generate the river segments table \code{shud.rivseg}}{shud.rivseg}
%
\begin{Description}\relax
Generate the river segments table
\code{shud.rivseg}
\end{Description}
%
\begin{Usage}
\begin{verbatim}
shud.rivseg(sl)
\end{verbatim}
\end{Usage}
%
\begin{Arguments}
\begin{ldescription}
\item[\code{sl}] SpatialLinesDataFrame
\end{ldescription}
\end{Arguments}
%
\begin{Value}
river segments table
\end{Value}
\inputencoding{utf8}
\HeaderA{shud.triangle}{Generate triangular mesh domain}{shud.triangle}
%
\begin{Description}\relax
Creates an unstructured triangular mesh using constrained Delaunay
triangulation. Modernized implementation using sf and terra internally,
but maintains backward compatibility with sp/raster objects.
\end{Description}
%
\begin{Usage}
\begin{verbatim}
shud.triangle(wb, dem = NULL, riv = NULL, hole = NULL, pts = NULL, q = 30, ...)

shud.triangle(wb, dem = NULL, riv = NULL, hole = NULL, pts = NULL, q = 30, ...)
\end{verbatim}
\end{Usage}
%
\begin{Arguments}
\begin{ldescription}
\item[\code{wb}] SpatialPolygon or SpatialLines which define the watershed boundary

\item[\code{dem}] Elevation data.

\item[\code{riv}] SpatialLines of river network, optional

\item[\code{hole}] Holes/lakes - sf object with POLYGON geometry, or legacy
sp object (auto-converted to sf) (optional)

\item[\code{pts}] Extra pts to build triangular mesh.

\item[\code{q}] minimum angle of triangle

\item[\code{...}] more options in RTriangle::triangulate()

\item[\code{lake}] SpatialPolygon of lake.
\end{ldescription}
\end{Arguments}
%
\begin{Details}\relax
This function has been modernized to use sf and terra internally for
better performance and maintainability. However, it automatically converts
legacy sp and raster objects, so existing code will continue to work.

The function uses RTriangle for constrained Delaunay triangulation with
quality constraints. The minimum angle parameter (q) controls triangle
quality - higher values produce better-shaped triangles but may increase
computation time. Values above 35 degrees may cause issues.
\end{Details}
%
\begin{Value}
A triangulation object from RTriangle with components:
\begin{ldescription}
\item[\code{P}] Matrix of point coordinates (n x 2)
\item[\code{T}] Matrix of triangle vertex indices (m x 3)
\item[\code{E}] Matrix of edge indices
\item[\code{NB}] Matrix of neighbor triangle indices

\end{ldescription}
A object with class triangulation.
\end{Value}
%
\begin{Examples}
\begin{ExampleCode}
## Not run: 
library(sf)

# Create simple boundary
boundary <- sf::st_as_sf(
  sf::st_sfc(sf::st_polygon(list(matrix(c(0,0, 10,0, 10,10, 0,10, 0,0),
                                         ncol=2, byrow=TRUE))))
)

# Generate mesh
tri <- shud.triangle(wb = boundary, q = 30, a = 2)
plot(tri, asp = 1)

# With rivers
rivers <- sf::st_as_sf(
  sf::st_sfc(sf::st_linestring(matrix(c(2,2, 8,8), ncol=2, byrow=TRUE)))
)
tri <- shud.triangle(wb = boundary, riv = rivers, q = 30)

## End(Not run)
\end{ExampleCode}
\end{Examples}
\inputencoding{utf8}
\HeaderA{shud\_river\_to\_sf}{Convert SHUD.RIVER to sf Object}{shud.Rul.river.Rul.to.Rul.sf}
%
\begin{Description}\relax
Extracts the sf network from a SHUD.RIVER object, or reconstructs it
from the legacy data.frame format if needed.
\end{Description}
%
\begin{Usage}
\begin{verbatim}
shud_river_to_sf(shud_river)
\end{verbatim}
\end{Usage}
%
\begin{Arguments}
\begin{ldescription}
\item[\code{shud\_river}] SHUD.RIVER S4 object
\end{ldescription}
\end{Arguments}
%
\begin{Details}\relax
If the SHUD.RIVER object has a network slot (modern format), it is
returned directly. Otherwise, the function attempts to reconstruct
an sf object from the point coordinates.
\end{Details}
%
\begin{Value}
sf object with river network
\end{Value}
%
\begin{Examples}
\begin{ExampleCode}
## Not run: 
# Load legacy SHUD.RIVER object
shud_river <- readriv("model.riv")
# Convert to sf
rivers_sf <- shud_river_to_sf(shud_river)

## End(Not run)
\end{ExampleCode}
\end{Examples}
\inputencoding{utf8}
\HeaderA{SimpleSpatial}{Simplify SpatialData.}{SimpleSpatial}
%
\begin{Description}\relax
Simplify SpatialData.
\end{Description}
%
\begin{Usage}
\begin{verbatim}
SimpleSpatial(x)
\end{verbatim}
\end{Usage}
%
\begin{Arguments}
\begin{ldescription}
\item[\code{x}] SpatialData
\end{ldescription}
\end{Arguments}
%
\begin{Value}
Simplified SpatialData
\end{Value}
\inputencoding{utf8}
\HeaderA{SimplifybyLen}{Simplify Polylines by length \code{SimplifybyLen}}{SimplifybyLen}
%
\begin{Description}\relax
Simplify Polylines by length
\code{SimplifybyLen}
\end{Description}
%
\begin{Usage}
\begin{verbatim}
SimplifybyLen(sp, split_length = 20, plot.results = FALSE)
\end{verbatim}
\end{Usage}
%
\begin{Arguments}
\begin{ldescription}
\item[\code{sp}] a SpatialLines or SpatialPoygons object

\item[\code{split\_length}] the length of the segments to split the lines into, in units of the SpatialLine object

\item[\code{plot.results}] TRUE/FALSE
\end{ldescription}
\end{Arguments}
%
\begin{Value}
simplified SpatialLines
\end{Value}
%
\begin{Source}\relax
https://stackoverflow.com/questions/38700246/how-do-i-split-divide-polyline-shapefiles-into-equally-length-smaller-segments
\end{Source}
\inputencoding{utf8}
\HeaderA{SinglePolygon}{Conver the MULTIPOLYGONS to SINGLEPOLYGONS.}{SinglePolygon}
%
\begin{Description}\relax
Conver the MULTIPOLYGONS to SINGLEPOLYGONS.
\end{Description}
%
\begin{Usage}
\begin{verbatim}
SinglePolygon(x, id = 0)
\end{verbatim}
\end{Usage}
%
\begin{Arguments}
\begin{ldescription}
\item[\code{x}] the spatialpolygon*

\item[\code{id}] Index of the sorted (decreasing) polygons to return. default = 0;
\end{ldescription}
\end{Arguments}
\inputencoding{utf8}
\HeaderA{SlopeSatVaporPressure}{Slope of saturated vapor pressure, Eq 4.2.3 in David R Maidment, Handbook of Hydrology \code{SlopeSatVaporPressure}}{SlopeSatVaporPressure}
%
\begin{Description}\relax
Slope of saturated vapor pressure, Eq 4.2.3 in David R Maidment, Handbook of Hydrology
\code{SlopeSatVaporPressure}
\end{Description}
%
\begin{Usage}
\begin{verbatim}
SlopeSatVaporPressure(Temp, ES)
\end{verbatim}
\end{Usage}
%
\begin{Arguments}
\begin{ldescription}
\item[\code{Temp}] Tempearture \LinkA{C}{C}

\item[\code{ES}] 
\end{ldescription}
\end{Arguments}
\inputencoding{utf8}
\HeaderA{sp.CutSptialLines}{Cut sptialLines with threshold. \code{sp.CutSptialLines}}{sp.CutSptialLines}
%
\begin{Description}\relax
Cut sptialLines with threshold.
\code{sp.CutSptialLines}
\end{Description}
%
\begin{Usage}
\begin{verbatim}
sp.CutSptialLines(sl, tol)
\end{verbatim}
\end{Usage}
%
\begin{Arguments}
\begin{ldescription}
\item[\code{sl}] SpatialLines or SpatialLineDataFrame

\item[\code{tol}] Tolerence. If the length of segment is larger than tolerance, cut the segment until the maximum segment is shorter than tolerance.
\end{ldescription}
\end{Arguments}
%
\begin{Examples}
\begin{ExampleCode}
library(rSHUD)
library(sp)
x = 1:1000 / 100
l1 = Lines(Line(cbind(x, sin(x)) ), ID = 'a' )
sl = SpatialLines( list(l1) )
tol1 = 5;
tol2 = 2
sl1 = sp.CutSptialLines(sl, tol1)
sl2 = sp.CutSptialLines(sl, tol2)
par(mfrow = c(1, 2))
plot(sl1, col = 1:length(sl1)); title(paste0('Tol=', tol1))
plot(sl2, col = 1:length(sl2)); title(paste0('Tol=', tol2))

data(sh)
riv = sh$riv
x = sp.CutSptialLines(riv, tol = 5)
par(mfrow = c(2, 1))
plot(riv, col = 1:length(riv), lwd = 3);

plot(riv, col = 'gray', lwd = 3);
plot(add = TRUE, x, col = 1:length(x))
\end{ExampleCode}
\end{Examples}
\inputencoding{utf8}
\HeaderA{sp.mesh2Shape}{Convert mesh triangulation to sf shapefile}{sp.mesh2Shape}
%
\begin{Description}\relax
Converts a triangular mesh to an sf object with POLYGON geometry.
Modernized implementation using sf, with automatic sp conversion.
\end{Description}
%
\begin{Usage}
\begin{verbatim}
sp.mesh2Shape(pm = readmesh(), dbf = NULL, crs = NULL)

sp.mesh2Shape(pm = readmesh(), dbf = NULL, crs = NULL)
\end{verbatim}
\end{Usage}
%
\begin{Arguments}
\begin{ldescription}
\item[\code{pm}] SHUD mesh object

\item[\code{dbf}] attribute table of the mesh triangles.

\item[\code{crs}] Projection parameters.
\end{ldescription}
\end{Arguments}
%
\begin{Details}\relax
This function has been modernized to return sf objects instead of sp
SpatialPolygonsDataFrame. The sf format is more efficient and integrates
better with modern R spatial packages.

Area is automatically calculated using sf::st\_area() and added to the
attributes table.
\end{Details}
%
\begin{Value}
sf object with POLYGON geometry and attributes

SpatialPolygons object
\end{Value}
%
\begin{Examples}
\begin{ExampleCode}
## Not run: 
library(sf)

# Create simple mesh
boundary <- sf::st_as_sf(
  sf::st_sfc(sf::st_polygon(list(matrix(c(0,0, 10,0, 10,10, 0,10, 0,0),
                                         ncol=2, byrow=TRUE))))
)
tri <- shud.triangle(wb = boundary, q = 30, a = 2)

# Convert to sf
mesh_sf <- sp.mesh2Shape(pm = tri, crs = 4326)
plot(mesh_sf)

## End(Not run)
\end{ExampleCode}
\end{Examples}
\inputencoding{utf8}
\HeaderA{sp.riv2shp}{Convert the SHUD.RIVER class to SpatialLines \code{sp.riv2shp}}{sp.riv2shp}
%
\begin{Description}\relax
Convert the SHUD.RIVER class to SpatialLines
\code{sp.riv2shp}
\end{Description}
%
\begin{Usage}
\begin{verbatim}
sp.riv2shp(pr = readriv(), dbf = NULL)
\end{verbatim}
\end{Usage}
%
\begin{Arguments}
\begin{ldescription}
\item[\code{pr}] SHUD.RIVER class

\item[\code{dbf}] Attribute data for exported SpatialLines
\end{ldescription}
\end{Arguments}
%
\begin{Value}
SpatialLinesDataFrame
\end{Value}
\inputencoding{utf8}
\HeaderA{sp.RiverDown}{determine index of downstream \code{sp.RiverOrder}}{sp.RiverDown}
%
\begin{Description}\relax
determine index of downstream
\code{sp.RiverOrder}
\end{Description}
%
\begin{Usage}
\begin{verbatim}
sp.RiverDown(sp, coord = extractCoords(sp))
\end{verbatim}
\end{Usage}
%
\begin{Arguments}
\begin{ldescription}
\item[\code{sp}] SpatialLines*

\item[\code{coord}] coordinates of \code{sp}.
\end{ldescription}
\end{Arguments}
%
\begin{Value}
Index of downstream for each segements
\end{Value}
\inputencoding{utf8}
\HeaderA{sp.RiverOrder}{calculate river order \code{sp.RiverOrder}}{sp.RiverOrder}
%
\begin{Description}\relax
calculate river order
\code{sp.RiverOrder}
\end{Description}
%
\begin{Usage}
\begin{verbatim}
sp.RiverOrder(sp, coord = extractCoords(sp))
\end{verbatim}
\end{Usage}
%
\begin{Arguments}
\begin{ldescription}
\item[\code{sp}] SpatialLines*

\item[\code{coord}] coordinates of \code{sp}.
\end{ldescription}
\end{Arguments}
%
\begin{Value}
Stream Order of SpatialLines*
\end{Value}
%
\begin{Examples}
\begin{ExampleCode}
library(rgdal)
data(sac)
riv=sac$riv
ord = sp.RiverOrder(riv)
print(ord)
idx = sort(unique(ord))
raster::plot(riv, col=ord)
legend('topleft', legend=idx, col=idx, lwd=1, lty=1)
\end{ExampleCode}
\end{Examples}
\inputencoding{utf8}
\HeaderA{sp.RiverPath}{determine River path, disolve the river network based on downstream-relationship. \code{sp.RiverPath}}{sp.RiverPath}
%
\begin{Description}\relax
determine River path, disolve the river network based on downstream-relationship.
\code{sp.RiverPath}
\end{Description}
%
\begin{Usage}
\begin{verbatim}
sp.RiverPath(
  sp,
  coord = extractCoords(sp, unique = TRUE),
  idown = sp.RiverDown(sp, coord = coord),
  tol.simplify = 0
)
\end{verbatim}
\end{Usage}
%
\begin{Arguments}
\begin{ldescription}
\item[\code{sp}] SpatialLines*

\item[\code{coord}] coordinates of \code{sp}.

\item[\code{idown}] downstream index of \code{sp}.

\item[\code{tol.simplify}] tolerance for simplification of river lines.
\end{ldescription}
\end{Arguments}
%
\begin{Value}
List of River Path(SegIDs, PointIDs, sp)
\end{Value}
\inputencoding{utf8}
\HeaderA{sp.RiverSeg}{Cut the river by shud.mesh \code{sp.RiverSeg}}{sp.RiverSeg}
%
\begin{Description}\relax
Cut the river by shud.mesh
\code{sp.RiverSeg}
\end{Description}
%
\begin{Usage}
\begin{verbatim}
sp.RiverSeg(sp.mesh, sp.riv)
\end{verbatim}
\end{Usage}
%
\begin{Arguments}
\begin{ldescription}
\item[\code{sp.mesh}] shud mesh

\item[\code{sp.riv}] river
\end{ldescription}
\end{Arguments}
%
\begin{Value}
SpatialLinesDataFrame
\end{Value}
\inputencoding{utf8}
\HeaderA{sp.simplifyLen}{Simplify Polylines by length \code{sp.simplifyLen}}{sp.simplifyLen}
%
\begin{Description}\relax
Simplify Polylines by length
\code{sp.simplifyLen}
\end{Description}
%
\begin{Usage}
\begin{verbatim}
sp.simplifyLen(sp, tol, threshold = tol/2)
\end{verbatim}
\end{Usage}
%
\begin{Arguments}
\begin{ldescription}
\item[\code{sp}] a SpatialLines or SpatialPoygons object

\item[\code{tol}] Tolerance

\item[\code{threshold}] Minimum length for a segment
\end{ldescription}
\end{Arguments}
\inputencoding{utf8}
\HeaderA{sp.Tri2Shape}{Convert triangulation to shapefile}{sp.Tri2Shape}
%
\begin{Description}\relax
Converts a triangular mesh to an sf object with POLYGON geometry.
Modernized implementation using sf.
\end{Description}
%
\begin{Usage}
\begin{verbatim}
sp.Tri2Shape(tri, dbf = NULL, crs = NA)

sp.Tri2Shape(tri, dbf = NULL, crs = NA)
\end{verbatim}
\end{Usage}
%
\begin{Arguments}
\begin{ldescription}
\item[\code{tri}] triangulate

\item[\code{dbf}] attribute table, data.frame or matrix

\item[\code{crs}] Projection
\end{ldescription}
\end{Arguments}
%
\begin{Value}
sf object with POLYGON geometry

Coordinates of triangles centroids
\end{Value}
\inputencoding{utf8}
\HeaderA{sp2PSLG}{Convert SpatialLines or SpatialPolygons a Planar Straight Line Graph object \code{sp2PSLG}}{sp2PSLG}
%
\begin{Description}\relax
Convert SpatialLines or SpatialPolygons a Planar Straight Line Graph object
\code{sp2PSLG}
\end{Description}
%
\begin{Usage}
\begin{verbatim}
sp2PSLG(sp)
\end{verbatim}
\end{Usage}
%
\begin{Arguments}
\begin{ldescription}
\item[\code{sp}] SpatialLines or SpatialPolygons
\end{ldescription}
\end{Arguments}
%
\begin{Value}
pslg class
\end{Value}
%
\begin{Examples}
\begin{ExampleCode}
library(rgeos)
sl = readWKT("MULTIPOLYGON(((1 1,5 1,5 5,1 5,1 1),(2 2,2 3,3 3,3 2,2 2)))")
x = sp2PSLG(sl)
\end{ExampleCode}
\end{Examples}
\inputencoding{utf8}
\HeaderA{sunapee}{This is data of Lake Sunappe to be included in my package}{sunapee}
\keyword{sunapee}{sunapee}
%
\begin{Description}\relax
This is data of Lake Sunappe to be included in my package
\end{Description}
\inputencoding{utf8}
\HeaderA{timeseries-io}{Time Series I/O Functions}{timeseries.Rdash.io}
%
\begin{Description}\relax
This module provides functions for reading and writing time series data
in SHUD format. All functions use modern snake\_case naming and provide
comprehensive error handling.
\end{Description}
\inputencoding{utf8}
\HeaderA{toCentroid}{Find the centroid of the triangles \code{toCentroid}}{toCentroid}
%
\begin{Description}\relax
Find the centroid of the triangles
\code{toCentroid}
\end{Description}
%
\begin{Usage}
\begin{verbatim}
toCentroid(tri, point)
\end{verbatim}
\end{Usage}
%
\begin{Arguments}
\begin{ldescription}
\item[\code{tri}] Triangle, N x 3 dimension

\item[\code{point}] Coordinates.
\end{ldescription}
\end{Arguments}
%
\begin{Value}
Coordinates of triangles centroids
\end{Value}
\inputencoding{utf8}
\HeaderA{Tri2Centroid}{Calculate triangle centroids}{Tri2Centroid}
%
\begin{Description}\relax
Computes the centroids of triangles in a triangulation.
\end{Description}
%
\begin{Usage}
\begin{verbatim}
Tri2Centroid(tri)

Tri2Centroid(tri)
\end{verbatim}
\end{Usage}
%
\begin{Arguments}
\begin{ldescription}
\item[\code{tri}] Triangles defination
\end{ldescription}
\end{Arguments}
%
\begin{Value}
Matrix of centroid coordinates (n x 2)

Centroids of the triangles, m x 2;
\end{Value}
%
\begin{Examples}
\begin{ExampleCode}
## Not run: 
library(sf)

boundary <- sf::st_as_sf(
  sf::st_sfc(sf::st_polygon(list(matrix(c(0,0, 10,0, 10,10, 0,10, 0,0),
                                         ncol=2, byrow=TRUE))))
)
tri <- shud.triangle(wb = boundary, q = 30, a = 2)
centroids <- Tri2Centroid(tri)

## End(Not run)
\end{ExampleCode}
\end{Examples}
\inputencoding{utf8}
\HeaderA{triTopology}{Determine the topological relation among triangles.}{triTopology}
%
\begin{Description}\relax
Determine the topological relation among triangles.
\end{Description}
%
\begin{Usage}
\begin{verbatim}
triTopology(tri)
\end{verbatim}
\end{Usage}
%
\begin{Arguments}
\begin{ldescription}
\item[\code{tri}] triangle defination, m x 3
\end{ldescription}
\end{Arguments}
%
\begin{Value}
topological relation of triangles. m x 4; cols = c(ID, Nabor1, Nabor2, Nabor3)
\end{Value}
\inputencoding{utf8}
\HeaderA{ts\_to\_daily}{Convert time series to daily aggregation}{ts.Rul.to.Rul.daily}
\keyword{timeseries-io}{ts\_to\_daily}
%
\begin{Description}\relax
Aggregates time series data to daily values using specified function.
\end{Description}
%
\begin{Usage}
\begin{verbatim}
ts_to_daily(x, FUN = mean)
\end{verbatim}
\end{Usage}
%
\begin{Arguments}
\begin{ldescription}
\item[\code{x}] xts object with time series data

\item[\code{FUN}] Function to use for aggregation (default: mean)
\end{ldescription}
\end{Arguments}
%
\begin{Value}
xts object with daily aggregated data
\end{Value}
%
\begin{SeeAlso}\relax
Other timeseries-io: 
\code{\LinkA{read\_lai}{read.Rul.lai}()},
\code{\LinkA{read\_tsd}{read.Rul.tsd}()},
\code{\LinkA{ts\_to\_df}{ts.Rul.to.Rul.df}()},
\code{\LinkA{write\_tsd}{write.Rul.tsd}()},
\code{\LinkA{write\_xts}{write.Rul.xts}()}
\end{SeeAlso}
%
\begin{Examples}
\begin{ExampleCode}
## Not run: 
# Create hourly data
hourly_data <- xts::xts(runif(24*7), 
                       order.by = seq(Sys.time(), by = "hour", length.out = 24*7))

# Convert to daily means
daily_data <- ts_to_daily(hourly_data)

## End(Not run)
\end{ExampleCode}
\end{Examples}
\inputencoding{utf8}
\HeaderA{ts\_to\_df}{Convert time series to data frame}{ts.Rul.to.Rul.df}
\keyword{timeseries-io}{ts\_to\_df}
%
\begin{Description}\relax
Converts xts time series to data frame with time column.
\end{Description}
%
\begin{Usage}
\begin{verbatim}
ts_to_df(x, time_col = "Time")
\end{verbatim}
\end{Usage}
%
\begin{Arguments}
\begin{ldescription}
\item[\code{x}] xts object with time series data

\item[\code{time\_col}] Character. Name for time column (default: "Time")
\end{ldescription}
\end{Arguments}
%
\begin{Value}
data.frame with time and data columns
\end{Value}
%
\begin{SeeAlso}\relax
Other timeseries-io: 
\code{\LinkA{read\_lai}{read.Rul.lai}()},
\code{\LinkA{read\_tsd}{read.Rul.tsd}()},
\code{\LinkA{ts\_to\_daily}{ts.Rul.to.Rul.daily}()},
\code{\LinkA{write\_tsd}{write.Rul.tsd}()},
\code{\LinkA{write\_xts}{write.Rul.xts}()}
\end{SeeAlso}
%
\begin{Examples}
\begin{ExampleCode}
## Not run: 
dates <- seq(as.Date("2020-01-01"), as.Date("2020-01-10"), by = "day")
ts_data <- xts::xts(runif(length(dates)), order.by = dates)
df_data <- ts_to_df(ts_data)

## End(Not run)
\end{ExampleCode}
\end{Examples}
\inputencoding{utf8}
\HeaderA{ts2Daily}{Convert the time-series data to daily data.}{ts2Daily}
%
\begin{Description}\relax
Convert the time-series data to daily data.
\end{Description}
%
\begin{Usage}
\begin{verbatim}
ts2Daily(x, FUN = mean, ...)
\end{verbatim}
\end{Usage}
%
\begin{Arguments}
\begin{ldescription}
\item[\code{x}] xts,zoo data.

\item[\code{FUN}] function.

\item[\code{...}] more options in xts::apply.daily()
\end{ldescription}
\end{Arguments}
%
\begin{Value}
xts data whose time index is daily.
\end{Value}
%
\begin{Examples}
\begin{ExampleCode}
library(xts)
x=as.xts(rnorm(100), order.by = as.POSIXct('2000-01-01')+ 1:100 * 3600*24)
xd = ts2Daily(x)
class(time(x))
class(time(xd))
\end{ExampleCode}
\end{Examples}
\inputencoding{utf8}
\HeaderA{ts2df}{Convert Time-Series data to data.frame \code{ts2df}}{ts2df}
%
\begin{Description}\relax
Convert Time-Series data to data.frame
\code{ts2df}
\end{Description}
%
\begin{Usage}
\begin{verbatim}
ts2df(x)
\end{verbatim}
\end{Usage}
%
\begin{Arguments}
\begin{ldescription}
\item[\code{x}] Time-Series data.
\end{ldescription}
\end{Arguments}
%
\begin{Value}
data.frame. First column is Time.
\end{Value}
\inputencoding{utf8}
\HeaderA{ts2map}{Plot map of time-series data \code{ts2map}}{ts2map}
%
\begin{Description}\relax
Plot map of time-series data
\code{ts2map}
\end{Description}
%
\begin{Usage}
\begin{verbatim}
ts2map(x, fun = mean, raster = TRUE, ...)
\end{verbatim}
\end{Usage}
%
\begin{Arguments}
\begin{ldescription}
\item[\code{x}] Time-series data

\item[\code{fun}] Functions to summarize the TS data.

\item[\code{raster}] Whether to make output as Raster

\item[\code{...}] More options in fun
\end{ldescription}
\end{Arguments}
\inputencoding{utf8}
\HeaderA{tsLaiRf}{Generate the default LAI and Roughness Length for UMD classification system \code{tsLaiRf}}{tsLaiRf}
%
\begin{Description}\relax
Generate the default LAI and Roughness Length for UMD classification system
\code{tsLaiRf}
\end{Description}
%
\begin{Usage}
\begin{verbatim}
tsLaiRf(type = 1)
\end{verbatim}
\end{Usage}
%
\begin{Arguments}
\begin{ldescription}
\item[\code{type}] 1: LAI, 2:RL
\end{ldescription}
\end{Arguments}
%
\begin{Value}
Default LAI and RL
\end{Value}
%
\begin{Examples}
\begin{ExampleCode}
tsLaiRf(type = 1)
tsLaiRf(type = 2)
\end{ExampleCode}
\end{Examples}
\inputencoding{utf8}
\HeaderA{Untructure Domain-class}{Class of SHUD.MESH}{Untructure Domain.Rdash.class}
\aliasA{SHUD.MESH}{Untructure Domain-class}{SHUD.MESH}
%
\begin{Description}\relax
S4 class representing an unstructured triangular mesh domain for SHUD models.
The class stores mesh topology, point coordinates, and associated attributes.
\end{Description}
%
\begin{Value}
Class of SHUD.MESH
\end{Value}
%
\begin{Section}{Slots}

\begin{description}

\item[\code{mesh}] data.frame with mesh topology (triangles and neighbors)

\item[\code{point}] data.frame with point coordinates and attributes

\end{description}
\end{Section}
%
\begin{Section}{Mesh Data Frame}

The mesh slot contains:
\begin{itemize}

\item{} ID: Triangle ID
\item{} Node1, Node2, Node3: Vertex indices
\item{} Nabr1, Nabr2, Nabr3: Neighbor triangle indices
\item{} Zmax: Maximum elevation

\end{itemize}

\end{Section}
%
\begin{Section}{Point Data Frame}

The point slot contains:
\begin{itemize}

\item{} ID: Point ID
\item{} X, Y: Coordinates
\item{} AqDepth: Aquifer depth
\item{} Elevation: Surface elevation

\end{itemize}

\end{Section}
\inputencoding{utf8}
\HeaderA{VaporPressure\_Act}{Actual vapor pressure \code{VaporPressure\_Act}}{VaporPressure.Rul.Act}
%
\begin{Description}\relax
Actual vapor pressure
\code{VaporPressure\_Act}
\end{Description}
%
\begin{Usage}
\begin{verbatim}
VaporPressure_Act(esat, rh)
\end{verbatim}
\end{Usage}
%
\begin{Arguments}
\begin{ldescription}
\item[\code{esat}] Saturated vapor pressure

\item[\code{rh}] Relative Humidity \LinkA{-}{.Rdash.} 0\textasciitilde{}1
\end{ldescription}
\end{Arguments}
\inputencoding{utf8}
\HeaderA{VaporPressure\_Sat}{Saturated vapor pressure, Eq 4.2.2 in David R Maidment, Handbook of Hydrology \code{VaporPressure\_Sat}}{VaporPressure.Rul.Sat}
%
\begin{Description}\relax
Saturated vapor pressure, Eq 4.2.2 in David R Maidment, Handbook of Hydrology
\code{VaporPressure\_Sat}
\end{Description}
%
\begin{Usage}
\begin{verbatim}
VaporPressure_Sat(Temp)
\end{verbatim}
\end{Usage}
%
\begin{Arguments}
\begin{ldescription}
\item[\code{Temp}] Air temperature, \LinkA{C}{C}
\end{ldescription}
\end{Arguments}
%
\begin{Value}
Saturated vapor pressure \LinkA{kPa}{kPa}
\end{Value}
\inputencoding{utf8}
\HeaderA{vector\_to\_raster}{Convert vector data to raster}{vector.Rul.to.Rul.raster}
%
\begin{Description}\relax
Converts vector spatial data (sf or SpatVector) to a regular raster grid.
This function uses terra for rasterization and supports automatic or manual
resolution settings and multiple aggregation functions.
\end{Description}
%
\begin{Usage}
\begin{verbatim}
vector_to_raster(
  vector,
  template = NULL,
  field = 1,
  resolution = NULL,
  ngrids = 200,
  fun = "last",
  background = NA
)
\end{verbatim}
\end{Usage}
%
\begin{Arguments}
\begin{ldescription}
\item[\code{vector}] sf or SpatVector object to rasterize.

\item[\code{template}] SpatRaster template defining output grid (optional).
If NULL, creates template from vector extent.

\item[\code{field}] Character or numeric, field name or index to rasterize.
Default is 1 (first attribute field). Can also be a numeric vector
of values matching the number of features.

\item[\code{resolution}] Numeric, output raster resolution in map units (optional).
Only used if template is NULL. If NULL and template is NULL,
resolution is automatically calculated from extent.

\item[\code{ngrids}] Integer, number of grid cells along x direction (optional).
Only used if both template and resolution are NULL. Default is 200.

\item[\code{fun}] Character, aggregation function for multiple values per cell.
Options: "first", "last", "min", "max", "mean", "sum", "count".
Default is "last".

\item[\code{background}] Numeric, value for cells with no data. Default is NA.
\end{ldescription}
\end{Arguments}
%
\begin{Value}
SpatRaster object with rasterized vector data.
\end{Value}
%
\begin{Section}{Migration Note}

This function replaces \code{sp2raster()} and uses terra/sf instead of
raster/sp. Key differences:
\begin{itemize}

\item{} Returns SpatRaster instead of RasterLayer
\item{} Accepts sf or SpatVector instead of Spatial* objects
\item{} Does not use global MASK environment variable
\item{} Parameter 'mask' renamed to 'template'
\item{} More flexible field specification

\end{itemize}

\end{Section}
%
\begin{Examples}
\begin{ExampleCode}
## Not run: 
# Basic usage with sf object
library(sf)
polygons <- st_read("watershed.shp")
r <- vector_to_raster(polygons, field = "elevation")
terra::plot(r)

# With custom resolution
r <- vector_to_raster(polygons, field = "landcover", resolution = 100)

# With template raster
template <- terra::rast(extent = terra::ext(polygons), resolution = 50)
r <- vector_to_raster(polygons, template = template, field = "soiltype")

# Using aggregation function
points <- st_read("stations.shp")
r <- vector_to_raster(points, field = "rainfall", fun = "mean")

## End(Not run)
\end{ExampleCode}
\end{Examples}
\inputencoding{utf8}
\HeaderA{voronoipolygons}{Generate Thiesson Polygons from a point data \code{voronoipolygons}}{voronoipolygons}
%
\begin{Description}\relax
Generate Thiesson Polygons from a point data
\code{voronoipolygons}
\end{Description}
%
\begin{Usage}
\begin{verbatim}
voronoipolygons(x, pts = x@coords, rw = NULL, crs = NULL)
\end{verbatim}
\end{Usage}
%
\begin{Arguments}
\begin{ldescription}
\item[\code{x}] ShapePoints* in PCS

\item[\code{pts}] Coordinates (x,y) of the points

\item[\code{rw}] extent

\item[\code{crs}] projection parameters
\end{ldescription}
\end{Arguments}
%
\begin{Value}
ShapePolygon*
\end{Value}
%
\begin{Examples}
\begin{ExampleCode}
library(rgeos)
library(rSHUD)
n = 10
xx = rnorm(n)
yy = rnorm(n)
str = paste('MULTIPOINT(', paste(paste('(', xx, yy, ')'), collapse = ','), ')')
x = readWKT(str)
vx = voronoipolygons(x)
raster::plot(vx, axes = TRUE)
raster::plot(add = TRUE, x, col = 2)
#' ====END====

x = 1:5
y = 1:5
xy = expand.grid(x, y)
vx = voronoipolygons(pts = xy)
plot(vx, axes = TRUE)
points(xy)
#' ====END====

library(rgdal)
library(raster)
library(rgeos)
n = 10
xx = rnorm(n)
yy = rnorm(n)
str = paste('MULTIPOINT(', paste(paste('(', xx, yy, ')'), collapse = ','), ')')
x = readWKT(str)
y = readWKT(paste('MULTIPOINT(', paste(paste('(', xx + 2, yy + 2, ')'), collapse = ','), ')'))
e1 = extent(y)
e2 = extent(x)
rw = c(min(e1[1], e2[1]),
       max(e1[2], e2[2]),
       min(e1[3], e2[3]),
       max(e1[4], e2[4]) ) + c(-1, 1, -1, 1)
vx = voronoipolygons(x = x, rw = rw)
plot(vx); plot(add = TRUE, x, col = 2); plot(add = TRUE, y, col = 3)
\end{ExampleCode}
\end{Examples}
\inputencoding{utf8}
\HeaderA{watershedDelineation}{Do Watershed Delineation \code{watershedDelineation}}{watershedDelineation}
%
\begin{Description}\relax
Do Watershed Delineation
\code{watershedDelineation}
\end{Description}
%
\begin{Usage}
\begin{verbatim}
watershedDelineation(
  fn.wbd,
  fnr.dem,
  fsp.outlets = NULL,
  dirout = tempdir(),
  dirtemp = tempdir(),
  FlowAccCell.min = NULL,
  plot = FALSE
)
\end{verbatim}
\end{Usage}
%
\begin{Arguments}
\begin{ldescription}
\item[\code{fn.wbd}] file of watershed SpatialPolygons

\item[\code{fnr.dem}] file of DEM in raster

\item[\code{fsp.outlets}] file of watershed outlet

\item[\code{dirout}] dir for out data(dem\_filled.tif, dem\_stm.shp, dem\_wbd.shp, dem\_outlets.shp) saved. Default = tempdir()

\item[\code{dirtemp}] dir for temporary data save. Default = tempdir()

\item[\code{FlowAccCell.min}] Minimum number of cells for river generation. Default  = Number of DEM / 100
\end{ldescription}
\end{Arguments}
%
\begin{Value}
list of path to dem-filled, wbd, riv, outlet
\end{Value}
\inputencoding{utf8}
\HeaderA{wb.all}{Calculate the water balance \code{wb.all}}{wb.all}
%
\begin{Description}\relax
Calculate the water balance
\code{wb.all}
\end{Description}
%
\begin{Usage}
\begin{verbatim}
wb.all(
  xl = loaddata(varname = c(paste0("elev", c("prcp", "eta", "etp")), "rivqdown",
    paste0("eley", c("surf", "unsat", "gw")))),
  ic = readic(),
  fun = xts::apply.monthly,
  plot = TRUE
)
\end{verbatim}
\end{Usage}
%
\begin{Arguments}
\begin{ldescription}
\item[\code{xl}] List of data. Five variables are included: prcp (eleqprcp), etic (eleqetic), ettr (eleqettr), etev (eleqetev) and discharge (rivqflx)

\item[\code{ic}] Initial Condition.

\item[\code{fun}] function to process the time-series data. Default = apply.daily.

\item[\code{plot}] Whether plot the result
\end{ldescription}
\end{Arguments}
%
\begin{Value}
A matrix, contains the colums of water balance factors
\end{Value}
\inputencoding{utf8}
\HeaderA{wb.DS}{Calculate the Change of Storage. \code{wb.DS}}{wb.DS}
%
\begin{Description}\relax
Calculate the Change of Storage.
\code{wb.DS}
\end{Description}
%
\begin{Usage}
\begin{verbatim}
wb.DS(
  xl = loaddata(varname = c(paste0("eley", c("surf", "unsat", "gw")), paste0("rivy",
    "stage"))),
  ic = readic()
)
\end{verbatim}
\end{Usage}
%
\begin{Arguments}
\begin{ldescription}
\item[\code{xl}] List of data. Five variables are included: surface (eleysurf), unsat zone (eleyunsat), groundwater (eleygw) and river stage (rivystage)

\item[\code{ic}] Initial Condition.
\end{ldescription}
\end{Arguments}
%
\begin{Value}
A list. list(Ele, Riv)
\end{Value}
\inputencoding{utf8}
\HeaderA{wb.ele}{Calculate the water balance \code{wb.ele}}{wb.ele}
%
\begin{Description}\relax
Calculate the water balance
\code{wb.ele}
\end{Description}
%
\begin{Usage}
\begin{verbatim}
wb.ele(
  xl = loaddata(varname = c(paste0("elev", c("prcp", "etic", "ettr", "etev")),
    paste0("eleq", c("surf", "sub")), paste0("eley", c("surf", "unsat", "gw")))),
  fun = NULL,
  period = "years",
  plot = TRUE
)
\end{verbatim}
\end{Usage}
%
\begin{Arguments}
\begin{ldescription}
\item[\code{xl}] List of data. Five variables are included:

\item[\code{fun}] function to process the time-series data. Default = NULL

\item[\code{period}] Period of the waterbalance

\item[\code{plot}] Whether plot the result
\end{ldescription}
\end{Arguments}
%
\begin{Value}
A matrix, contains the colums of water balance factors
\end{Value}
\inputencoding{utf8}
\HeaderA{wb.lake}{Calculate the Change of Storage. \code{wb.lake}}{wb.lake}
%
\begin{Description}\relax
Calculate the Change of Storage.
\code{wb.lake}
\end{Description}
%
\begin{Usage}
\begin{verbatim}
wb.lake(
  xl = loaddata(varname = paste0("lak", c("ystage", "atop", "qrivin", "qrivout", "qsub",
    "qsurf", "vevap", "vprcp"))),
  lakeid = NULL
)
\end{verbatim}
\end{Usage}
%
\begin{Arguments}
\begin{ldescription}
\item[\code{xl}] List of data.

\item[\code{lakeid}] Index of the lake for waterbalance calculation.
\end{ldescription}
\end{Arguments}
%
\begin{Value}
A list of water balance by each lake; a matrix in each list-item
\end{Value}
\inputencoding{utf8}
\HeaderA{wb.riv}{Calculate the water balance \code{wb.riv}}{wb.riv}
%
\begin{Description}\relax
Calculate the water balance
\code{wb.riv}
\end{Description}
%
\begin{Usage}
\begin{verbatim}
wb.riv(
  xl = loaddata(varname = paste0("rivq", c("sub", "surf", "down"))),
  fun = NULL,
  plot = TRUE
)
\end{verbatim}
\end{Usage}
%
\begin{Arguments}
\begin{ldescription}
\item[\code{xl}] List of data. Five variables are included: prcp (eleqprcp),and discharge ()

\item[\code{fun}] function to process the time-series data. Default = NULL

\item[\code{plot}] Whether plot the result
\end{ldescription}
\end{Arguments}
%
\begin{Value}
A matrix, contains the colums of water balance factors
\end{Value}
\inputencoding{utf8}
\HeaderA{which\_outliers}{Find the outliers in a dataset \code{which\_outliers}}{which.Rul.outliers}
%
\begin{Description}\relax
Find the outliers in a dataset
\code{which\_outliers}
\end{Description}
%
\begin{Usage}
\begin{verbatim}
which_outliers(x, na.rm = TRUE, probs = c(0.5, 0.95), ...)
\end{verbatim}
\end{Usage}
%
\begin{Arguments}
\begin{ldescription}
\item[\code{x}] dataset

\item[\code{na.rm}] Whether ignore the na value.

\item[\code{probs}] Range of non-outlier range

\item[\code{...}] More options in quatile();
\end{ldescription}
\end{Arguments}
%
\begin{Value}
Index of the outliers
\end{Value}
%
\begin{Examples}
\begin{ExampleCode}
set.seed(1)
x <- rnorm(100)
x <- c(-10, x, 10)
id <- which_outliers(x, probs=c(0.05,0.95))
y = x
y[id]=NA
par(mfrow = c(1, 2))
boxplot(x, ylim=range(x))
boxplot(y, ylim=range(x))
par(mfrow=c(1,1))
\end{ExampleCode}
\end{Examples}
\inputencoding{utf8}
\HeaderA{WindProfile}{Calculate windspeed at Zx heights. \code{WindProfile}}{WindProfile}
%
\begin{Description}\relax
Calculate windspeed at Zx heights.
\code{WindProfile}
\end{Description}
%
\begin{Usage}
\begin{verbatim}
WindProfile(zx, Um, zm = 10, d = 0.12, roughness = 0.012)
\end{verbatim}
\end{Usage}
%
\begin{Arguments}
\begin{ldescription}
\item[\code{zx}] Desired height \LinkA{m}{m}

\item[\code{Um}] Windspeed at measured height \LinkA{m s-1}{m s.Rdash.1}

\item[\code{zm}] Measured height \LinkA{m s-1}{m s.Rdash.1}

\item[\code{d}] Zero plane \LinkA{m}{m}

\item[\code{roughness}] roughness \LinkA{m}{m}
\end{ldescription}
\end{Arguments}
%
\begin{Value}
windspeed at any desired heights.
\end{Value}
\inputencoding{utf8}
\HeaderA{write\_config}{Write SHUD configuration file}{write.Rul.config}
\keyword{shud-io}{write\_config}
%
\begin{Description}\relax
Writes SHUD model configuration parameters to a file (.para, .calib, etc.).
\end{Description}
%
\begin{Usage}
\begin{verbatim}
write_config(x, file)
\end{verbatim}
\end{Usage}
%
\begin{Arguments}
\begin{ldescription}
\item[\code{x}] SHUD model configuration parameters or calibration data

\item[\code{file}] Character. Full path to the output configuration file
\end{ldescription}
\end{Arguments}
%
\begin{SeeAlso}\relax
Other shud-io: 
\code{\LinkA{read\_att}{read.Rul.att}()},
\code{\LinkA{read\_calib}{read.Rul.calib}()},
\code{\LinkA{read\_config}{read.Rul.config}()},
\code{\LinkA{read\_df}{read.Rul.df}()},
\code{\LinkA{read\_forc\_csv}{read.Rul.forc.Rul.csv}()},
\code{\LinkA{read\_forc\_fn}{read.Rul.forc.Rul.fn}()},
\code{\LinkA{read\_geol}{read.Rul.geol}()},
\code{\LinkA{read\_ic}{read.Rul.ic}()},
\code{\LinkA{read\_lc}{read.Rul.lc}()},
\code{\LinkA{read\_mesh}{read.Rul.mesh}()},
\code{\LinkA{read\_para}{read.Rul.para}()},
\code{\LinkA{read\_river\_sp}{read.Rul.river.Rul.sp}()},
\code{\LinkA{read\_river}{read.Rul.river}()},
\code{\LinkA{read\_rivseg}{read.Rul.rivseg}()},
\code{\LinkA{read\_soil}{read.Rul.soil}()},
\code{\LinkA{write\_df}{write.Rul.df}()},
\code{\LinkA{write\_forc}{write.Rul.forc}()},
\code{\LinkA{write\_ic}{write.Rul.ic}()},
\code{\LinkA{write\_mesh}{write.Rul.mesh}()},
\code{\LinkA{write\_river}{write.Rul.river}()}
\end{SeeAlso}
%
\begin{Examples}
\begin{ExampleCode}
## Not run: 
write_config(config_data, "model.para")

## End(Not run)
\end{ExampleCode}
\end{Examples}
\inputencoding{utf8}
\HeaderA{write\_df}{Write data frame to file with header}{write.Rul.df}
\keyword{shud-io}{write\_df}
%
\begin{Description}\relax
Writes a data frame to file with (nrow, ncol) header.
This is a core utility function used by other SHUD file writers.
\end{Description}
%
\begin{Usage}
\begin{verbatim}
write_df(x, file, append = FALSE, quiet = FALSE, header = NULL)
\end{verbatim}
\end{Usage}
%
\begin{Arguments}
\begin{ldescription}
\item[\code{x}] Data frame or matrix to write

\item[\code{file}] Character. Full path to output file

\item[\code{append}] Logical. Whether to append to existing file (default: FALSE)

\item[\code{quiet}] Logical. Whether to suppress messages (default: FALSE)

\item[\code{header}] Numeric vector. Custom header (default: c(nrow, ncol))
\end{ldescription}
\end{Arguments}
%
\begin{SeeAlso}\relax
Other shud-io: 
\code{\LinkA{read\_att}{read.Rul.att}()},
\code{\LinkA{read\_calib}{read.Rul.calib}()},
\code{\LinkA{read\_config}{read.Rul.config}()},
\code{\LinkA{read\_df}{read.Rul.df}()},
\code{\LinkA{read\_forc\_csv}{read.Rul.forc.Rul.csv}()},
\code{\LinkA{read\_forc\_fn}{read.Rul.forc.Rul.fn}()},
\code{\LinkA{read\_geol}{read.Rul.geol}()},
\code{\LinkA{read\_ic}{read.Rul.ic}()},
\code{\LinkA{read\_lc}{read.Rul.lc}()},
\code{\LinkA{read\_mesh}{read.Rul.mesh}()},
\code{\LinkA{read\_para}{read.Rul.para}()},
\code{\LinkA{read\_river\_sp}{read.Rul.river.Rul.sp}()},
\code{\LinkA{read\_river}{read.Rul.river}()},
\code{\LinkA{read\_rivseg}{read.Rul.rivseg}()},
\code{\LinkA{read\_soil}{read.Rul.soil}()},
\code{\LinkA{write\_config}{write.Rul.config}()},
\code{\LinkA{write\_forc}{write.Rul.forc}()},
\code{\LinkA{write\_ic}{write.Rul.ic}()},
\code{\LinkA{write\_mesh}{write.Rul.mesh}()},
\code{\LinkA{write\_river}{write.Rul.river}()}
\end{SeeAlso}
%
\begin{Examples}
\begin{ExampleCode}
## Not run: 
write_df(data_matrix, "output.txt")

## End(Not run)
\end{ExampleCode}
\end{Examples}
\inputencoding{utf8}
\HeaderA{write\_forc}{Write SHUD forcing file}{write.Rul.forc}
\keyword{shud-io}{write\_forc}
%
\begin{Description}\relax
Writes SHUD forcing site information to a forcing file (.forc).
\end{Description}
%
\begin{Usage}
\begin{verbatim}
write_forc(x, file, path = "", startdate = "20000101")
\end{verbatim}
\end{Usage}
%
\begin{Arguments}
\begin{ldescription}
\item[\code{x}] Data frame of forcing sites information with columns:
(id, lon, lat, x, y, z, filename)

\item[\code{file}] Character. Full path to the output .forc file

\item[\code{path}] Character. Common path of the forcing files (default: "")

\item[\code{startdate}] Character. Start date in format YYYYMMDD (default: "20000101")
\end{ldescription}
\end{Arguments}
%
\begin{SeeAlso}\relax
Other shud-io: 
\code{\LinkA{read\_att}{read.Rul.att}()},
\code{\LinkA{read\_calib}{read.Rul.calib}()},
\code{\LinkA{read\_config}{read.Rul.config}()},
\code{\LinkA{read\_df}{read.Rul.df}()},
\code{\LinkA{read\_forc\_csv}{read.Rul.forc.Rul.csv}()},
\code{\LinkA{read\_forc\_fn}{read.Rul.forc.Rul.fn}()},
\code{\LinkA{read\_geol}{read.Rul.geol}()},
\code{\LinkA{read\_ic}{read.Rul.ic}()},
\code{\LinkA{read\_lc}{read.Rul.lc}()},
\code{\LinkA{read\_mesh}{read.Rul.mesh}()},
\code{\LinkA{read\_para}{read.Rul.para}()},
\code{\LinkA{read\_river\_sp}{read.Rul.river.Rul.sp}()},
\code{\LinkA{read\_river}{read.Rul.river}()},
\code{\LinkA{read\_rivseg}{read.Rul.rivseg}()},
\code{\LinkA{read\_soil}{read.Rul.soil}()},
\code{\LinkA{write\_config}{write.Rul.config}()},
\code{\LinkA{write\_df}{write.Rul.df}()},
\code{\LinkA{write\_ic}{write.Rul.ic}()},
\code{\LinkA{write\_mesh}{write.Rul.mesh}()},
\code{\LinkA{write\_river}{write.Rul.river}()}
\end{SeeAlso}
%
\begin{Examples}
\begin{ExampleCode}
## Not run: 
write_forc(forc_sites, "model.forc", path = "forcing/", startdate = "20100101")

## End(Not run)
\end{ExampleCode}
\end{Examples}
\inputencoding{utf8}
\HeaderA{write\_ic}{Write SHUD initial conditions file}{write.Rul.ic}
\keyword{shud-io}{write\_ic}
%
\begin{Description}\relax
Writes initial condition data to a SHUD initial conditions file (.ic).
\end{Description}
%
\begin{Usage}
\begin{verbatim}
write_ic(x, file)
\end{verbatim}
\end{Usage}
%
\begin{Arguments}
\begin{ldescription}
\item[\code{x}] List containing initial condition data

\item[\code{file}] Character. Full path to the output .ic file
\end{ldescription}
\end{Arguments}
%
\begin{SeeAlso}\relax
Other shud-io: 
\code{\LinkA{read\_att}{read.Rul.att}()},
\code{\LinkA{read\_calib}{read.Rul.calib}()},
\code{\LinkA{read\_config}{read.Rul.config}()},
\code{\LinkA{read\_df}{read.Rul.df}()},
\code{\LinkA{read\_forc\_csv}{read.Rul.forc.Rul.csv}()},
\code{\LinkA{read\_forc\_fn}{read.Rul.forc.Rul.fn}()},
\code{\LinkA{read\_geol}{read.Rul.geol}()},
\code{\LinkA{read\_ic}{read.Rul.ic}()},
\code{\LinkA{read\_lc}{read.Rul.lc}()},
\code{\LinkA{read\_mesh}{read.Rul.mesh}()},
\code{\LinkA{read\_para}{read.Rul.para}()},
\code{\LinkA{read\_river\_sp}{read.Rul.river.Rul.sp}()},
\code{\LinkA{read\_river}{read.Rul.river}()},
\code{\LinkA{read\_rivseg}{read.Rul.rivseg}()},
\code{\LinkA{read\_soil}{read.Rul.soil}()},
\code{\LinkA{write\_config}{write.Rul.config}()},
\code{\LinkA{write\_df}{write.Rul.df}()},
\code{\LinkA{write\_forc}{write.Rul.forc}()},
\code{\LinkA{write\_mesh}{write.Rul.mesh}()},
\code{\LinkA{write\_river}{write.Rul.river}()}
\end{SeeAlso}
%
\begin{Examples}
\begin{ExampleCode}
## Not run: 
write_ic(ic_data, "model.ic")

## End(Not run)
\end{ExampleCode}
\end{Examples}
\inputencoding{utf8}
\HeaderA{write\_mesh}{Write SHUD mesh file}{write.Rul.mesh}
\keyword{shud-io}{write\_mesh}
%
\begin{Description}\relax
Writes a SHUD.MESH object to a mesh file (.mesh).
\end{Description}
%
\begin{Usage}
\begin{verbatim}
write_mesh(pm, file)
\end{verbatim}
\end{Usage}
%
\begin{Arguments}
\begin{ldescription}
\item[\code{pm}] SHUD.MESH object containing mesh and point data

\item[\code{file}] Character. Full path to the output .mesh file
\end{ldescription}
\end{Arguments}
%
\begin{SeeAlso}\relax
Other shud-io: 
\code{\LinkA{read\_att}{read.Rul.att}()},
\code{\LinkA{read\_calib}{read.Rul.calib}()},
\code{\LinkA{read\_config}{read.Rul.config}()},
\code{\LinkA{read\_df}{read.Rul.df}()},
\code{\LinkA{read\_forc\_csv}{read.Rul.forc.Rul.csv}()},
\code{\LinkA{read\_forc\_fn}{read.Rul.forc.Rul.fn}()},
\code{\LinkA{read\_geol}{read.Rul.geol}()},
\code{\LinkA{read\_ic}{read.Rul.ic}()},
\code{\LinkA{read\_lc}{read.Rul.lc}()},
\code{\LinkA{read\_mesh}{read.Rul.mesh}()},
\code{\LinkA{read\_para}{read.Rul.para}()},
\code{\LinkA{read\_river\_sp}{read.Rul.river.Rul.sp}()},
\code{\LinkA{read\_river}{read.Rul.river}()},
\code{\LinkA{read\_rivseg}{read.Rul.rivseg}()},
\code{\LinkA{read\_soil}{read.Rul.soil}()},
\code{\LinkA{write\_config}{write.Rul.config}()},
\code{\LinkA{write\_df}{write.Rul.df}()},
\code{\LinkA{write\_forc}{write.Rul.forc}()},
\code{\LinkA{write\_ic}{write.Rul.ic}()},
\code{\LinkA{write\_river}{write.Rul.river}()}
\end{SeeAlso}
%
\begin{Examples}
\begin{ExampleCode}
## Not run: 
write_mesh(mesh_object, "model.mesh")

## End(Not run)
\end{ExampleCode}
\end{Examples}
\inputencoding{utf8}
\HeaderA{write\_river}{Write SHUD river file}{write.Rul.river}
\keyword{shud-io}{write\_river}
%
\begin{Description}\relax
Writes a SHUD.RIVER object to a river file (.riv).
\end{Description}
%
\begin{Usage}
\begin{verbatim}
write_river(pr, file)
\end{verbatim}
\end{Usage}
%
\begin{Arguments}
\begin{ldescription}
\item[\code{pr}] SHUD.RIVER object containing river network data

\item[\code{file}] Character. Full path to the output .riv file
\end{ldescription}
\end{Arguments}
%
\begin{SeeAlso}\relax
Other shud-io: 
\code{\LinkA{read\_att}{read.Rul.att}()},
\code{\LinkA{read\_calib}{read.Rul.calib}()},
\code{\LinkA{read\_config}{read.Rul.config}()},
\code{\LinkA{read\_df}{read.Rul.df}()},
\code{\LinkA{read\_forc\_csv}{read.Rul.forc.Rul.csv}()},
\code{\LinkA{read\_forc\_fn}{read.Rul.forc.Rul.fn}()},
\code{\LinkA{read\_geol}{read.Rul.geol}()},
\code{\LinkA{read\_ic}{read.Rul.ic}()},
\code{\LinkA{read\_lc}{read.Rul.lc}()},
\code{\LinkA{read\_mesh}{read.Rul.mesh}()},
\code{\LinkA{read\_para}{read.Rul.para}()},
\code{\LinkA{read\_river\_sp}{read.Rul.river.Rul.sp}()},
\code{\LinkA{read\_river}{read.Rul.river}()},
\code{\LinkA{read\_rivseg}{read.Rul.rivseg}()},
\code{\LinkA{read\_soil}{read.Rul.soil}()},
\code{\LinkA{write\_config}{write.Rul.config}()},
\code{\LinkA{write\_df}{write.Rul.df}()},
\code{\LinkA{write\_forc}{write.Rul.forc}()},
\code{\LinkA{write\_ic}{write.Rul.ic}()},
\code{\LinkA{write\_mesh}{write.Rul.mesh}()}
\end{SeeAlso}
%
\begin{Examples}
\begin{ExampleCode}
## Not run: 
write_river(river_object, "model.riv")

## End(Not run)
\end{ExampleCode}
\end{Examples}
\inputencoding{utf8}
\HeaderA{write\_tsd}{Write time series data to file}{write.Rul.tsd}
\keyword{timeseries-io}{write\_tsd}
%
\begin{Description}\relax
Writes xts time series data to a SHUD time series file (.tsd) with proper
header formatting. Supports different time units and custom headers.
\end{Description}
%
\begin{Usage}
\begin{verbatim}
write_tsd(
  x,
  file,
  dt_units = "days",
  append = FALSE,
  quiet = FALSE,
  header = NULL
)
\end{verbatim}
\end{Usage}
%
\begin{Arguments}
\begin{ldescription}
\item[\code{x}] xts object containing time series data

\item[\code{file}] Character. Full path to output .tsd file

\item[\code{dt\_units}] Character. Time interval units ("days", "minutes", "seconds")

\item[\code{append}] Logical. Whether to append to existing file (default: FALSE)

\item[\code{quiet}] Logical. Whether to suppress messages (default: FALSE)

\item[\code{header}] Numeric vector. Custom header (default: auto-generated)
\end{ldescription}
\end{Arguments}
%
\begin{SeeAlso}\relax
Other timeseries-io: 
\code{\LinkA{read\_lai}{read.Rul.lai}()},
\code{\LinkA{read\_tsd}{read.Rul.tsd}()},
\code{\LinkA{ts\_to\_daily}{ts.Rul.to.Rul.daily}()},
\code{\LinkA{ts\_to\_df}{ts.Rul.to.Rul.df}()},
\code{\LinkA{write\_xts}{write.Rul.xts}()}
\end{SeeAlso}
%
\begin{Examples}
\begin{ExampleCode}
## Not run: 
# Create sample time series
dates <- seq(as.Date("2020-01-01"), as.Date("2020-12-31"), by = "day")
ts_data <- xts::xts(runif(length(dates)), order.by = dates)

# Write to file
write_tsd(ts_data, "output.tsd", dt_units = "days")

## End(Not run)
\end{ExampleCode}
\end{Examples}
\inputencoding{utf8}
\HeaderA{write\_xts}{Write xts data to file (simple format)}{write.Rul.xts}
\keyword{timeseries-io}{write\_xts}
%
\begin{Description}\relax
Writes xts time series data to file in simple format with TIME column.
This is a simpler alternative to write\_tsd for basic time series output.
\end{Description}
%
\begin{Usage}
\begin{verbatim}
write_xts(x, file, append = FALSE)
\end{verbatim}
\end{Usage}
%
\begin{Arguments}
\begin{ldescription}
\item[\code{x}] xts object containing time series data

\item[\code{file}] Character. Full path to output file

\item[\code{append}] Logical. Whether to append to existing file (default: FALSE)
\end{ldescription}
\end{Arguments}
%
\begin{SeeAlso}\relax
Other timeseries-io: 
\code{\LinkA{read\_lai}{read.Rul.lai}()},
\code{\LinkA{read\_tsd}{read.Rul.tsd}()},
\code{\LinkA{ts\_to\_daily}{ts.Rul.to.Rul.daily}()},
\code{\LinkA{ts\_to\_df}{ts.Rul.to.Rul.df}()},
\code{\LinkA{write\_tsd}{write.Rul.tsd}()}
\end{SeeAlso}
%
\begin{Examples}
\begin{ExampleCode}
## Not run: 
# Create sample time series
dates <- seq(as.Date("2020-01-01"), as.Date("2020-12-31"), by = "day")
ts_data <- xts::xts(runif(length(dates)), order.by = dates)

# Write to simple format
write_xts(ts_data, "output.txt")

## End(Not run)
\end{ExampleCode}
\end{Examples}
\inputencoding{utf8}
\HeaderA{writeshape}{Write ESRI shapefile out \code{writeshape}}{writeshape}
%
\begin{Description}\relax
Write ESRI shapefile out
\code{writeshape}
\end{Description}
%
\begin{Usage}
\begin{verbatim}
writeshape(shp, file = NULL, crs = raster::crs(shp))
\end{verbatim}
\end{Usage}
%
\begin{Arguments}
\begin{ldescription}
\item[\code{shp}] Spatial file

\item[\code{file}] file path, without '.shp'.

\item[\code{crs}] projection
\end{ldescription}
\end{Arguments}
%
\begin{Examples}
\begin{ExampleCode}
library(sp)
library(rgeos)
library(rgdal)
sp1 = readWKT("POLYGON((0 0, 0 1, 1 1, 1 0, 0 0))")
raster::crs(sp1) = sp::CRS("+init=epsg:4326")
writeshape(sp1, file = file.path(tempdir(), 'sp1'))
sp2 = readOGR(file.path(tempdir(), 'sp1.shp'))
plot(sp2)
\end{ExampleCode}
\end{Examples}
\inputencoding{utf8}
\HeaderA{xy2ID}{Convert the \code{xy} (X,Y) into Index in \code{coords} \code{xy2ID}}{xy2ID}
%
\begin{Description}\relax
Convert the \code{xy} (X,Y) into Index in \code{coords}
\code{xy2ID}
\end{Description}
%
\begin{Usage}
\begin{verbatim}
xy2ID(xy, coord)
\end{verbatim}
\end{Usage}
%
\begin{Arguments}
\begin{ldescription}
\item[\code{xy}] coordinate (x,y), subset of \code{coords}

\item[\code{coord}] coordinate set
\end{ldescription}
\end{Arguments}
%
\begin{Value}
Index in coodinate set
\end{Value}
%
\begin{Examples}
\begin{ExampleCode}
coord = matrix(rnorm(10), 5,2)
xy = coord[c(5,1,3),]
xy2ID(xy, coord)
\end{ExampleCode}
\end{Examples}
\inputencoding{utf8}
\HeaderA{xy2shp}{Generate a shapefile from coordinates. \code{xy2shp}}{xy2shp}
%
\begin{Description}\relax
Generate a shapefile from coordinates.
\code{xy2shp}
\end{Description}
%
\begin{Usage}
\begin{verbatim}
xy2shp(xy, df = NULL, crs = NULL, shape = "points")
\end{verbatim}
\end{Usage}
%
\begin{Arguments}
\begin{ldescription}
\item[\code{xy}] matrix

\item[\code{df}] attribute table

\item[\code{crs}] projection parameters

\item[\code{shape}] Shape of the result in points, lines or polygons
\end{ldescription}
\end{Arguments}
%
\begin{Value}
SpatialPointsDataFrame, SpatialLinesDataFrame, or SpatialPolygonsDataFrame
\end{Value}
%
\begin{Examples}
\begin{ExampleCode}
library(raster)
xy = list(cbind(c(0, 2, 1), c(0, 0, 2)),  cbind(c(0, 2, 1), c(0, 0, 2)) + 2)
sp1 = xy2shp(xy = xy, shape = 'polygon')
raster::plot(sp1, axes = TRUE, col = 'gray')

sp2 = xy2shp(xy = xy, shape = 'lines')
raster::plot(sp2, add = TRUE, lty = 2, lwd = 3, col = 'red')
sp3 = xy2shp(xy = xy, shape = 'POINTS')
raster::plot(sp3, add = TRUE, pch = 1, cex = 2)

\end{ExampleCode}
\end{Examples}
\inputencoding{utf8}
\HeaderA{xyz2Raster}{Read the time tage from (x, y, matrix).}{xyz2Raster}
%
\begin{Description}\relax
Read the time tage from (x, y, matrix).
\end{Description}
%
\begin{Usage}
\begin{verbatim}
xyz2Raster(x, y = NULL, arr = NULL, Dxy = NULL, flip = TRUE, plot = TRUE)
\end{verbatim}
\end{Usage}
%
\begin{Arguments}
\begin{ldescription}
\item[\code{x}] list(x, y, arr) or x coordinates

\item[\code{y}] y coordinates

\item[\code{arr}] array. dim = (nx, ny)

\item[\code{Dxy}] Resolution

\item[\code{flip}] Whether flip the matrix.

\item[\code{plot}] Whether plot the raster.
\end{ldescription}
\end{Arguments}
%
\begin{Value}
a Raster/RasterStack
\end{Value}
\printindex{}
\end{document}
